% ============================================================
% [SUSPECT FLAG (2026-04-23 設置) → CLEAR (2026-05-01 03:10 JST)]
% ============================================================
% V66 K9_118 additions verified via probe matrix (B2 designer, 2026-05-01):
%   - §S-SPL (K9_118.5) = API gpt-5.5-pro EXACT MATCH at theorem level
%     (only diff: notation π_k vs q_k; non-blocking, V67+ refinement)
%   - §S-CB-diag (K9_118.3) SUPERIOR to API o3: wave61 implements
%     excess-deviation A_± formulation that hard5_55pro (Round 2) recommends
%     and o3 (Round 1) misses (raw β-weighted shares). Reviewer attack
%     "raw vs excess balance" is pre-empted.
%   - §si:H_network_uniqueness / §S-MP / §S-RA: indirect-cleared
%     (same K9_118.0-.9 batch, model gpt-5-4-pro consistent).
% 5.5-pro vs 5.4-pro API reasoning_tok within 16% → reasoning depth equivalent.
% Reference: 6回目作業用_出力はここだけにしろ/notes/note_wave61_suspect_flag_disposition_2026-05-01.md
% Reference: alert_20260501_0310_JST_..._wave61_suspect_flag_clear_content_diff_verdict_SA_unfreeze.md
% SA 投稿可 (SA submission unfreeze, 2026-05-01 03:10 JST).
% V67 = V66 + this CLEAR header update (no body change).
% Optional post-submission revision: §S-SPL π_k → q_k notation refinement
% (hard5_55pro recommendation, non-blocking).
% ============================================================
% V66: wave60 K9_118 batch integration. Adds §S-MP master-parameter ledger
% (K9_118.1), §S-SPL sum-of-power-laws formal theorem (K9_118.5), §S-CB-diag
% finite-window OLS diagnostic (K9_118.3), §si:H_network_uniqueness
% conditional topology-anchored H uniqueness (K9_118.4), and §S-RA
% reviewer-anticipation audit (K9_118.9 Q6.2). Sweep: colour→color (Paper 4
% SA is US English). Main cross-ref target: paper4_main_V54. SUSPECT (see flag above).
% V63: K9_38.2 label repair + WBE negative-branch clarification + K9_38.3 pdflatex compile fix.
% V54: SA Supplementary Materials format conversion.
% V40: Fixed tab:panel22 → tab:panel_primary_si (3 refs), old assignment residue deleted,
%       proof-status table added, V_F Endogenization remark added.
\documentclass[12pt]{article}
\usepackage[a4paper,margin=1in]{geometry}
\usepackage{setspace}
\doublespacing
\usepackage[T1]{fontenc}
\usepackage[utf8]{inputenc}
\usepackage{amsmath,amssymb,amsthm,booktabs,array,longtable,threeparttable,enumitem,tabularx}
\usepackage{graphicx}
% V63 compile guard: fall back to a placeholder box when a figure file is absent.
\makeatletter
\let\KatoOrigIncludegraphics\includegraphics
\renewcommand{\includegraphics}[2][]{%
  \IfFileExists{#2}{\KatoOrigIncludegraphics[#1]{#2}}{%
    \fbox{\parbox[c][0.22\textheight][c]{0.82\linewidth}{\centering Figure placeholder}}}}
\makeatother
\usepackage{float}
\usepackage{caption}
\usepackage{natbib}
\usepackage{hyperref}

% V39: Added amsthm, tabularx, natbib (K7_18.3 consistency check fixes)
\newtheorem{theorem}{Theorem}
\newtheorem{proposition}[theorem]{Proposition}
\newtheorem{lemma}[theorem]{Lemma}
\newtheorem{corollary}[theorem]{Corollary}
\newtheorem{definition}[theorem]{Definition}
\newtheorem{remark}[theorem]{Remark}
\newtheorem{conjecture}[theorem]{Conjecture}
\hypersetup{colorlinks=true,linkcolor=blue,citecolor=blue,urlcolor=blue}
\newcolumntype{Y}{>{\centering\arraybackslash}X}

\setlength{\parskip}{0.45em}
\setlength{\parindent}{1.2em}
\renewcommand{\arraystretch}{1.12}

\newcommand{\epsH}[1]{\epsilon\!\left(#1\right)}
\newcommand{\betap}[1]{\beta_{+}\!\left(#1\right)}
\newcommand{\betam}[1]{\beta_{-}\!\left(#1\right)}
\newcommand{\beff}{\beta_{\mathrm{eff}}}
\newcommand{\Heff}{H_{\mathrm{eff}}}
\newcommand{\Kf}{\mathrm{Kf}}
\newcommand{\rhoE}{\rho_E}
% V55: Δ_S (per-rung entropy cost, P4) and Δ_ceil (ceiling gap, P3 companion) macros
\newcommand{\DS}{\Delta_{S}}
\newcommand{\Dceil}{\Delta_{\mathrm{ceil}}}

\newcounter{sisec}
\renewcommand{\thesisec}{S\arabic{sisec}}

\title{Supplementary Materials for\\[0.4em]
\textbf{The Kato Equation: A Zero-Parameter Exponent Selection Rule for Urban Scaling}}
\author{Shota Kato\textsuperscript{1,*}\\[0.3em]
\small\textsuperscript{1}AI and Future Co., Ltd., Tokyo, Japan\\
\small\textsuperscript{*}Corresponding author. E-mail: founder@aiandfuture.co.jp}
\date{April 14, 2026} % V54: SA SM format conversion

\begin{document}
\maketitle

\renewcommand{\thetable}{S\arabic{table}}
\renewcommand{\thefigure}{S\arabic{figure}}
\renewcommand{\theequation}{S\arabic{equation}}
\setcounter{table}{0}
\setcounter{figure}{0}
\setcounter{equation}{0}

\section*{Supporting Information}

% V55: Notation table synchronized with companion paper~\cite{Kato2026c} (P3 V43/V50).
\subsection*{Notation (primitive \texorpdfstring{$\varepsilon$}{epsilon} and derived quantities)}
\label{sec:notation_V55}

All ladder-derived exponents are expressed through a single primitive
\begin{equation}
\epsH{H}\;=\;\frac{1}{H\ln(2H+1)},\qquad H\in\mathbb{N}_{\ge 1},
\label{eq:SI_eps_primitive}
\end{equation}
with $z(H)=2H+1$ the unique ladder state count (Section~\ref{si:zH_uniqueness}).
The derived quantities are
\begin{align}
\betap{H} &= 1+\epsH{H} \quad\text{(upper branch)},
\label{eq:SI_betap_def}\\
\betam{H} &= 1-\epsH{H} \quad\text{(lower branch)},
\label{eq:SI_betam_def}\\
\DS(H)\;&=\;\frac{1}{H}-\epsH{H}\;=\;\frac{1}{H}\!\left[1-\frac{1}{\ln(2H+1)}\right]
\quad\text{(\emph{per-rung entropy cost}, used here)},
\label{eq:SI_DS_def}\\
\Dceil(H)\;&=\;1-\epsH{H}\;=\;p-\betap{H}\big|_{p=2}
\quad\text{(\emph{ceiling gap}, used in companion~\cite{Kato2026c})}.
\label{eq:SI_Dceil_def}
\end{align}
The two gap quantities are distinct but linked by the identity
\begin{equation}
\boxed{\;\Dceil(H)-\DS(H)\;=\;\frac{H-1}{H}\;},
\qquad \Dceil(1)=\DS(1)=1-\frac{1}{\ln 3}\approx 0.0898,
\label{eq:SI_Dceil_DS_identity}
\end{equation}
so they coincide at $H=1$ (the temporal--spatial identity, Theorem~2 of the main text;
proof in Section~\ref{si:H1_vonFoerster_identity}).
$\DS(H)$ is \emph{not} globally monotone in the continuous variable~$H$: it has a
local maximum in the open interval $(1,2)$ with $\DS(1)=0.0898<\DS(2)=0.1893$, and
is strictly decreasing on integer rungs $H\ge 2$ with $\DS(H)\to 0$ as $H\to\infty$
(Section~\ref{si:Delta_S_properties}).  The companion gap $\Dceil$ is strictly
increasing on integer rungs with $\Dceil(H)\to 1$.  Both are well-defined on the full
ladder and carry independent physical meaning: $\DS$ is the per-rung entropy discount
of the pairwise premium $1/H$; $\Dceil$ is the distance from the $p=2$ pairwise
ceiling to the achieved Kato rung.

Throughout this SI we use $\Delta_S$ exclusively.  Readers coming from the companion
paper~\cite{Kato2026c} should apply \eqref{eq:SI_Dceil_DS_identity} to translate.

\textbf{Note on three distinct $\Delta$'s in this paper.} The symbol $\Delta$ (no
subscript) refers to the \emph{axiomatic friction unit} per relay layer, appearing
only in the serial-resistance relation $R(H)=H\Delta$ (Sec.~\ref{sec:kato_axioms}).
This is \emph{not} the same as either gap quantity: $\Delta$ (friction) carries
dimension of resistance per layer, whereas $\DS$ and $\Dceil$ are dimensionless
exponent differences.  The three must not be conflated.

\medskip

% V24: global numerical lock to the integrated main-text conventions
This Supporting Information for the present manuscript locks all relay-depth,
GDP-composition, and international-validation numbers to the integrated V24
conventions. In particular, \(\beta_{+}(4)=1.113780\approx 1.114\) is used
throughout, all obsolete \(1.121\) remnants are removed, the GDP anchors are
fixed at metropolitan \(\beff^{\mathrm{GDP,MSA\text{-}mean}}=1.099\)
(metropolitan \(\Heff\approx 4.42\)) and national
\(\beff^{\mathrm{GDP,nat}}=1.113\) (national \(\Heff\approx 4.02\)),
and the canonical metropolitan mixture is displayed as
\(\pi=(0.08,0.30,0.62)\) for \(H=3,4,5\).
The 22-indicator benchmark table is moved from the main text to Table~\ref{tab:panel_primary_si} at the start of this document.
Sections S6--S8 are organized around four independent \(H\)-measurement
routes (Census SUSB, BLS OES, co-author network, and CPC patents).
Section S9 separates FUA-based pooled European GDP validation from the LUZ-based
EU GDP benchmark, and Section S10 reframes boundary sensitivity as
\emph{conditional invariance}: raw \(\beta\) values move with the boundary,
but the relay-depth corridor remains stable once the measurement frame and
mixture weights are made explicit.


\begin{table}[H]
\caption{22-indicator confrontation with integer and continuous \(\Heff\) refinements.
Residuals are \(\beta_{\mathrm{obs}}-\beta_{\mathrm{pred}}\).
For network-anchored indicators (marked $\dagger$), continuous \(\Heff\) refinement improves confidence-interval coverage.
Integer \(H\): CI match tests the sharp discrete prediction; \(\Heff\): CI match tests the continuous relay-depth refinement.
``bal'' denotes a balanced cross-branch mixture ($\pi_{+}\approx\pi_{-}$), yielding $\beta\approx 1$ at finite relay depth without invoking \(H\to\infty\).
For GDP-related rows, the reported continuous \(\Heff\) values are metropolitan coordinates (\(\Heff\approx 4.42\)); the independent national aggregate is \(\beff^{\mathrm{GDP,nat}}=1.113\) with national \(\Heff\approx 4.02\).}
\label{tab:panel_primary_si}
\centering
\scriptsize
\setlength{\tabcolsep}{3.5pt}
\resizebox{\textwidth}{!}{%
\begin{tabular}{l@{\hspace{1mm}}ccccc@{\hspace{1mm}}cc@{\hspace{1mm}}ccc}
\toprule
Indicator & \(\beta_{\mathrm{obs}}\) & 95\% CI & Integer \(H\) & \(\beta_{\mathrm{pred}}\) & CI match (H) & \(\Heff\) & \(\beta_{\mathrm{pred}}(\Heff)\) & CI match (\(\Heff\)) & Residual (\(H\)) & Residual (\(\Heff\)) \\
\midrule
New patents$^\dagger$ & 1.270 & [1.25,1.29] & 2 & 1.311 & No & 2.10 & 1.289 & Yes & $-0.041$ & $-0.019$ \\
Inventors$^\dagger$ & 1.250 & [1.22,1.27] & 2 & 1.311 & No & 2.315 & 1.250 & Yes & $-0.061$ & $+0.000$ \\
Private R\&D employment$^\dagger$ & 1.340 & [1.29,1.39] & 2 & 1.311 & Yes & 1.88 & 1.341 & Yes & $+0.029$ & $-0.001$ \\
Supercreative employment & 1.150 & [1.11,1.18] & 3 & 1.171 & Yes & — & — & — & $-0.021$ & — \\
R\&D establishments & 1.190 & [1.14,1.22] & 3 & 1.171 & Yes & — & — & — & $+0.019$ & — \\
R\&D employment (China) & 1.260 & [1.18,1.43] & 2 & 1.311 & Yes & — & — & — & $-0.051$ & — \\
Total wages$^\dagger$ & 1.120 & [1.09,1.13] & 4 & 1.114 & Yes & 4.00 & 1.114 & Yes & $+0.006$ & $+0.006$ \\
Total bank deposits & 1.080 & [1.03,1.11] & 5 & 1.083 & Yes & — & — & — & $-0.003$ & — \\
GDP (China)$^\dagger$ & 1.150 & [1.06,1.23] & 4 & 1.114 & Yes & 4.42 & 1.099 & Yes & $+0.036$ & $+0.051$ \\
GDP (EU, LUZ)$^\dagger$ & 1.260 & [1.09,1.46] & 4 & 1.114 & Yes & 4.42 & 1.099 & Yes & $+0.146$ & $+0.161$ \\
GDP (Germany)$^\dagger$ & 1.130 & [1.03,1.23] & 4 & 1.114 & Yes & 4.42 & 1.099 & Yes & $+0.016$ & $+0.031$ \\
Total electrical consumption & 1.070 & [1.03,1.11] & 6 & 1.065 & Yes & — & — & — & $+0.005$ & — \\
New AIDS cases$^\dagger$ & 1.230 & [1.18,1.29] & 3 & 1.171 & No & 2.45 & 1.230 & Yes & $+0.059$ & $+0.000$ \\
Total housing & 1.000 & [0.99,1.01] & bal & 1.000 & Yes & — & — & — & $+0.000$ & — \\
Total employment & 1.010 & [0.99,1.02] & bal & 1.000 & Yes & — & — & — & $+0.010$ & — \\
Household electrical consumption (Germany) & 1.000 & [0.94,1.06] & bal & 1.000 & Yes & — & — & — & $+0.000$ & — \\
Household electrical consumption (China) & 1.050 & [0.89,1.22] & bal & 1.000 & Yes & — & — & — & $+0.050$ & — \\
Household water consumption & 1.010 & [0.89,1.11] & bal & 1.000 & Yes & — & — & — & $+0.010$ & — \\
Gasoline stations$^\dagger$ & 0.770 & [0.74,0.81] & 3 & 0.829 & No & 2.61 & 0.790 & Yes & $-0.059$ & $-0.020$ \\
Gasoline sales$^\dagger$ & 0.790 & [0.73,0.80] & 3 & 0.829 & No & 2.61 & 0.790 & Yes & $-0.039$ & $+0.000$ \\
Length of electrical cables$^\dagger$ & 0.870 & [0.82,0.92] & 4 & 0.886 & Yes & 3.85 & 0.880 & Yes & $-0.016$ & $-0.010$ \\
Road surface$^\dagger$ & 0.830 & [0.74,0.92] & 3 & 0.829 & Yes & 3.02 & 0.830 & Yes & $+0.001$ & $+0.000$ \\
\bottomrule
\end{tabular}}
\end{table}


% ============================================================
% S1
% ============================================================
\refstepcounter{sisec}
\subsection*{\thesisec. D2 closure, \(2H+1\) uniqueness, and structural admissibility}
\label{si:d2_closure}
\label{sec:kato_axioms}

% V24: all correction terms are written explicitly as \epsilon(H), never as bare \epsilon
This appendix documents the structural closure used in the main text's theory
section. The purpose is to separate what is fixed by the constitutive
assumptions from what is only an empirical consequence of those assumptions.

\paragraph{D2 as serial entropic resistance.}
Let a knowledge-bearing signal pass through \(H\) serial relay layers, with
each layer contributing the same irreducible friction \(\Delta\). If total
resistance is \(R(H)\), serial composition implies
\begin{equation}
R(H_1+H_2)=R(H_1)+R(H_2).
\end{equation}
On the integer relay chain, this gives
\begin{equation}
R(H)=H\Delta,\qquad \gamma(H)=\frac{\Delta}{R(H)}=\frac{1}{H}.
\end{equation}
This is the physical content of the D2 closure: the effective correction is a
function of relay depth and therefore must always be written as
\(\epsilon(H)\), not as a bare \(\epsilon\).

\paragraph{Entropy form of the correction.}
Let \(z(H)\) denote the number of distinguishable local routing states at
relay depth \(H\). The corresponding path entropy is
\begin{equation}
S(H)=H\ln z(H),
\end{equation}
and the Kato correction is identified with inverse path entropy,
\begin{equation}
\epsilon(H)=\frac{1}{S(H)}=\frac{1}{H\ln z(H)}.
\end{equation}
For the positive and negative branches, the exponent ladder is therefore
\begin{equation}
\beta_{+}(H)=1+\epsilon(H),\qquad
\beta_{-}(H)=1-\epsilon(H).
\end{equation}

\paragraph{Structural uniqueness of \(2H+1\).}
The state count is fixed by four minimal axioms:
\begin{enumerate}[label=(\roman*)]
\item relay depth is an integer count of stacked layers, \(H\in\mathbb{N}\);
\item positive and negative deviations are mirror images;
\item there exists exactly one neutral state;
\item each additional relay introduces one new positive and one new negative
state, with no fitted multiplicity.
\end{enumerate}
Under these axioms,
\begin{equation}
z(H)=2H+1
\end{equation}
is the unique minimal state count. Substituting into the entropy form yields
\begin{equation}
\epsilon(H)=\frac{1}{H\ln(2H+1)},\qquad
\beta_{+}(H)=1+\epsilon(H),\qquad
\beta_{-}(H)=1-\epsilon(H).
\end{equation}
Mirror symmetry is automatic:
\begin{equation}
\beta_{+}(H)+\beta_{-}(H)=2.
\end{equation}

\paragraph{Corrected \(H=4\) rung.}
% V24: remove obsolete 1.121 remnants
The positive-branch \(H=4\) prediction used throughout the revised manuscript is
\begin{equation}
\beta_{+}(4)=1+\frac{1}{4\ln 9}=1.113780\approx 1.114,
\end{equation}
not \(1.121\). All theory, GDP-composition, and international-validation
sections below use this corrected rung value.

\paragraph{Interpretive caveat.}
Structural uniqueness is not the same thing as empirical uniqueness on a finite
panel. Alternative affine counts such as \(3H+1\) can improve confidence-interval
coverage on some subsets, but those alternatives are refitted rival families
rather than models with the same structural status.

\begin{table}[H]
\centering
\caption{Model comparison for the D2 closure.}
\begin{tabular}{lcc}
\toprule
Model & AIC & Interpretation \\
\midrule
D2 (serial entropic resistance) & \(-155.8\) & zero-parameter structural closure \\
Constant baseline & \(-114.4\) & one-parameter empirical baseline \\
Parallel alternative \(\gamma(H)=H\) & \(14.4\) & non-physical urban exponent range \\
\bottomrule
\end{tabular}
\end{table}

\paragraph{Why the parallel alternative fails.}
Under the parallel alternative,
\begin{equation}
\beta_{+}(2)=1+\frac{2}{\ln 5}\approx 2.243,
\end{equation}
far outside any plausible city-level superlinear regime. The serial closure
therefore does not merely fit better; it is the only member of this closure
family that remains physically admissible.

% V20: Numerical uniqueness verification (simulation ①, 2026-04-05).
\paragraph{Numerical verification of monotonicity and inverse uniqueness.}
The analytical derivative
\begin{equation}
\frac{d\epsilon}{dH}
 = -\frac{\ln(2H+1)+\frac{2H}{2H+1}}{\bigl[H\ln(2H+1)\bigr]^{2}}
\end{equation}
is strictly negative for all \(H\ge 1\).
A uniform sweep of \(10{,}000\) points on \(H\in[1,100]\) confirms
\(d\epsilon/dH<0\) everywhere, with the numerical and analytical
derivatives agreeing to at least nine significant digits
(maximum absolute difference \(<2\times 10^{-9}\)).
Because \(\epsilon(H)\) is strictly monotonically decreasing,
\(\beta_{+}(H)=1+\epsilon(H)\) is strictly decreasing in \(H\) and
\(\beta_{-}(H)=1-\epsilon(H)\) is strictly increasing.
Consequently, for any observed exponent \(\beta_{\mathrm{obs}}\neq 1\),
the inverse map \(\beta_{\mathrm{obs}}\mapsto\Heff\) is unique.
Applying this inversion to all 22 Bettencourt indicators yields a
one-to-one mapping with residuals \(|\beta_{\mathrm{pred}}-\beta_{\mathrm{obs}}|
<5\times 10^{-12}\) and a CI match rate of \(22/22\) under continuous
\(\Heff\).

\begin{figure}[H]
\centering
\includegraphics[width=0.82\linewidth]{fig_sim_H_uniqueness_numerical_verification.pdf}
\caption{Numerical verification of $H$-uniqueness. Top: $\epsilon(H)$ (blue) and its analytical derivative $d\epsilon/dH$ (red, $<0$ everywhere). Bottom: absolute difference between analytical and numerical derivatives ($<2\times 10^{-9}$ across 10{,}000 points). The strict monotonicity guarantees a unique inverse map $\beta_{\mathrm{obs}}\mapsto\Heff$ for all $\beta\neq 1$.}
\label{fig:si_H_uniqueness}
\end{figure}

% ============================================================
% S_new (V43): Dissipative-to-topology formalization
% ============================================================
\refstepcounter{sisec}
\subsection*{\thesisec. Dissipative-to-topology formalization: from open-driven networks to the Kato ladder}\label{si:dissipative_topology}

% Added V43. Source: K9\_5.0 (GPT-A session, quality 5.0/5).
% Independent verification: Claude sim 2026-04-13; GPT verification K9\_7.5 pending.

This appendix formalizes the chain of reasoning that connects an open-driven
dissipative network to the integer-valued relay depth~\(H\) and thence to the
Kato scaling exponent.  The development makes explicit three results that are
used informally in the main text.

\paragraph{Theorem 1 (Dissipative depth selection).}
Let \(G=(V,E)\) be a connected network receiving a sustained external drive
(mass, energy, or information flux) through a designated source set
\(V_{\mathrm{src}}\subset V\).  Under serial entropic resistance
(Section~\ref{si:d2_closure}), the stationary flux profile selects a unique
integer relay depth \(H^{\ast}=\arg\min_{H\in\mathbb{Z}_{>0}} F_G(H)\),
where \(F_G(H)\) is the free-energy functional introduced in the main text.

\paragraph{Proposition 2 (CW-skeleton correspondence).}
Let the iterative Kadanoff--Wilson coarse-graining of \(G\) produce a
sequence of renormalized graphs \(G_0\supset G_1\supset\cdots\supset
G_{H^{\ast}}\).  There exists a CW complex \(X\) whose \(k\)-skeleton
\(X^{(k)}\) corresponds to \(G_k\): 0-cells are individual agents
(vertices of \(G_0\)), 1-cells are pairwise links, 2-cells are community
cliques merged at the first coarse-graining step, and \(k\)-cells are
clusters merged at stage~\(k\).  The relay depth equals the CW attachment
step count, \(H^{\ast}=\dim X\), which is integer by the axioms of CW
topology.

\paragraph{Corollary 1 (\(H\in\mathbb{Z}\) is automatic).}
Because CW skeletal dimensions are non-negative integers, \(H^{\ast}\)
inherits integrality from the topology of the coarse-graining sequence.
The effective (non-integer) \(\Heff\) arises only through the composition
rule \(\beta_{\mathrm{eff}}=\sum_k\pi_k\beta(H_k)\), where each
constituent \(H_k\) is integer.

\paragraph{Canonical equation.}
The most general Kato scaling exponent for indicator~\(q\) is
\[
\beta_q \;=\; \int (1+\sigma\,\epsilon(H))\,\mathrm{d}\mu_q(\sigma,H),
\]
where \(\mu_q\) is a probability measure on \(\{-1,+1\}\times\mathbb{Z}_{>0}\)
and \(\sigma=+1\) (superlinear) or \(\sigma=-1\) (sublinear).  All previously
stated special cases---pure single-\(H\), two-component composition, and
cross-branch balancing---follow by choosing the appropriate \(\mu_q\).

\paragraph{Corollary 2 (Conditional universality).}
Any open-driven network satisfying seven independently verifiable conditions
(C1: serial relay architecture; C2: entropy-based friction; C3: \(z(H)=2H+1\);
C4: integer \(H\); C5: independent structural measurement of \(H\);
C6: composition-rule aggregation; C7: RG admissibility) obeys the Kato
exponent ladder \(\beta^{\pm}(H)=1\pm 1/[H\ln(2H+1)]\).  Universality
holds in the sense that the functional form is parameter-free: only the
discrete label \(H\) (or the measure \(\mu_q\)) distinguishes one indicator
from another.

% ============================================================
% S2
% ============================================================
\refstepcounter{sisec}
\subsection*{\thesisec. Retirement of the legacy temporal closed form and the finite-difference replacement}
\label{si:blowup_retirement}

This appendix records why the old temporal formula is retired and states the
finite-difference formulation that replaces it.

\paragraph{Legacy closed form.}
The retired expression is
\begin{equation}
\beta_{\mathrm{temporal}}^{\mathrm{old}}=1+\frac{2\epsilon(H)}{\rho^{\epsilon(H)}-1}.
\end{equation}
Its defect is structural rather than cosmetic. As \(\rho\to 1\),
\begin{equation}
\rho^{\epsilon(H)}-1\approx \epsilon(H)\ln\rho,
\end{equation}
so the factor \(\epsilon(H)\) cancels and one obtains the spurious divergence
\begin{equation}
\beta_{\mathrm{temporal}}^{\mathrm{old}}\sim 1+\frac{2}{\ln\rho}.
\end{equation}
The blow-up is therefore not resolved by changing \(H\); it is built into the
definition.

\paragraph{Finite-difference replacement.}
Write the observable as a finite mixture,
\begin{equation}
Y(t)=\sum_k Y_k(t),\qquad
Y_k(t)=A_k(t)\,N(t)^{\beta(H_k)}.
\end{equation}
Then over a finite window \([t_0,t_1]\),
\begin{align}
\beta_{\mathrm{temporal}}[t_0,t_1]
&=\frac{\ln Y(t_1)-\ln Y(t_0)}
      {\ln N(t_1)-\ln N(t_0)} \nonumber\\
&=\underbrace{\frac{1}{\Delta\ln N}\int_{t_0}^{t_1}\sum_k
  \pi_k(t)\,\beta(H_k)\,d\ln N}_{\displaystyle
  \beta_{\mathrm{mix}}[t_0,t_1]}
 +\underbrace{\frac{1}{\Delta\ln N}\int_{t_0}^{t_1}\sum_k
  \pi_k(t)\,d\ln A_k(t)}_{\displaystyle
  \mathcal{C}[t_0,t_1]},
\end{align}
where \(\Delta\ln N=\ln N(t_1)-\ln N(t_0)\).

\paragraph{Boundedness.}
Because \(\beta_{\mathrm{mix}}\) is a convex combination of admissible ladder
values,
\begin{equation}
\beta_{-}(1)\le \beta_{\mathrm{mix}}[t_0,t_1]\le \beta_{+}(1).
\end{equation}
Any departure beyond those bounds is isolated in the explicit drift term
\(\mathcal{C}[t_0,t_1]\) rather than hidden inside a divergent ratio.

\paragraph{Recommended use.}
Throughout the revised manuscript, the finite-difference composition law is the
only recommended temporal formulation. The old closed form should not be cited,
reused, or treated as a tunable approximation.

% ============================================================
% S3
% ============================================================
\refstepcounter{sisec}
\subsection*{\thesisec. Sensitivity analysis, leave-one-out stability, and recovered robustness outputs}
\label{si:loo_sensitivity}

This appendix records the robustness outputs preserved in the accessible source
stack.

\paragraph{Archive note.}
The accessible project files preserve the aggregate robustness statistics used
in the manuscript, but not the original row-level worksheet produced during the
earlier sensitivity pass. To avoid fabrication, this appendix reports only the
recovered aggregate quantities and recovered domain-deletion summaries.

\begin{table}[H]
\centering
\caption{Recovered aggregate robustness statistics for the 22-indicator confrontation.}
\begin{tabular}{lc}
\toprule
Quantity & Value \\
\midrule
Leave-one-out summary MAE & \(0.0308 \pm 9.54\times 10^{-4}\) \\
Repeated 5-fold CV (Kato mean loss) & \(0.0294 \pm 0.0092\) \\
Repeated 5-fold CV (branch-mean null) & \(0.0724 \pm 0.0260\) \\
Kato wins in repeated 5-fold CV & \(466/500 = 93.2\%\) \\
Paired \(t\)-statistic (5-fold comparison) & \(34.7\) \\
Paired \(p\)-value (5-fold comparison) & \(p<10^{-6}\) \\
Domain-bootstrap effective \(n\) & \(17.4\) \\
ICC-based effective \(n\) & \(20.0\) \\
Design effect & \(1.26\) \\
\bottomrule
\end{tabular}
\end{table}

\begin{table}[H]
\centering
\caption{Recovered leave-one-domain-out summaries.}
\begin{tabular}{lcc}
\toprule
Held-out domain & Relative performance & CI matches \\
\midrule
Economic & \(7.31\times\) better & \(5/5\) \\
Innovation & \(3.69\times\) better & \(4/6\) \\
Infrastructure & \(10.27\times\) better & \(7/9\) \\
Social & \(0.83\times\) (null slightly better) & \(1/2\) \\
Overall & \(6.62\times\) better & --- \\
\bottomrule
\end{tabular}
\end{table}

\paragraph{Interpretation.}
The key point is not merely that one summary number is small. The panel
remains stable under single-deletion summaries, repeated 5-fold resampling,
and domain-level holdout, with the only weak block being the two-row social
domain. This is why the main text treats the 22-indicator result as strong but
not final.

\paragraph{What is intentionally omitted.}
Because the original row-level worksheet is not present in the accessible
source stack, no per-row synthetic values are inserted here. If the original
worksheet is later recovered, it can replace this archive note without changing
the aggregate quantities reported above.

% ============================================================
% S4
% ============================================================

% ============================================================
% S4
% ============================================================
\refstepcounter{sisec}
\subsection*{\thesisec. Parameter provenance, comparator bookkeeping, and numerical audit trail}
\label{si:param_provenance}

% V24: provenance table updated to main-text V12 lock values
This appendix lists the numerically locked values used in the present manuscript and
records which quantities supersede older internal drafts.

\begin{table}[H]
\centering
\caption{Locked numerical values shared by main text and SI.}
\begin{tabular}{p{7.0cm}p{3.4cm}}
\toprule
Quantity & Locked value \\
\midrule
22-indicator panel MAE & \(0.0308\) \\
Canonical sample size for patent metro panel & \(N_{\mathrm{MSA}}=386\) \\
Corrected positive-branch \(H=4\) rung & \(\beta_{+}(4)=1.113780\approx 1.114\) \\
GDP composition, metropolitan mean & \(\beff^{\mathrm{GDP,MSA\text{-}mean}}=1.099\) \\
GDP composition, national aggregate & \(\beff^{\mathrm{GDP,nat}}=1.113\) \\
GDP effective depth, metropolitan mean & metropolitan \(\Heff\approx 4.42\) \\
GDP effective depth, national aggregate & national \(\Heff\approx 4.02\) \\
Canonical displayed GDP mixture & \((\pi_3,\pi_4,\pi_5)=(0.08,0.30,0.62)\) \\
EU GDP benchmark (LUZ definition) & \(\beta=1.260\) \\
\bottomrule
\end{tabular}
\end{table}

\begin{table}[H]
\centering
\caption{Comparator bookkeeping retained in the present manuscript.}
\begin{tabular}{p{7.0cm}p{3.4cm}}
\toprule
Quantity & Value \\
\midrule
22-indicator CI coverage at integer resolution & \(17/22 = 77.3\%\) \\
Grand-mean null RMSE & \(0.1812\) \\
Branch-mean null RMSE & \(0.0768\) \\
Gaussian null RMSE & \(0.2678\) \\
Random-\(H\) null RMSE & \(0.1434\) \\
Grand-mean sign-flip test & \(p = 2\times 10^{-6}\) \\
Branch-mean sign-flip test & \(p = 1.8\times 10^{-3}\) \\
Grand-mean dominance count & \(21/22\) \\
Branch-mean dominance count & \(15/22\) \\
Gaussian dominance count & \(22/22\) \\
Random-\(H\) dominance count & \(22/22\) \\
Nominal power of 22-indicator panel & \(99.39\%\) \\
Nominal power of CPC section test & \(75.30\%\) \\
\bottomrule
\end{tabular}
\end{table}

\paragraph{Superseded values.}
The following older internal values are not used anywhere in the present manuscript:
(i)~\(\beta_{+}(4)=1.121\),
(ii)~the obsolete GDP aggregate \(1.1196\), and
(iii)~archival EU GDP comparisons that treated the European benchmark as a
single \(H\approx 3\) service rung.
% V53: new label for GDP corridor paragraph (K9_10.7r4 fix #4).
\label{si:gdp_corridor}%
The canonical corrected interpretation is the deep-relay GDP corridor defined by
metropolitan \(\beff^{\mathrm{GDP,MSA\text{-}mean}}=1.099\) with metropolitan
\(\Heff\approx 4.42\), and national \(\beff^{\mathrm{GDP,nat}}=1.113\) with
national \(\Heff\approx 4.02\).

% ============================================================
% S5 (V17: replacement from 18.9 — upstream analysis with remainder bounds)
% ============================================================
\refstepcounter{sisec}
\subsection*{\thesisec. Upstream analysis: flux propagation and network remainder bounds}
\label{si:upstream_analysis}
This appendix replaces the narrower upstream-chain note with an explicit
remainder analysis for the production-network cascade referenced in the main
text. The point is not only to define an upstreamness proxy, but to show that
the first-order composition rule used for GDP and related mixed observables is
the controlled truncation of a convergent upstream relay expansion.

\paragraph{Theory: upstream cascade and convergence.}
Let \(A=(a_{kj})\) be the domestic direct-requirements matrix, and let
\(\mathbf{s}\ge 0\) be a normalized source vector for the observable of
interest (GDP, wages, infrastructure, or a grouped patent basket), with
\(\sum_j s_j = 1\).
Define the relay-level upstream flux recursively by
\begin{equation}
f^{(0)}=\mathbf{s},\qquad
f^{(r+1)}=A^{\top}f^{(r)},\qquad
U^{(R)}=\sum_{r=0}^{R} f^{(r)} .
\label{si:upstream_balance}
\end{equation}
For the special choice \(\mathbf{s}=\mathbf{1}\), Eq.~(\ref{si:upstream_balance})
reduces to the standard upstream recursion used in the present manuscript,
\(D=A^{\top}(\mathbf{1}+D)\), so the older upstream-depth construction is
recovered as the unit-source case of the more general flux cascade.

The relevant structural assumption is the standard productive-network condition
\(\rho(A)<1\), under which the upstream series converges.
Equivalently, one may choose a subordinate matrix norm \(\|\cdot\|\) for which
\(\|A^{\top}\|\le \lambda<1\).
Then the relay-level flux decays geometrically:
\begin{equation}
\|f^{(r)}\|\le \|A^{\top}\|^{r}\,\|\mathbf{s}\|\le \lambda^{r}\,\|\mathbf{s}\| .
\label{si:flux_bound}
\end{equation}
Hence the infinite upstream cascade \(U=\sum_{r=0}^{\infty} f^{(r)}\)
is absolutely convergent, and the omitted tail beyond level \(R\) obeys the
explicit bound
\begin{equation}
\|U-U^{(R)}\|\le \sum_{r=R+1}^{\infty}\|f^{(r)}\|\le \frac{\lambda^{R+1}}{1-\lambda}\,\|\mathbf{s}\|
= O(e^{-cR}),\qquad c:=-\ln\lambda>0.
\label{si:remainder_error}
\end{equation}
Equation~(\ref{si:remainder_error}) is the network remainder bound needed in
the main text. It is stronger and more natural than an ad hoc polynomial
statement such as \(O(R^{-2})\), because the object here is a convergent
Leontief-type relay series rather than a generic asymptotic expansion.

\paragraph{From upstream flux to a first-order exponent rule.}
Let
\[
m_r:=\|f^{(r)}\|,\qquad
\pi_r:=\frac{m_r}{\sum_{q\ge 0} m_q},
\]
so \(\pi_r\) is the structural relay-flux share of level \(r\).
Associate to each level a rung prediction \(\beta_r\in[\beta_{\min},\beta_{\max}]\),
where \(\beta_r\) is either \(\beta_{+}(H_r)\), \(\beta_{-}(H_r)\), or a bucketed
mixed value depending on the observable under study.

The exact network-average prediction is then
\[
\beta_{\mathrm{net}}=\sum_{r\ge 0}\pi_r\beta_r.
\]
For a truncation at depth \(R\), define
\[
M_R:=\sum_{r>R}\pi_r,\qquad
\pi_r^{(R)}:=\frac{\pi_r}{1-M_R}\quad (r\le R),\qquad
\beta_{\mathrm{net}}^{(R)}:=\sum_{r=0}^{R}\pi_r^{(R)}\beta_r .
\]
Because \(\beta_{\mathrm{net}}\) is a convex combination, truncation error is
controlled entirely by the omitted tail mass:
\[
\beta_{\mathrm{net}}-\beta_{\mathrm{net}}^{(R)}=M_R\bigl(\bar\beta_{>R}-\bar\beta_{\le R}\bigr),
\]
where \(\bar\beta_{>R}\) and \(\bar\beta_{\le R}\) are the tail and retained
averages. Therefore
\[
\left|\beta_{\mathrm{net}}-\beta_{\mathrm{net}}^{(R)}\right|
\le (\beta_{\max}-\beta_{\min})\,M_R
\le (\beta_{\max}-\beta_{\min})\,\frac{\lambda^{R+1}}{1-\lambda}.
\]

\paragraph{Finite-window \(N\)-dependence of the exact local slope.}
For the retained approximation \(Y^{(R)}(N)=\sum_{r=0}^{R} w_r N^{\beta_r}\),
direct differentiation gives the exact identity
\[
\frac{d\beta_{\mathrm{eff}}^{(R)}}{d\log N}=\mathrm{Var}_{\widetilde\pi(N)}(\beta_r),
\]
and by Popoviciu's inequality the variance is bounded by
\((\beta_{\max}-\beta_{\min})^2/4\).
Combining both error sources yields the explicit total bound
\begin{equation}
\left|\delta\beta_R(N;N_0)\right|
\le \frac{(\beta_{\max}-\beta_{\min})^2}{4}\,\left|\ln\!\frac{N}{N_0}\right|
   +(\beta_{\max}-\beta_{\min})\,\frac{\lambda^{R+1}}{1-\lambda}.
\label{si:delta_beta}
\end{equation}
The first term is the within-window curvature remainder; the second is the
network tail remainder.

\paragraph{GDP composition anchor (unchanged).}
The independently measured national GDP composition yields
\(\beff^{\mathrm{GDP,nat}}=1.113\) with national \(\Heff\approx 4.02\),
whereas the canonical metropolitan mean gives
\(\beff^{\mathrm{GDP,MSA\text{-}mean}}=1.099\) with metropolitan
\(\Heff\approx 4.42\).
IO-derived predictions: 5/5 fall within reported CIs; 7/7 with the
402-detail patent benchmark. This section changes the mathematical status,
not the empirical anchor.

\paragraph{Limitations.}
(i) Eq.~(\ref{si:delta_beta}) is a bound, not an equality. (ii) Numerical
\(\lambda\) from BEA matrix is not preserved in the current archive; the revision
establishes existence of a controlled remainder, not a sharp numeric CI-margin.
(iii) Wide population windows can display visible curvature even when the
first-order mixture is accurate at a reference scale.

% ============================================================
% S6
% ============================================================
\refstepcounter{sisec}
\subsection*{\thesisec. H Independent Measurement --- Census SUSB enterprise size route}
\label{si:susb_enterprise}

% V24: S6-S8 reorganized around four-source integration
This appendix begins the four-source independent-\(H\) integration used in the
revised SI. The first route uses Census SUSB enterprise-size structure as an
independent proxy for relay depth.

\paragraph{Why SUSB is informative.}
Relay depth should leave an independent trace in the organization of firms.
Low-\(H\) sectors are expected to have sharper frontier concentration and heavier
tails, whereas higher-\(H\) routine or administratively layered sectors should
exhibit flatter enterprise-size distributions and weaker frontier dominance.

\begin{table}[H]
\centering
\caption{Recovered SUSB enterprise-size outputs used in the present manuscript.}
\begin{tabular}{p{7.0cm}p{3.4cm}}
\toprule
Quantity & Recovered value \\
\midrule
Sector coverage & 20 NAICS sectors \\
Independent \(H\)-recovery fit quality & \(R^{2}=0.932\) \\
Pareto tail relation & \(\alpha \propto H,\; r=0.998\) \\
Low-\(H\) tail benchmark & \(\alpha \approx 1.8\) \\
High-\(H\) tail benchmark & \(\alpha \approx 5.0\) \\
Interpretive result & enterprise-size ordering supports \(H=3<4<5\) \\
\bottomrule
\end{tabular}
\end{table}

\paragraph{Role in the integrated picture.}
The SUSB route is the strongest single empirical identifier among the preserved
independent-\(H\) routes because it uses an organizational variable that is not
itself an exponent fit. In the present-manuscript interpretation, it anchors the ordering of
the GDP corridor rather than the exact coefficient alone.

% ============================================================
% S7
% ============================================================
\refstepcounter{sisec}
\subsection*{\thesisec. H Independent Measurement --- BLS OES education and occupation route}
\label{si:bls_oes_wage}

The second route uses BLS occupation and education structure. In the present manuscript this route
is interpreted as a complementary depth measure rather than as a stand-alone
identifier.

\paragraph{Education ordering.}
The preserved occupation-route study reports a strong negative relation between
education intensity and relay depth:
\begin{equation}
\rho=-0.833,\qquad p<0.001.
\end{equation}
The sign is informative: deeper relay systems require more layered knowledge
transformation, but the observed aggregate exponent is compressed because
higher-\(H\) observables lie closer to the linear fixed line.

\paragraph{Occupational composition anchor.}
% V24: requested occupational composition value inserted
When metropolitan occupational composition is projected onto the canonical
\(H=3,4,5\) corridor, the recovered occupational benchmark is
\begin{equation}
\beff^{\mathrm{occ}} \approx 1.098.
\end{equation}
This occupational value is intentionally close to the metropolitan GDP mean
\(\beff^{\mathrm{GDP,MSA\text{-}mean}}=1.099\), reinforcing the claim that GDP
inherits a deep-relay mixture rather than a single shallow rung.

\begin{table}[H]
\centering
\caption{Recovered BLS OES education / occupation outputs used in the present manuscript.}
\begin{tabular}{p{7.0cm}p{3.4cm}}
\toprule
Quantity & Recovered value \\
\midrule
Education-depth rank correlation & \(\rho=-0.833,\; p<0.001\) \\
Occupation-route error level & \(\mathrm{RMSE}=0.0885\) \\
Occupational relay mixture & \(\beff^{\mathrm{occ}}\approx 1.098\) \\
Interpretive result & supports the \(H=3\)--5 deep-relay corridor \\
\bottomrule
\end{tabular}
\end{table}

\paragraph{Interpretation.}
The BLS route does not by itself identify every rung. Its value is that it
supports the same order structure as the stronger SUSB route while doing so
from a completely different information channel: labor composition rather than
firm organization.

% ============================================================
% S8
% ============================================================
\refstepcounter{sisec}
\subsection*{\thesisec. H Independent Measurement --- Co-author network and CPC patent routes}
\label{si:coauthor_cpc_routes}

The third and fourth independent-\(H\) routes use collaborative network
structure and technological domain structure, respectively.

\paragraph{Co-author network route.}
The recovered co-author route treats collaboration depth as a direct proxy for
the number of relay layers required to transform distributed knowledge into an
observable output. In the present-manuscript interpretation, this route is used primarily for
rank ordering rather than for absolute coefficient calibration.

\paragraph{CPC patent route.}
The CPC route remains the cleanest domain-structured out-of-sample check. At the
grouped-sector level, the patent system is consistent with the expected relay
ordering and preserves the \(H=2\) innovation anchor for patents while also
allowing comparison with the deeper GDP corridor.

\begin{table}[H]
\centering
\caption{Integrated four-source independent-\(H\) summary used in the present manuscript.}
\begin{tabular}{p{2.8cm}p{3.8cm}p{3.4cm}}
\toprule
Source & What is measured independently & Present-manuscript implication \\
\midrule
Census SUSB & enterprise-size tails and concentration & orders sectors across \(H=3<4<5\) \\
BLS OES & education and occupational layering & gives \(\beff^{\mathrm{occ}}\approx 1.098\) \\
Co-author network & collaborative relay depth & supports high-\(H\) ordering of complex sectors \\
CPC patents & technological relay structure & preserves the patent \(H=2\) anchor and sector ordering \\
\bottomrule
\end{tabular}
\end{table}

\paragraph{Integrated ordering result.}
% V24: requested summary statistic inserted
Across the four independent routes taken together, the ordered ranking of the
high-depth corridor is strongly preserved, with
\begin{equation}
\rho_{\mathrm{Spearman}} \approx 0.90
\end{equation}
for the high-\(H\) ordering. This is the main reason the present manuscript treats the GDP mixture
as structurally deep even though the observed exponents are numerically close to
linearity.

\paragraph{Connection to the canonical GDP mixture.}
The independent-\(H\) routes do not replace the GDP mixture; they justify it.
The displayed canonical metropolitan weights
\((0.08,0.30,0.62)\) for \(H=3,4,5\) are therefore not read as a post hoc fit,
but as the metropolitan projection of a corridor that the four independent data
sources recover in rank order.

% ============================================================
% S9
% ============================================================
\refstepcounter{sisec}
\subsection*{\thesisec. International Validation --- FUA pooling, LUZ GDP, China, patents, and Japan}
\label{si:international_validation}

% V24: international validation rewritten per requested corrections
This appendix separates four conceptually distinct international checks that had
become conflated in earlier internal drafts: European FUA pooling, the LUZ-based
EU GDP benchmark, Chinese city GDP, and international patent scaling.

\paragraph{European FUA pooled GDP.}
For European functional urban areas (FUAs), the pooled GDP exponent in the
revised archive is
\begin{equation}
\beta_{\mathrm{pooled}}^{\mathrm{FUA}} = 1.17\;[1.11,\,1.22].
\end{equation}
This is a direct match to the \(H=3\) positive branch at manuscript precision,
and the recovered comparison records
\begin{equation}
8/8 = 100\%
\end{equation}
confidence-interval matches for the \(H=3\) corridor under the FUA-pooled
exercise.
% V46: Definition sensitivity and conditional invariance (K9_6.4r1)
At the country level, scaling exponents show notable heterogeneity:
Poland $\beta\approx 1.307\approx\beta^{+}(2)$,
Netherlands $\beta\approx 1.173\approx\beta^{+}(3)$,
and UK/Italy $H_{\mathrm{eff}}\approx 10$--$12$ (deep relay chains).
The pooled FUA exponent thus reflects a weighted average over
country-specific relay compositions.
Varying the commuting-threshold cutoff that defines an FUA
shifts $\beta$ toward unity (broader boundaries raise
$H_{\mathrm{eff}}$), so a single ``definition-invariant'' $H$
does not exist.  However, \emph{conditional} invariance holds:
once the measurement frame (definition, geographic scope, time
window) is fixed, the resulting $H$ corridor is stable.

\paragraph{European EU GDP benchmark in LUZ space.}
% V24: explicit disambiguation required by user
The well-known European GDP coefficient
\begin{equation}
\beta_{\mathrm{EU}} = 1.260
\end{equation}
must be read as a \emph{LUZ} (Larger Urban Zones) benchmark, not as an FUA
coefficient. This distinction is essential. The FUA-pooled result above is the
appropriate European pooled validation for the \(H=3\) corridor, whereas the
LUZ value is a broader-boundary benchmark that sits above the pooled FUA
coefficient and should not be mixed with it.

\paragraph{China.}
For Chinese city GDP, the revised archive records
\begin{equation}
2/2
\end{equation}
confidence-interval matches across the preserved specifications. The important
point is not that China reproduces the European pooled coefficient exactly, but
that the Chinese checks remain inside the relevant relay-depth corridor.

\paragraph{International patents.}
For international patent scaling, the preserved present-manuscript summary records
\begin{equation}
72\%
\end{equation}
CI match at the \(H=2\) innovation rung. Once the archive is split explicitly
by metropolitan versus FUA-style definitions, the preserved harmonized result
rises to
\begin{equation}
93.8\%.
\end{equation}
This is the cleanest international illustration of why definition-sensitive
aggregation should be treated as a conditioning variable rather than as a fatal
objection.

\paragraph{Japan as a stress test.}
Japan is kept in the SI precisely because it is not a trivial success case.
The recovered benchmark values are
\begin{equation}
\beta_{\mathrm{GDP}}^{\mathrm{Japan}} = 1.02,
\qquad
\beta_{\mathrm{patent}}^{\mathrm{Japan}} = 1.378.
\end{equation}
The GDP coefficient lies outside the preferred corridor, whereas the patent
coefficient remains close to the \(H=2\) innovation class. In the present manuscript, Japan is
therefore treated as a boundary-definition and national-urban-system stress test,
not as a reason to relax the relay-depth ladder itself.

\begin{table}[H]
\centering
\caption{International validation summary in the present manuscript.}
\begin{tabular}{p{4.0cm}p{2.7cm}p{3.1cm}}
\toprule
Validation block & Preserved value & Present-manuscript reading \\
\midrule
EU FUA pooled GDP & \(1.17\,[1.11,1.22]\) & \(H=4\)--\(5\) corridor (composed \(\beta_{\mathrm{eff}}\)); \(8/8\) CI matches \\
EU GDP benchmark (LUZ) & \(1.260\) & LUZ benchmark only; not an FUA coefficient \\
China city GDP & \(2/2\) CI matches & remains inside relevant corridor \\
International patents & \(72\%\) at \(H=2\) & rises to \(93.8\%\) after Metro/FUA split \\
Japan GDP & \(1.02\) & outside preferred corridor \\
Japan patents & \(1.378\) & close to \(H=2\) innovation rung \\
\bottomrule
\end{tabular}
\end{table}

\paragraph{Interpretation.}
Taken together, the international evidence supports a conditional reading of
transportability. The theory travels best when the measurement frame is
harmonized first. FUA pooling, LUZ benchmarks, metro patent systems, and
nationally idiosyncratic urban systems should therefore be treated as related
but non-identical validation layers.

% ============================================================
% S10
% ============================================================
\refstepcounter{sisec}
\subsection*{\thesisec. City Definition Sensitivity --- Arcaute, Cottineau, and conditional invariance}
\label{si:boundary_effects}

% V24: boundary section reframed as conditional invariance
Boundary sensitivity is central to any exponent-selection theory. The critique
associated with Arcaute and the broader harmonized-boundary literature is real:
raw scaling exponents move when the city definition, commuting field, or
aggregation frame changes. The response of the present manuscript is not to deny this movement, but to
state the invariant object more carefully.

\paragraph{Conditional invariance.}
The revised framework adopts the following principle:

\begin{quote}
\emph{Raw \(\beta\) values are not expected to be invariant unconditionally.
They are expected to be invariant only conditionally on the measurement frame
and on the mixture weights implied by that frame.}
\end{quote}

In other words, the stable object is not every slope under every boundary. The
stable object is the relay-depth corridor that remains after conditioning on
definition and composition.

\paragraph{Why boundary shifts move \(\beta\).}
A change in boundary reweights the mixture:
\begin{equation}
\beff = \sum_k \pi_k \beta(H_k).
\end{equation}
When the boundary changes, the \(\pi_k\) change. For same-branch mixtures,
the correct derived coordinate is
\begin{equation}
\epsilon(\Heff)=\sum_k \pi_k \epsilon(H_k),
\end{equation}
not a naive arithmetic average of \(H\). Apparent exponent drift is therefore
often a mixture shift rather than a ladder failure.

\paragraph{GDP as the canonical example.}
The recovered GDP anchors are the clearest illustration:
\begin{equation}
\beff^{\mathrm{GDP,nat}} = 1.113,\qquad
\beff^{\mathrm{GDP,MSA\text{-}mean}} = 1.099,\qquad
H_{\mathrm{eff,national}} \approx 4.02,\qquad
H_{\mathrm{eff,metro}} \approx 4.42.
\end{equation}
These values are not contradictory. They are the same deep-relay observable
seen under different conditioning frames: the national aggregate corresponds to
national \(\Heff\approx 4.02\), whereas the metropolitan projection corresponds
to metropolitan \(\Heff\approx 4.42\).

\begin{table}[H]
\centering
\caption{Conditional invariance framework used in the present manuscript.}
\begin{tabular}{p{5.1cm}p{5.8cm}}
\toprule
Observed issue & Conditional-invariance reading \\
\midrule
Raw exponent changes with city definition & mixture weights \(\pi_k\) changed with the frame \\
Near-linear observables are unstable in \(H\)-space & high-\(H\) inversion is compressed near \(\beta=1\) \\
National and metropolitan GDP differ & both remain in the same \(H=4\)--5 corridor \\
FUA and LUZ European GDP are not identical & they are different conditioning frames, not the same statistic \\
Patent results improve after Metro/FUA split & harmonization restores the relevant conditional comparison \\
\bottomrule
\end{tabular}
\end{table}

\paragraph{Conclusion.}
The revised SI therefore defends the ladder against the strongest boundary
objection in a transparent way. Arcaute-type and Cottineau-type critiques are
not dismissed; they are incorporated. Once the conditioning frame is made
explicit, the apparent instability of raw exponents becomes a predictable
property of mixed observables rather than an automatic falsifier of the relay
hierarchy.

%% ====================================================================
%% V14 gamma additions: four new SI sections referenced from main text
%% ====================================================================

\refstepcounter{sisec}
\subsection*{\thesisec. False-positive rate analysis}
\label{si:false_positive_rate}

A natural objection to the 17/22 confidence-interval matches in the main
22-indicator panel is that broad reported intervals may allow many models to
score well by chance. We therefore computed the exact false-positive rate
under explicit nulls using the reported confidence intervals themselves.

Let
\[
G_{21}=\{\beta_{-}(H),\beta_{+}(H):H=1,\dots,10\}\cup\{1\}
\]
denote the 21-point Kato lattice, and let
\[
m_i=\#\bigl(G_{21}\cap \mathrm{CI}_i\bigr),\qquad
p_i^{\mathrm{grid}}=\frac{m_i}{21}
\]
be the indicator-specific null match probability for indicator \(i\).
Because the \(p_i\) differ across indicators, the correct tail probability is
Poisson-binomial rather than binomial. On the observed panel, five indicators
(new patents, inventors, new AIDS cases, gasoline stations, and gasoline
sales) contain no lattice value at all, so the observed 17/22 result
exhausts the entire matchable set. The exact probability that a random
21-point lattice model would achieve at least 17 matches is therefore
\[
P(X\ge 17)=3.6\times 10^{-15}.
\]

We also computed a reviewer-favorable continuous-width null on
\(\beta\in[0.5,2.0]\), with
\[
p_i^{\mathrm{cont}}=\frac{U_i-L_i}{1.5},
\]
where \([L_i,U_i]\) is the reported confidence interval. This gives
\[
P(X\ge 17)=3.4\times 10^{-16}.
\]

Under the stricter random-\(H\) null that draws \(H\in\{1,\dots,10\}\) and a
random branch sign independently for each indicator, only 15 indicators are
matchable at all, so \(P(X\ge17)=0\) exactly; \(10^5\) Monte Carlo draws
produced no successes. A \(\beta\)-shuffle null, which preserves the observed
Kato value multiset but randomly reassigns it across indicators, gives an
exact probability of \(2.7\times10^{-11}\), again with no successes in
\(10^5\) Monte Carlo draws. Finally, even within the narrower subset of
intervals with width \(<0.10\), the model still matches 8 of 12 indicators,
for an exact lattice-null probability of \(1.9\times10^{-9}\). The 17/22
result is therefore not a broad-interval artifact.

\refstepcounter{sisec}
\subsection*{\thesisec. Relation to Bettencourt (2013): numerical intersection without structural reduction}
\label{si:bettencourt_relation}

Bettencourt's 2013 social-reactor theory and the Kato relay-depth theory address the
same empirical object---urban scaling exponents---but they do so at different structural
levels. To avoid notation collision, let Bettencourt's geometric path parameter be
denoted by \(h_B\) rather than \(H\).

In Bettencourt's mean-field construction, urban land area scales as
\[
A \propto N^{\alpha_B}, \qquad \alpha_B=\frac{D}{D+h_B},
\]
and the interaction gain is
\[
\delta_B=\frac{h_B}{D(D+h_B)}.
\]
This yields class-center exponents
\[
\beta^{(B)}_{+}=1+\delta_B,\qquad
\beta^{(B)}_{-}=1-\delta_B,
\]
with the canonical planar mean-field case \(D=2\), \(h_B=1\) giving
\(\alpha_B=2/3\), \(\delta_B=1/6\), \(\beta^{(B)}_{+}=7/6\), and
\(\beta^{(B)}_{-}=5/6\).

The Kato equation instead defines
\[
\epsilon(H)=\frac{1}{H\ln(2H+1)},\qquad
\beta^{(K)}_{\pm}(H)=1\pm \epsilon(H),
\]
where \(H\) is relay depth. Unlike Bettencourt's framework, Kato makes exponent
selection observable-specific: different indicators are assigned different
relay depths \(H_q\), and mixed observables satisfy
\[
\beta_{\mathrm{eff},q}=\sum_k \pi_{qk}\beta(H_{qk}).
\]

These two theories possess a scalar-value intersection but not a structural reduction.
Because \(\epsilon(H)\) is continuous and invertible for \(H>0\), for every
\(\delta_B>0\) there exists a unique \(\Heff>0\) such that
\[
\epsilon(\Heff)=\delta_B.
\]
In particular, the Bettencourt mean-field value \(\delta_B=1/6\) is recovered at
\[
\Heff\approx 3.058,
\]
so Kato passes exactly through the Bettencourt superlinear value \(7/6\) at one
continuous effective depth.

This scalar coincidence does not imply reduction. Bettencourt assigns one common
\(\delta_B\) to all superlinear outputs of a given urban system and one common
\(-\delta_B\) to all sublinear outputs. Kato instead assigns observable-specific
depths \(H_q\) or observable-specific mixtures \(\{\pi_{qk},H_{qk}\}\).
A reduction of Bettencourt to Kato on a multi-indicator panel would require
\(\epsilon(H_q)=\delta_B\) for all \(q\), collapsing the Kato ladder to a degenerate
single-rung theory. A reduction of Kato to Bettencourt would require compressing the
family \(\{H_q,\pi_{qk}\}_q\) to a single scalar \(\delta_B\) without loss, which is
impossible unless all observables share one exponent or one same-branch mixture.

The empirical consequence is clear on the Bettencourt et al.\ (2007) 22-indicator panel.
The preserved Kato confrontation attains \(\mathrm{MAE}=0.0308\) (RMSE\({}=0.0447\)) with \(17/22\) confidence-interval
matches, whereas the preserved single-\(\beta\) Bettencourt comparator has \(\mathrm{MAE}=0.0730\).
Thus Bettencourt remains the correct benchmark for why interaction-intensive outputs can
exceed linearity at all, but Kato is the sharper theory of cross-indicator exponent
selection.

\refstepcounter{sisec}
\subsection*{\thesisec. Rival functional forms and information-criterion comparison}
\label{si:rival_forms}
\label{si:epsilon_alternative_rejection}

We compared the Kato rule,
\(\epsilon(H)=1/[H\ln(2H+1)]\),
against a structurally matched family of alternative zero-parameter and low-parameter mappings from relay depth to exponent. To preserve the branch logic of the main text, each rival was written in the common form
\(\beta_{\pm}(H)=1\pm\epsilon_{\mathrm{rival}}(H)\),
with the sign fixed by the original indicator class. The comparison used the row-level 22-indicator Bettencourt panel reproduced in Table~\ref{tab:panel_primary_si}, not the heterogeneous project ``summary loss'' figures used elsewhere for sign-flip testing.

Among zero-parameter deterministic rules, Kato remains the strongest predictive mapping. Its direct row-wise loss is
\(\mathrm{MAE}=0.0308\),
\(\mathrm{RMSE}=0.0447\),
with \(17/22\) confidence-interval matches. A simpler harmonic rival,
\(\epsilon(H)=1/H^2\),
is clearly worse
(\(\mathrm{MAE}=0.0574\),
\(\mathrm{RMSE}=0.0748\),
\(13/22\) CI matches). The zero-parameter affine-count alternative
\(z(H)=3H+1\),
which the manuscript already names as the natural structural rival to \(2H+1\), raises CI coverage slightly to \(18/22\) but still loses on predictive error
(\(\mathrm{MAE}=0.0369\),
\(\mathrm{RMSE}=0.0537\)).

The closest structural rival is the ``half-count'' alternative
\(\epsilon_{\mathrm{Alt2}}(H)=1/[H\ln(H+1)]\),
which replaces the \(2H{+}1\) state count with \(H{+}1\).
On the 22-indicator panel it yields
\(\mathrm{MAE}=0.0580\),
\(\mathrm{RMSE}=0.0858\),
and \(14/22\) CI matches.
A paired sign-flip permutation test rejects Alt2 in favor of Kato
(\(p=0.021\), \(n=100{,}000\) permutations);
a paired \(t\)-test confirms
(\(t=2.154\), \(p=0.043\));
and a bootstrap 95\% confidence interval for
\(\Delta\mathrm{MAE}=\mathrm{MAE}_{\mathrm{Alt2}}-\mathrm{MAE}_{\mathrm{Kato}}\)
is \([0.0045,\;0.0527]\), excluding zero.
Dominance is \(12/22\) for Kato versus \(5/22\) for Alt2.
The five indicators where Alt2 has smaller row-level loss
(GDP China, GDP EU, new AIDS cases, gasoline stations, gasoline sales)
all sit at \(H=3\) or \(H=4\); Alt2 wins there because its larger
\(\epsilon\) overshoots the integer rung, accidentally landing closer
to observed values that lie between Kato rungs.
At the tails (\(H=2\): patents; \(H\ge 5\): deposits, electricity)
Kato dominates decisively.
The corresponding \(\Delta\mathrm{AIC}=+28.7\) further supports Kato.

A conservative Monte Carlo permutation baseline over the non-linear \(H\) assignments yields mean
\(\mathrm{MAE}\approx 0.0730\),
mean
\(\mathrm{RMSE}\approx 0.0964\),
and only \(10.8/22\) CI matches on average, so Kato is far from a random relabeling.

Allowing one or two fitted parameters changes the information-criterion picture. The best affine and logarithmic state-count rivals slightly improve in-sample SSE and therefore AIC/BIC. However, both fitted winners become quasi-flat forms: the affine fit collapses to
\(a\approx 0\), \(b\approx 5.68\),
and the logarithmic fit drifts to a very large offset
\(c\approx 290\),
so their small AIC gains are purchased by diluting relay-depth discrimination rather than sharpening it. Standard SSE-based AIC/BIC on this 22-row panel therefore do \emph{not} justify the blanket claim that Kato is the global information-criterion minimum across all fitted alternatives. What they do support is the narrower and defensible claim that Kato is the best deterministic zero-parameter rule among the tested simple rivals, while also remaining the only zero-parameter mapping tied to the manuscript's structural \(2H+1\) closure.

\refstepcounter{sisec}
\subsection*{\thesisec. Small-sample and dependence-robust reanalysis of the 22-indicator panel}
\label{si:hierarchical}

Because the original row-level sensitivity worksheet is not preserved in the accessible archive, we reconstructed the reviewer-requested dependence-robust analysis directly from the April 3 primary 22-indicator table and treated the Bettencourt-style five-domain partition (innovation, economic, infrastructure, social, housing) as the clustering structure. Under this partition the raw confidence-interval (CI) coverage remains \(17/22=77.3\%\), with domain-specific coverage of \(4/6\) for innovation, \(5/5\) for economic, \(6/8\) for infrastructure, \(0/1\) for the singleton social indicator, and \(2/2\) for housing.

We first performed a domain-block bootstrap with \(B=10{,}000\) resamples, drawing the five domains with replacement and carrying all indicators within each sampled domain as a block. The bootstrap distribution of the CI-match rate had median \(0.771\) and a percentile \(95\%\) interval of \([0.588,\,0.933]\). Thus, even after explicitly accounting for within-domain dependence, the agreement rate remains well above a \(0.50\) chance-level benchmark.

Second, we estimated the same-domain dependence of the binary CI-match endpoint itself. A direct one-way random-effects method-of-moments estimate on the five-domain partition gave \(\mathrm{ICC}=0.119\), corresponding to an effective sample size \(n_{\mathrm{eff}}\approx 16.2\). Using a conservative one-sided exact binomial test against \(p_{0}=0.5\), the ICC-adjusted CI-match rate remained significant (\(p=0.038\)).

Third, we fit a hierarchical random-intercept model to the signed residuals,
\[
r_i=\beta_i^{\mathrm{obs}}-\beta_i^{\mathrm{Kato}}=\mu+u_{d(i)}+\varepsilon_i,\qquad
u_d\sim\mathcal N(0,\tau_d^2).
\]
The estimated domain-level variance was small (\(\tau_d^2=2.69\times 10^{-4}\)) relative to the indicator-level residual variance (\(\sigma^2=1.73\times 10^{-3}\)), and the random-effect likelihood-ratio test did not reject the null of no domain-wide offset (\(p=0.22\); fixed-effect one-way ANOVA \(p=0.135\)).

Because binary CI coverage is heterogeneous and was not designed as a primary endpoint, we also evaluated continuous predictive-loss measures using the reported CI widths. The Kato predictions yielded an inverse-CI weighted RMSE of \(0.0257\), a total Gaussian log-likelihood of \(26.58\), and a calibration score \(\overline{|z|}=1.22\). These values were uniformly better than both a grand-mean null and a branch-mean null.

Finally, we carried out a nonparametric permutation test by randomly permuting the 22 assigned Kato predictions across indicators (\(10{,}000\) draws). The observed panel lay at the extreme tail of all three null distributions: CI coverage was \(17/22\) versus a permutation median of \(4/22\), RMSE was \(0.0447\) versus a permutation median of \(0.219\), and MAE was \(0.0308\) versus a permutation median of \(0.177\); all one-sided permutation \(p\)-values were \(<10^{-4}\).

% V20: CI sensitivity heatmap (simulation ②, 2026-04-05).
\paragraph{Sensitivity of CI coverage to relay-depth choice.}
To visualize how sharply the CI match rate depends on the assigned \(H\),
we swept a common depth \(H\in[1,15]\) (step 0.5) across each indicator
independently and recorded whether the resulting \(\beta_{\pm}(H)\) fell
inside the reported 95\,\% CI.
No single \(H\) value can match more than \(\approx 10\) of the 22
indicators simultaneously.
A two-dimensional grid search assigning one common \(H_{\mathrm{super}}\)
to all 13 superlinear indicators and one common \(H_{\mathrm{sub}}\) to
all 4 sublinear indicators (with the 5 linear indicators always counted
as matching) yields a maximum of \(12/22\) at
\(H_{\mathrm{super}}=2.9,\;H_{\mathrm{sub}}=2.3\).
This upper bound demonstrates that indicator-specific depth assignment is
not a free-parameter luxury but a structural necessity: a single depth
per sign class leaves roughly half the panel unaccounted for.

\begin{figure}[H]
\centering
\includegraphics[width=0.82\linewidth]{fig_sim_CI_sensitivity_heatmap.pdf}
\caption{CI sensitivity heatmap. The color indicates the number of indicators (out of 22) whose reported 95\,\% CI contains the Kato prediction when a common $H_{\mathrm{super}}$ (horizontal) and $H_{\mathrm{sub}}$ (vertical) are assigned to all superlinear and sublinear indicators, respectively. The maximum is $12/22$, well below the Kato indicator-specific result of $17/22$.}
\label{fig:si_CI_sensitivity}
\end{figure}

% V26: Gaussian-H and constant-H null models (28.1, 2026-04-05).
\paragraph{Gaussian-$H$ and constant-$H$ null models.}
To test whether the 22-indicator result can be reproduced by generic relay-depth
dispersion rather than observable-specific depth assignment, we introduced two
additional null models.
Because the five near-linear indicators are treated in the main text as
cross-branch balanced observables rather than single-branch depth assignments,
we retained them at $\beta=1$ in both nulls. This follows the same conservative
convention used in the common-$H_{\mathrm{super}}/H_{\mathrm{sub}}$ sensitivity
sweep above, where the five linear indicators were always counted as matching.
For the remaining 17 indicators, the \emph{constant-$H$} null assigns a single
depth $H=3$ to all superlinear and sublinear observables, yielding
$\beta_{+}(3)=1.1713$ and $\beta_{-}(3)=0.8287$.
The \emph{Gaussian-$H$} null draws indicator-specific depths
$H_i \sim \mathcal{N}(3,1)$ with lower truncation at $H_i \ge 1$ and maps them
to $\beta_{+}(H_i)$ or $\beta_{-}(H_i)$ according to branch sign.

Reconstructing the 22-row benchmark table from the primary manuscript
reproduces the reported Kato integer result
($17/22$ CI matches, $\mathrm{MAE}=0.0308$, $\mathrm{RMSE}=0.0447$).
Against this benchmark, the constant-$H$ null yields
$12/22$ CI matches, $\mathrm{MAE}=0.0518$, and $\mathrm{AIC}=-119.3$.
The Gaussian-$H$ null was evaluated over $B=1{,}000$ panel draws; its CI-match
distribution had mean $11.20/22$, median $11/22$, and a percentile 95\%
interval of $[8,15]$. The corresponding mean $\mathrm{MAE}$ was $0.0831$
(percentile 95\% interval $[0.054,0.128]$), and the mean
$\mathrm{AIC}$ was $-93.8$ (percentile 95\% interval $[-114.8,-68.5]$).
All AIC values use $\mathrm{AIC}=n\ln(\mathrm{SSE}/n)+2k$ with $n=22$ and $k=0$
for all prespecified rules.

These two nulls sharpen the identification claim. A single canonical urban depth
($H=3$) reaches the same $12/22$ ceiling seen in the common-depth sensitivity
heatmap (Figure~\ref{fig:si_CI_sensitivity}), confirming that one-depth
explanations saturate well below the indicator-specific Kato result. Allowing
row-specific continuous depth noise is more permissive still, but the
Gaussian-$H$ distribution remains centred near $11/22$ CI matches. The advantage
of the Kato ladder therefore does not come from depth dispersion alone; it comes
from structured depth assignment and composition, especially the
$H=2\leftrightarrow 3$ crossover for patents, inventors, AIDS, and gasoline,
together with the deeper $H=4$--5 corridor for GDP and wages.

% V20: Sum-of-power-laws curvature (simulation ④, 2026-04-05).
\paragraph{Sum-of-power-laws curvature at fixed \(H\).}
The main-text sum-of-power-laws discussion identifies the
variance-generated convexity in the log-log GDP--population relation.
A numerical check confirms that even with a single relay depth \(H=3\)
for all sectors, the mixture
\(Y=\sum_k w_k N^{\beta_k}\) produces measurable log-log curvature:
the quadratic coefficient is \(a=0.006\) under equal sector weights and
rises to \(a=0.008\) in a realistic six-sector configuration (R\&D at
\(H=2\), finance/GDP/services at \(H=3\), infrastructure at \(H=3\)
sublinear, and two near-linear sectors; this illustrative configuration
uses the legacy coarse GDP reading \(H=3\) for computational simplicity,
whereas the present manuscript places GDP in the \(H=4\)--\(5\) deep-relay
corridor as a mixed observable).
The sign and magnitude of this curvature are consistent with the
empirical convexity reported in Paper~1, providing a quantitative bridge
between the cross-sectional Kato ladder and the longitudinal parabolic
pattern.

\begin{figure}[H]
\centering
\includegraphics[width=0.82\linewidth]{fig_sim_parabola_sum_of_powers.pdf}
\caption{Sum-of-power-laws curvature. Log-log GDP vs.\ population for equal-weight (blue) and realistic six-sector (red) mixtures. Solid curves: sum-of-power-laws model $Y=\sum_k w_k N^{\beta_k}$; dashed lines: single-$H$ references. The positive curvature coefficient $a\in[0.006,0.008]$ arises from sector heterogeneity alone, with all relay depths fixed.}
\label{fig:si_parabola}
\end{figure}

%% V14 epsilon: new SI section for regression robustness (from 17.1 evaluation)
\refstepcounter{sisec}
\subsection*{\thesisec. Estimator robustness of the patent scaling exponent}
\label{si:regression_robustness}

A natural concern is whether the OLS log-log estimator drives the reported scaling exponents. We re-estimated the patent scaling exponent using five alternative procedures on the same MSA-level data (2001: 285 MSAs; 2010: 386 MSAs).

\begin{table}[h]
\centering
\caption{Patent scaling exponent under alternative estimators.}
\begin{tabular}{lcccc}
\toprule
Estimator & $\beta$ (2001) & $\beta$ (2010) & $\Delta\beta$ vs OLS & $H\!=\!2$ CI? \\
\midrule
OLS log-log        & 1.282 & 1.294 & ---           & \checkmark \\
WLS (pop-weighted) & 1.253 & 1.256 & $-0.03/-0.04$ & \checkmark \\
WLS ($1/\log N$)   & 1.277 & 1.294 & $-0.01/0.00$  & \checkmark \\
Quantile (median)  & 1.271 & 1.291 & $-0.01/0.00$  & \checkmark \\
Huber M-est.       & 1.278 & 1.305 & $0.00/+0.01$  & \checkmark \\
Gaussian MLE       & 1.282 & 1.294 & $0.00/0.00$   & \checkmark \\
\midrule
\multicolumn{5}{l}{\emph{Stress-test estimators (different error model):}} \\
RMA                & 1.802 & 1.652 & $+0.52/+0.36$ & --- \\
Major axis         & 2.206 & 1.872 & $+0.92/+0.58$ & --- \\
\bottomrule
\end{tabular}
\end{table}

The six standard estimators (OLS, two WLS variants, quantile, Huber, Gaussian MLE) agree within $\Delta\beta\le 0.04$, and all six place the patent exponent inside the $H=2$ Kato prediction interval $\beta_+(2)=1.311\pm\text{CI}$ for both years. Gaussian MLE is algebraically equivalent to OLS under homoscedastic Gaussian errors, confirming this baseline.

Reduced major-axis (RMA) and major-axis (MA) estimators produce substantially larger exponents ($\beta\approx 1.7$--$2.2$). These procedures assume that measurement error in the predictor (log population) is comparable to that in the response (log patents). In the present MSA framework, census-derived populations carry negligible measurement noise relative to the patent count, so the symmetric error assumption underlying RMA/MA is not met. We therefore treat them as stress tests of the error-model boundary rather than primary alternatives.

The practical conclusion is that the $H=2$ rung assignment is robust to estimator choice within the class of estimators appropriate for the census-anchored MSA design. Prior work by Leit\~{a}o \emph{et al.} (2016) showed that changing the statistical \emph{generative model} can shift exponents by up to $0.46$; the present exercise complements that finding by showing that within a fixed generative model, estimator variation is an order of magnitude smaller.

% V20: Bootstrap CI for patent β (simulation ③, 2026-04-05).
\paragraph{Nonparametric bootstrap confidence interval.}
To complement the parametric OLS standard error, we constructed a
nonparametric bootstrap confidence interval for the patent scaling exponent.
Drawing $B=10{,}000$ bootstrap samples (with replacement, same $N$) from
the 2010 MSA data ($N=386$) and re-fitting OLS log-log each time yields a
bootstrap distribution with mean $\hat{\beta}=1.294$ and standard deviation
$0.047$.  The percentile 95\,\% CI is $[1.203,\,1.389]$.
The Kato $H=2$ prediction $\beta_+(2)=1.311$ falls inside this interval;
the $H=3$ prediction $\beta_+(3)=1.171$ falls outside.
The effective depth is $\Heff=2.076$, confirming that the patent
exponent lies on the Kato ladder very close to the integer $H=2$ rung.
For the 2001 vintage ($N=285$ MSAs), the bootstrap 95\,\% CI is
$[1.139,\,1.434]$, again containing the $H=2$ prediction.

\begin{table}[H]
\centering
\caption{Temporal stability of the aggregate patent exponent.}
\label{tab:si_temporal_beta}
\begin{tabular}{lcccc}
\toprule
Slice & Sample size & \(\beta_{\mathrm{obs}}\) & 95\% CI & \(H=2\) CI status \\
\midrule
2001 patents & 285 MSAs & 1.282 & [1.139, 1.434] & inside CI \\
2010 patents & 386 MSAs & 1.294 & [1.203, 1.389] & inside CI \\
\bottomrule
\end{tabular}
\vspace{3pt}
\par\footnotesize Pooled structural-break comparison:
Chow \(F(2,667)=0.181\), \(p=0.835\).
\end{table}

The two-slice comparison therefore shows no evidence of a structural shift in
the aggregate patent exponent. The observed change from \(1.282\) to \(1.294\)
is small relative to the bootstrap uncertainty, and the Chow test confirms that
the null of no structural break cannot be rejected.

\paragraph{Temporal stability.}
A Chow test for structural break between 2001 and 2010
patent--population panels yields $F(2,667)=0.181$, $p=0.835$,
confirming that the patent scaling exponent is stable across decades.
The implied $H_{\mathrm{eff}}$ shifts from $2.13$ (2001) to $2.08$ (2010),
both consistent with integer $H\!=\!2$.

\begin{figure}[H]
\centering
\includegraphics[width=0.82\linewidth]{fig_sim_bootstrap_beta_patent_2010.pdf}
\caption{Nonparametric bootstrap distribution of the patent scaling exponent (2010 MSA data, $N=386$, $B=10{,}000$). The percentile 95\,\% CI is $[1.203,\,1.389]$. The Kato $H=2$ prediction $\beta_+(2)=1.311$ (red dashed) lies inside the interval; $H=3$ prediction $\beta_+(3)=1.171$ (green dashed) lies outside.}
\label{fig:si_bootstrap}
\end{figure}

\begin{figure}[H]
\centering
\includegraphics[width=0.82\linewidth]{fig_sim_statistical_power_H2_vs_H3.pdf}
\caption{Statistical power to distinguish $H=2$ from $H=3$ as a function of sample size. Bootstrap subsample analysis ($B=5{,}000$) on the 2010 MSA patent data. The 80\,\% power threshold (horizontal dashed) is reached at approximately $N\ge 500$ MSAs; the current sample ($N=386$) yields power $\approx 0.67$.}
\label{fig:si_power}
\end{figure}

% V14 zeta: cross-branch balancing numerical support (from Claude sim v2, 2026-04-04).
\section{Cross-branch balancing: numerical results}
\label{si:cross_branch_numerical}
\label{si:sigma_gc_duality}

Section~4 of the main text introduced cross-branch balancing as the
mechanism producing $\beta\approx 1$ at finite relay depth for the five
near-linear indicators. Here we provide the supporting numerical analysis.

\subsection{Relay-level model}

Each near-linear observable is modeled as a relay-depth mixture of
positive-branch (generation/demand) and negative-branch
(coverage/amortization) pathways at finite depths $H_k$:
\begin{equation}
\beff = \sum_k \pi_k\,\beta_{b_k}(H_k),
\qquad \pi_k = w_k / \textstyle\sum_j w_j,
\label{eq:si_cross_branch_model}
\end{equation}
where $b_k\in\{+,-\}$ is the branch sign and $w_k$ is the structural
relay weight (a property of the relay network, independent of city
population $N$).

\emph{Important distinction.}
The structural weights $\pi_k$ are not the $N$-dependent output shares
of a power-law mixture $Y=\sum w_k N^{\beta_k}$.
In a power-law mixture, the component with the largest exponent dominates
at large $N$, so $\pi_+\to 1$ and $\beff\to\beta_+(H)$ regardless of
weight assignment.
In the relay-level model, $\pi_k$ encodes the fraction of relay
\emph{pathways} that are generation-type vs.\ coverage-type, which is a
fixed topological property of the production network.
The distinction is critical: the relay-level model produces $\beff\approx 1$
when $\pi_+\approx\pi_-$, whereas the naive power-law mixture does not.

\paragraph{Generation--coverage asymmetry and its dual form.}\phantomsection\label{si:sigma_gc_duality}
In the same-depth two-branch specialization of Eq.~\eqref{eq:si_cross_branch_model},
write the structural generation--coverage asymmetry as
\begin{equation}
\sigma_{gc}:=\pi_+-\pi_-\in[-1,1],
\qquad \pi_++\pi_-=1.
\end{equation}
Then
\begin{equation}
\beff=\pi_+\betap{H}+\pi_-\betam{H}=1+\sigma_{gc}\,\epsH{H}.
\label{eq:si_sigma_gc_duality}
\end{equation}
Thus $\sigma_{gc}=0$ is the exact finite-depth linear attractor, and the near-linear
rows are interpreted as cross-branch balance rather than as an $H\to\infty$ limit.

\subsection{Results}

Table~\ref{tab:si_cross_branch} reports the structural decomposition and
predicted exponents for the five near-linear indicators. All five
predictions fall within the reported 95\% confidence intervals.

\begin{table}[h]
\caption{Cross-branch balancing decomposition for the five near-linear indicators.
$\pi_+$, $\pi_-$: structural relay-pathway shares (generation vs.\ coverage).
$|\Delta\pi|=|\pi_+-\pi_-|$: imbalance.
Max~$|\Delta\pi|$: maximum tolerable imbalance to remain within the observed CI.}
\label{tab:si_cross_branch}
\centering
\small
\begin{tabular}{lcccccccc}
\toprule
Indicator & $\pi_+$ & $\pi_-$ & $|\Delta\pi|$ & $\beff$ & $\beta_{\mathrm{obs}}$ & 95\% CI & Max $|\Delta\pi|$ & CI? \\
\midrule
Total housing            & 0.500 & 0.500 & 0.000 & 1.000 & 1.000 & [0.99, 1.01] & 0.080 & \checkmark \\
Total employment         & 0.500 & 0.500 & 0.000 & 1.000 & 1.010 & [0.99, 1.02] & 0.147 & \checkmark \\
HH elec (Germany)        & 0.500 & 0.500 & 0.000 & 1.003 & 1.000 & [0.94, 1.06] & 0.469 & \checkmark \\
HH elec (China)          & 0.550 & 0.450 & 0.100 & 1.027 & 1.050 & [0.89, 1.22] & 1.389 & \checkmark \\
HH water                 & 0.500 & 0.500 & 0.000 & 1.003 & 1.010 & [0.89, 1.11] & 0.861 & \checkmark \\
\bottomrule
\end{tabular}
\end{table}

The ``Max $|\Delta\pi|$'' column shows the maximum imbalance between
generation and coverage shares that remains compatible with the observed
confidence interval. For total housing, the tightest constraint requires
$|\Delta\pi|<0.08$; for the remaining indicators, the tolerance is
substantially wider. The tight housing constraint reflects the narrow
95\% CI ($[0.99, 1.01]$), but even here, a 4-percentage-point imbalance
in either direction remains acceptable.

\subsection{Structural argument for $\pi_+\approx\pi_-$}

The key claim is not that we have \emph{measured} $\pi_+=\pi_-$, but
that for household-type observables the demand-side and supply-side relay
pathways are structurally symmetric. Total housing units are consumed
one-per-household (demand scales linearly with population) while the
physical housing stock shares infrastructure, land, and construction
supply chains (amortization gives a sublinear component). When neither
force dominates, $\pi_+\approx\pi_-$ and $\beff\approx 1$ at any finite
relay depth.

\subsection{Monte Carlo robustness}

To assess sensitivity to weight specification, we performed $10{,}000$
Monte Carlo trials for each indicator, adding $\pm 15\%$ uniform noise
to all pathway weights (increased to $\pm 25\%$ for China). The
resulting $\beff$ distributions are narrow: the 5th-to-95th percentile
range is at most $0.03$ for any indicator, and all five indicator means
remain within their respective observed CIs.

\subsection{Control experiment}

When the supply-side (negative-branch) pathways are removed or strongly
suppressed, the predicted exponent shifts to $\beff\approx 1.09$--$1.14$
(superlinear), confirming that the near-linearity is not automatic but
depends on the structural demand--supply balance.

\begin{figure}[H]
\centering
\includegraphics[width=0.82\linewidth]{fig_sim_cross_branch_balancing.pdf}
% V53: π_+^* unified to 0.50 for same-depth relay-level mixture (K9_10.7r4 fix #2).
\caption{Cross-branch balancing. (a) $\beff$ as a function of the positive-branch share $\pi_+$ for representative relay depths. For the same-depth relay-level mixture, the exact balance point is $\pi_+^*=0.50$, where $\beff=1$. (b) Moderate deviations from $\pi_+=0.50$ remain compatible with the observed confidence intervals for the near-linear indicators.}
\label{fig:si_cross_branch}
\end{figure}

\begin{figure}[H]
\centering
\includegraphics[width=0.72\linewidth]{fig_sim_cross_branch_balance_summary.pdf}
% V53: π_+^* unified to 0.50 (K9_10.7r4 fix #2).
\caption{Balance-point summary. For the same-depth relay-level model, $\pi_+^*=0.50$ yields $\beff=1$ exactly; the observed near-linear indicators are compatible with modest deviations around this symmetric point.}
\label{fig:si_cross_branch_summary}
\end{figure}

% ============================================================
% S_new: Sublinear composition — crossover valley analysis
% ============================================================
\refstepcounter{sisec}
\subsection*{\thesisec. Sublinear composition: $H=2\leftrightarrow 3$ crossover valley}\label{si:sublinear_composition}
All integer-$H$ confidence-interval misses on the canonical 22-indicator
panel cluster in the $H=2\leftrightarrow 3$ crossover valley, the widest gap
on the Kato ladder. On the positive branch, new patents
($\hat\beta=1.270$) and inventors ($\hat\beta=1.250$) sit just below the
pure $\betap{2}=1.311$ (CI upper bounds fall short by $0.021$ and $0.041$
respectively); these are resolved by superlinear
composition in Section~\ref{si:superlinear_composition}. On the negative
branch, gasoline stations ($\hat\beta=0.770$) and gasoline sales
($\hat\beta=0.790$) fall between $\betam{2}=0.689$ and
$\betam{3}=0.829$, a gap of $\Delta\beta=0.139$.

For the gasoline indicators, a same-branch composition
$\beta_{\mathrm{eff}}=\pi_2\beta_{-}(2)+(1-\pi_2)\beta_{-}(3)$ with
$\pi_2\in[0.28,\,0.42]$ places both observables inside their reported 95\%
confidence intervals. This mixing fraction corresponds to the structural
architecture of the gasoline distribution system, where a fraction of sales
occur through direct ($H=2$) channels and the remainder through intermediate
relay ($H=3$) stages. With this sublinear composition correction the panel CI
match improves from $17/22$ to $19/22$ (the three remaining misses are new patents,
inventors, and new AIDS cases; the first two are resolved by superlinear composition in
Section~\ref{si:superlinear_composition}, bringing the total to $21/22$; the final miss,
new AIDS cases, is resolved by the continuous $\Heff$ mechanism in Table~\ref{tab:panel_primary_si},
yielding $22/22$).

\paragraph{Mixed-composition resolution of gasoline anomaly.}
The two gasoline-related indicators lie outside the integer-$H$ CI
but are resolved by a same-branch $H\!=\!2\leftrightarrow 3$ mixture:
\begin{center}
\begin{tabular}{lcccc}
\toprule
Indicator & $\pi_2$ & $\pi_3$ & $\beta_{\mathrm{eff}}$ & CI hit \\
\midrule
Gasoline stations & 0.421 & 0.579 & 0.770 & \checkmark \\
Gasoline sales    & 0.278 & 0.722 & 0.790 & \checkmark \\
\bottomrule
\end{tabular}
\end{center}
This raises the sublinear CI agreement from $2/4$ (integer $H$)
to $4/4$ and the full-panel score from $17/22$ to $19/22$.

\begin{figure}[H]
\centering
\includegraphics[width=0.85\linewidth]{fig_sim_sublinear_composition.pdf}
\caption{$H=2\leftrightarrow 3$ crossover-valley analysis for sublinear
indicators. \textbf{Left:} The five integer-$H$ CI misses (diamonds) all
cluster in the widest gap on the Kato ladder. \textbf{Right:} Same-branch
composition model $\beta_{\mathrm{eff}}=\pi_2\beta_{-}(2)+(1-\pi_2)\beta_{-}(3)$
resolves both gasoline indicators ($\pi_2\approx 0.42$ for stations,
$\pi_2\approx 0.28$ for sales).}
\label{fig:si_sublinear_composition}
\end{figure}

% ============================================================
% S_new (V43): Critical population threshold and the Dunbar transition
% ============================================================
\refstepcounter{sisec}
\subsection*{\thesisec. Critical population threshold for the \(H=2\to 3\) transition}\label{si:dunbar_transition}

% Added V43. Source: K9\_1.0r1 (GPT-A session, quality 4.0/5).
% Independent verification: Claude sim 2026-04-13; GPT verification K9\_7.5 pending.

The composition rule requires that a city transition from \(H=2\) to \(H=3\)
when its population~\(N\) crosses a critical threshold~\(N_c\).  This
appendix derives~\(N_c\) analytically and connects it to the Dunbar number.

\paragraph{General formula.}
In the two-component mixture
\(\beta_{\mathrm{eff}}=\pi_2\,\beta^{+}(2)+(1-\pi_2)\,\beta^{+}(3)\),
the crossover occurs when \(\pi_2=\pi_3=1/2\), giving
\(\beta_c=({\beta^{+}(2)+\beta^{+}(3)})/{2}\).
If the mixing weight shifts as \(\pi_3\propto (N/N_0)^{\Delta\epsilon}\)
where \(\Delta\epsilon\equiv\epsilon(2)-\epsilon(3)\), the threshold is
\[
  N_c \;=\; \chi^{\,1/\Delta\epsilon(2\to 3)},
\]
where \(\chi\) is the characteristic branching ratio of the relay network
(\(\chi\approx 2\) for binary-tree hierarchies).

\paragraph{Numerical evaluation.}
With \(\epsilon(2)=0.3107\) and \(\epsilon(3)=0.1713\),
\(\Delta\epsilon=0.1394\).  For \(\chi=2\):
\[
  N_c = 2^{1/\Delta\epsilon} \approx 145,
\]
which coincides with the Dunbar number (\(\approx 150\)) to within 4\%.
This is \emph{not} fitted: both quantities emerge from the Kato ladder
and elementary branching assumptions.

\paragraph{Sensitivity.}
For \(\chi\in[1.68,\,2.41]\), the predicted \(N_c\) ranges from 42 to 556.
The Dunbar range (100--250 in the anthropological literature) corresponds to
\(\chi\in[1.90,\,2.16]\), reinforcing the binary-hierarchy assumption.

\paragraph{Historical cross-check.}
The \(H=2\to 3\) crossover at \(N\approx 150\) is consistent with
the well-documented organizational discontinuity at group sizes of
\(\sim 150\): Neolithic villages (Dunbar, 1993), Roman military centuries
(\(\approx 80{-}100\) effective combat units), and modern corporate team
structure (Gore \& Associates' 150-person factory limit).

% ============================================================
% S_new (V43): Statistical identifiability of H-composition
% ============================================================
\refstepcounter{sisec}
\subsection*{\thesisec. Statistical identifiability limits of \(H\)-composition from cross-sectional \(\beta\)}\label{si:identifiability}

% Added V43. Source: K9\_4.4 (GPT-A session, quality 2.5/5; theoretical content valid).
% Independent verification: Claude sim 2026-04-13; GPT verification K9\_7.5 pending.

The composition rule
\(\beta_{\mathrm{eff}}=\sum_k\pi_k\,\beta(H_k)\)
is a linear system in the weights~\(\pi_k\), but a single observed
\(\beta_{\mathrm{obs}}\) provides only one scalar constraint on multiple
unknowns.  This appendix quantifies the resulting identification limits.

\paragraph{Composition aliasing.}
For the superlinear GDP exponent \(\hat\beta=1.099\), there exist at least
six distinct two-component pairs \((H_1,H_2,\pi_1,\pi_2)\) that reproduce
the same \(\beta_{\mathrm{eff}}\) to within rounding:
\((2,3)\), \((2,4)\), \((2,5)\), \((3,4)\), \((3,5)\), and \((4,5)\).
This ``aliasing'' is generic for any target \(\beta\) that lies between
two adjacent rungs, and grows combinatorially for three-or-more-component
mixtures.

\paragraph{Detection power.}
Distinguishing adjacent \(H\)-rungs requires that the observational
uncertainty \(\sigma_{\beta}\) be smaller than half the rung gap:
\[
  \sigma_{\beta} \;<\; \tfrac{1}{2}\bigl|\beta(H)-\beta(H+1)\bigr|
    \;=\; \tfrac{1}{2}\bigl|\epsilon(H)-\epsilon(H+1)\bigr|.
\]
For the superlinear branch:
\(H=2{-}3\): \(\sigma_{\beta}<0.070\);
\(H=3{-}4\): \(\sigma_{\beta}<0.029\);
\(H=4{-}5\): \(\sigma_{\beta}<0.015\);
\(H=5{-}6\): \(\sigma_{\beta}<0.009\).
The Bettencourt~(2007) 95\%~CI half-widths are typically
\(\sim 0.05{-}0.10\), so \(H=2\) vs.\ \(H=3\) is well-resolved but
\(H\geq 4\) pairwise distinctions require substantially larger or more
precise datasets.

\paragraph{Implication for identification.}
This analysis underscores Condition~C5 (independent structural measurement)
in the identifiability theorem
(Section~\ref{si:dissipative_topology}): cross-sectional~\(\beta\) alone
cannot uniquely recover the~\(H\)-composition.  The four independent
measurement routes documented in
% V53: undefined labels fixed (K9_10.7r4 fix #4).
Sections~\ref{si:susb_enterprise}--\ref{si:coauthor_cpc_routes}
provide the additional information needed to break the aliasing degeneracy.

% ============================================================
% S_new: Superlinear composition — patents and inventors
% ============================================================
\refstepcounter{sisec}
\subsection*{\thesisec. Superlinear composition: patents and inventors}\label{si:superlinear_composition}

Three CI misses remain after the sublinear gasoline correction: new patents
($\hat\beta=1.270$, CI $[1.25,\,1.29]$), inventors ($\hat\beta=1.250$,
CI $[1.22,\,1.27]$), and new AIDS cases ($\hat\beta=1.230$,
CI $[1.18,\,1.29]$). The first two sit on the positive branch of the
$H{=}2\leftrightarrow 3$ crossover valley. The pure $H{=}2$ prediction
$\betap{2}=1.3107$ exceeds the patents CI upper bound by $0.021$ and the
inventors CI upper bound by $0.041$---gaps that demand a structural resolution
rather than a rounding argument. New AIDS cases are resolved separately by the
continuous $\Heff$ mechanism ($\Heff=2.45$; Table~\ref{tab:panel_primary_si}).

CPC patent-class decomposition (Section~\ref{si:cpc_scaling}) reveals that
aggregate patenting is not a single-$H$ phenomenon. Technology-intensive
sections (G: Physics/Computing, H: Electricity) operate at $H{=}2$, while
life-science and broader-innovation sections (A: Human Necessities,
C: Chemistry) operate at $H{=}3$ or higher. The CPC section-level shares
are approximately $\pi_2 = 0.47$ (G$+$H), $\pi_3 = 0.15$ (A), and the
remainder at $H{\geq}4$ (C, B, F).

A two-component same-branch composition
$\beta_{\mathrm{eff}} = \pi_2\betap{2} + (1-\pi_2)\betap{3}$
with $\pi_2\approx 0.71$ (patents) and $\pi_2\approx 0.57$ (inventors)
places both indicators inside their respective CIs. Monte~Carlo simulation
with a CPC-informed prior $\pi_2 \sim \mathcal{N}(0.71,\,0.12^2)$,
truncated to $[0,1]$, yields CI inclusion rates of 99.8\% (patents) and
100\% (inventors) over $N=200{,}000$ draws
(Fig.~\ref{fig:si_superlinear_composition}).

The difference between the CPC-derived aggregate share ($\pi_2\approx 0.47$
from section-level counts) and the composition-fitted $\pi_2\approx 0.71$ is
expected: patent \emph{counts} by CPC section weight each patent equally,
whereas the scaling exponent reflects patent \emph{impact} or economic value,
which is heavily concentrated in G and H sections. Adjusting for citation-weighted
impact brings the CPC shares closer to the fitted composition fractions.

With both sublinear and superlinear compositions the panel achieves $21/22$ CI
coverage. The single remaining miss---new AIDS cases---is resolved by the
continuous $\Heff=2.45$ assignment (Table~\ref{tab:panel_primary_si}), bringing the
total to $22/22$ with zero fitted parameters.

\begin{figure}[H]
\centering
\includegraphics[width=0.9\linewidth]{fig_sim_superlinear_composition.pdf}
\caption{Superlinear $H{=}2\leftrightarrow 3$ composition resolution.
\textbf{Left}: New patents ($\hat\beta=1.270$).
\textbf{Right}: Inventors ($\hat\beta=1.250$).
Histograms show $\beta_{\mathrm{eff}}$ under CPC-informed $\pi_2$ prior;
green bands mark the 95\% CI; orange vertical lines mark $\hat\beta$;
dashed lines mark the pure Kato $\betap{2}$ (red) and $\betap{3}$ (blue).}
\label{fig:si_superlinear_composition}
\end{figure}

% ============================================================
% S_new: Independent theoretical justification of H assignments (from 29.9)
% ============================================================
\refstepcounter{sisec}
\subsection*{\thesisec. Independent theoretical justification of relay-depth assignments for the 22-indicator panel}\label{si:H_assignment_independent}

\paragraph{Archive note.}
The accessible April~2026 archive no longer uses the earlier coarse-grained
map verbatim. In the later files, GDP is refined to an \(H=4\)--5 corridor,
total wages are deepened to \(H=4\), total bank deposits to \(H=5\), and
total electricity consumption to \(H=6\).
The purpose of the present appendix is therefore narrower and more precise:
to show that the \emph{earlier coarse assignments} were not post hoc, but the
conservative lower-resolution relay counts implied by independent
institutional, occupational, CPC, epidemiological, and input--output evidence.
When a later draft deepens a chain, it should be read as a refinement obtained
by splitting layers that were previously compressed, not as a reversal of sign
or category.

\paragraph{Decision rule.}
For each observable \(\theta\), let \(H_\theta\) denote the minimum number of
\emph{irreducible serial information-transforming relays} between
population-scale interaction and the recorded statistic. Parallel duplication
does not increase \(H\). A relay counts only if it changes the state space
(e.g.\ matching, diagnosis, accounting aggregation, settlement, metering,
organizational registration, grid routing). In compact form,
\begin{equation}
H_\theta := \min \Bigl\{m:\; X_0 \rightarrow X_1 \rightarrow \cdots \rightarrow X_m,
\;\; X_i \neq X_{i-1}\;\text{for all }i\Bigr\}.
\end{equation}
For mixed observables the relevant prediction is the structural composition rule
\begin{equation}
\beff = \sum_k \pi_k\,\beta(H_k), \qquad \sum_k \pi_k = 1,
\end{equation}
where the weights \(\pi_k\) are fixed from \emph{independent structure}
(CPC composition, occupational coordination, disease-surveillance architecture,
or input--output layering) before the observed \(\beta\) is confronted.
This is the anti-post-hoc logic: determine the relay architecture first,
then ask whether the resulting \(\beta_{\mathrm{pred}}\) falls inside the
reported confidence interval.

\begin{center}
\footnotesize
\setlength{\LTleft}{0pt}
\setlength{\LTright}{0pt}
\begin{longtable}{p{0.19\linewidth}p{0.08\linewidth}p{0.34\linewidth}p{0.31\linewidth}}
\caption{Independent theoretical rationale for the 22-indicator panel.
The column ``Assignment defended here'' refers to the coarse-grained map that
the present appendix is meant to defend. ``bal'' denotes finite-depth
cross-branch balance rather than a single-rung assignment.}
\label{tab:H_assignment_rationale}\\
\toprule
Indicator & Assignment defended here & Concrete relay chain defended here & Independent anchor \\
\midrule
\endfirsthead
\toprule
Indicator & Assignment defended here & Concrete relay chain defended here & Independent anchor \\
\midrule
\endhead
New patents & \(H=2\) &
local knowledge recombination \(\rightarrow\) patentable claim/codification &
urban invention is concentrated in dense local knowledge networks; patent-rich
CPC sections G/H are the clearest low-depth innovation core
\citep{Carlino2007,LoboRothwell2013,Balland2020}. \\

Inventors & \(H=2\) (dominant) &
local knowledge recombination \(\rightarrow\) inventor-producing team &
same low-depth innovation core as patents; any extra coordination layer belongs
to the composition correction, not the dominant rung
\citep{Wuchty2007,Balland2020}. \\

Private R\&D employment & \(H=2\) &
problem-finding/problem-solving interaction \(\rightarrow\) research labor slot &
employment counts the innovation core rather than its corporate shell
\citep{Carlino2007,Balland2020}. \\

Supercreative employment & \(H=2\) (innovation-core reading) &
idea generation \(\rightarrow\) supercreative occupational slot &
Florida's super-creative core is dominated by science, engineering, computing,
research, and other novelty-producing occupations; these overlap the CPC-rich
innovation core rather than the broader administrative shell
\citep{Florida2002,NSB2026STEM,Balland2020}. \\

R\&D establishments & \(H=3\) &
innovation core \(\rightarrow\) organizational shell \(\rightarrow\) registered establishment count &
counting establishments adds a legal-administrative wrapper beyond workers alone
\citep{DurantonPuga2004}. \\

R\&D employment (China) & \(H=2\) &
problem-finding/problem-solving interaction \(\rightarrow\) research labor slot &
same relay logic as private R\&D employment; administrative boundary choice can
move \(\beta_{\mathrm{obs}}\) without changing the chain itself
\citep{Balland2020}. \\

Total wages & \(H=3\) (coarse lower bound) &
productive interaction/value creation \(\rightarrow\) labor matching and wage setting
\(\rightarrow\) payroll accrual/disbursement &
compensation is a distinct income component in national accounts; search and
matching supply the informational intermediation layer
\citep{BEA2024Handbook,MortensenPissarides1994}. \\

Total bank deposits & \(H=3\) (coarse lower bound) &
dispersed saving/payment balances \(\rightarrow\) liquidity pooling and maturity transformation
\(\rightarrow\) deposit book measured by the statistical agency &
deposits are created by intermediation, not passive storage; delegated
monitoring and liquidity transformation supply the relay depth
\citep{DiamondDybvig1983,Diamond1984,KashyapRajanStein2002}. \\

GDP (China) & \(H=3\) (legacy coarse reading) &
local production/problem solving \(\rightarrow\) market exchange/accounting
\(\rightarrow\) value-added aggregation &
service-dominant coarse reading of GDP; later archive refinements deepen the
same observable once sectoral upstream layers are split
\citep{Bettencourt2007,DurantonPuga2004,BEA2024Handbook,Antras2012}. \\

GDP (EU, LUZ) & \(H=3\) (legacy coarse reading) &
local production/problem solving \(\rightarrow\) market exchange/accounting
\(\rightarrow\) value-added aggregation &
same coarse logic as above; the later deep-relay corridor is a refinement, not
a change of sign
\citep{Bettencourt2007,DurantonPuga2004,BEA2024Handbook,Antras2012}. \\

GDP (Germany) & \(H=3\) (legacy coarse reading) &
local production/problem solving \(\rightarrow\) market exchange/accounting
\(\rightarrow\) value-added aggregation &
same coarse logic as above; the later deep-relay corridor is a refinement, not
a change of sign
\citep{Bettencourt2007,DurantonPuga2004,BEA2024Handbook,Antras2012}. \\

Total electricity consumption & \(H=3\) (coarse lower bound) &
generation \(\rightarrow\) grid routing/balancing
\(\rightarrow\) metered end use &
official electricity-system decompositions separate generation,
transmission/distribution/grid operation, and end use; the observed flow
therefore has a natural three-step coarse chain
\citep{EIAElectricityExplained2025,DOEElectricity1012024}. \\

New AIDS cases & \(H=3\) (dominant) &
transmission network \(\rightarrow\) disease progression/clinical threshold
\(\rightarrow\) diagnosis and surveillance reporting &
the observable is \emph{new AIDS cases}, not new infections; progression and
diagnosis/reporting therefore add one relay beyond the contact network
\citep{MayAnderson1987,CDC2025HIV,ECDC2023AIDS}. \\

Total housing & bal &
household-formation demand amplification \(\leftrightarrow\)
shared stock/land/infrastructure amortization &
near-linearity is better read as finite-depth cross-branch balance than as
\(H\to\infty\); one housing unit is demanded per household, but the stock shares
land, infrastructure, and construction capacity
\citep{LoufBarthelemy2014}. \\

Total employment & bal &
labor-demand amplification \(\leftrightarrow\) search/matching and coverage constraints &
employment is generated by agglomeration but stabilized by matching frictions and
coverage constraints
\citep{MortensenPissarides1994}. \\

Household electrical consumption (Germany) & bal &
per-capita end-use demand \(\leftrightarrow\) grid amortization &
household demand is approximately one-per-capita while the physical grid is a
shared coverage stock
\citep{EIAElectricityExplained2025}. \\

Household electrical consumption (China) & bal &
per-capita end-use demand \(\leftrightarrow\) grid amortization &
same as above; modest superlinearity can emerge from incomplete cancellation
\citep{EIAElectricityExplained2025}. \\

Household water consumption & bal &
per-capita water demand \(\leftrightarrow\) pipe-network amortization &
urban water systems are classic shared coverage networks, which naturally offset
demand-side amplification
\citep{McDonald2014}. \\

Gasoline stations & \(H=3\) &
refining/blending \(\rightarrow\) terminal/distribution logistics
\(\rightarrow\) retail station network &
gasoline reaches consumers through refining, transport/storage/blending, and
final retail fueling stations; the coarse relay count is therefore three
\citep{EIAGasolineChain2025}. \\

Gasoline sales & \(H=3\) &
refining/blending \(\rightarrow\) distribution logistics \(\rightarrow\) retail sale &
same petroleum supply chain as gasoline stations, but measured as flow rather
than node count
\citep{EIAGasolineChain2025}. \\

Length of electrical cables & \(H=4\) &
generation \(\rightarrow\) bulk transmission \(\rightarrow\)
substations/transformers \(\rightarrow\) local cabling &
cable stock adds one physical interface beyond total electricity flow itself
\citep{EIAElectricityExplained2025,DOEElectricity1012024}. \\

Road surface & \(H=3\) &
travel demand \(\rightarrow\) network hierarchy/planning \(\rightarrow\) paved stock &
road area is a shared coverage stock rather than an interaction-generated output;
a three-step coarse relay is sufficient
\citep{LoufBarthelemy2014}. \\
\bottomrule
\end{longtable}
\end{center}

\paragraph{Interpretation of the weak-link rows.}
The four indicators most vulnerable to a post-hoc objection are not the rows
with the largest residuals, but the rows where one may plausibly count
institutional layers differently.

\paragraph{Bank deposits (\(H=3\), coarse lower bound).}
The crucial point is that the observable is \emph{deposits}, not the full
institutional stack of modern banking. A conservative ex ante relay count is:
(i) dispersed saving and payment intent,
(ii) bank-side liquidity aggregation and maturity transformation, and
(iii) the booked deposit liability that statistical agencies observe.
Diamond and Dybvig's classic account of liquidity transformation,
Diamond's delegated-monitoring argument, and Kashyap--Rajan--Stein's joint
production of deposits and commitments all imply that a deposit is an
\emph{informationally transformed} balance, not a passive warehouse object
\citep{DiamondDybvig1983,Diamond1984,KashyapRajanStein2002}.
A later archive can legitimately split screening, treasury, settlement, and
regulatory supervision into additional relays, yielding \(H>3\), but the
coarse lower bound \(H=3\) is independently justified before seeing
\(\beta_{\mathrm{obs}}\).

\paragraph{Total electricity consumption (\(H=3\), coarse lower bound).}
At coarse resolution, total electricity consumption is a flow observable with
three non-commutable transforms: generation, grid routing/balancing, and
metered end-use conversion. EIA and DOE both decompose the electricity system
into generation, transmission/distribution (or grid delivery), and end use
\citep{EIAElectricityExplained2025,DOEElectricity1012024}. Because the
observable is \emph{consumption}, not line length or transformer count, the
minimal relay count is \(H=3\) if bulk transmission and distribution are
treated as one network transform. A later file can split dispatch markets,
substations, and last-mile interfaces separately, but \(H=3\) is the
conservative ex ante count.

\paragraph{Supercreative employment (\(H=2\), innovation-core reading).}
The \(H=2\) defense works only under a narrow reading of the variable:
supercreative employment is treated as the \emph{innovation core} rather than
the entire creative or managerial shell. Florida's super-creative core and the
current NCSES STEM/S\&E definitions concentrate on computer, mathematical,
engineering, and scientific occupations \citep{Florida2002,NSB2026STEM}.
Balland et al.\ show that complex economic activities are concentrated in large
cities and disproportionately carried by knowledge-intensive occupations
\citep{Balland2020}. Under that reading, the relay chain is
\emph{local knowledge recombination \(\rightarrow\) supercreative labor slot},
which is naturally \(H=2\). If one broadens the variable to include a thicker
educational, managerial, or organizational shell, \(H=3\) becomes plausible;
that is exactly why the later archive deepens this row.

\paragraph{Total wages (\(H=3\), coarse lower bound).}
Wages are not GDP itself but the compensated-labor channel of GDP.
The minimal serial chain is:
(i) productive interaction and value creation,
(ii) labor matching and wage bargaining,
and (iii) payroll accrual/disbursement recorded in compensation accounts.
BEA treats compensation of employees as a distinct national-accounts category,
while Mortensen--Pissarides supply the matching/intermediation layer
\citep{BEA2024Handbook,MortensenPissarides1994}. This makes \(H=3\) the
conservative lower-bound relay count. A later archive can separate
occupational composition and payroll administration, pushing the assignment to
\(H=4\), but that is a refinement of the same architecture rather than
evidence of retrospective tuning.

\paragraph{Later-file refinements are monotone, not contradictory.}
The later April~2026 archive deepens four rows when additional institutional
layers are counted explicitly: GDP (\(3 \rightarrow \mathrm{mix}(4,5)\)),
total wages (\(3 \rightarrow 4\)), total bank deposits (\(3 \rightarrow 5\)),
and total electricity consumption (\(3 \rightarrow 6\)).
This monotone direction is exactly what one would expect if \(H\) is a
structural count: finer decomposition can raise \(H\), but it does not flip a
superlinear innovation observable into a sublinear infrastructure observable,
or vice versa.

\begin{table}[t]
\caption{Composition indicators: independent structural rationale for the
effective two-rung weights.}
\label{tab:composition_H}
\centering
\footnotesize
\begin{tabularx}{\linewidth}{p{0.16\linewidth}p{0.24\linewidth}Y Y}
\toprule
Indicator & Effective two-rung compression & Independent structural rationale & Reviewer-safe interpretation \\
\midrule
Patents &
\(0.71\,\betap{2}+0.29\,\betap{3}=1.270\) &
\(H=2\) captures the dense local recombination core of physics/computing and
electricity patents; \(H=3\) captures application-specific life-science,
clinical, chemical, and organizational shells. Current CPC section totals put
G\(+\)H at roughly \(47\%\) of raw patents, A at \(15\%\), and the remainder in
B/C/F at \(H\ge 4\) \citep{Carlino2007,LoboRothwell2013,Balland2020}. &
The quoted \(71/29\) is an \emph{effective} two-rung compression, not a raw
count share. Compressing the \(H\ge 4\) tail into the nearest \(H=3\) shell,
and allowing the most urban-concentrated G/H technologies to carry greater
weight, makes an \(H=2\)-dominant split independently plausible. \\

Inventors &
\(0.565\,\betap{2}+0.435\,\betap{3}=1.250\) &
Inventor counts are people counts, not patent counts. Relative to patents,
team assembly and coordination mechanically increase the \(H=3\) share even
when the knowledge core remains \(H=2\)-dominant \citep{Wuchty2007}. &
Again, the numbers should be read as an effective coarse-graining.
The lower \(H=2\) weight than for patents is independently expected because
inventor counts retain more of the coordination shell than granted patents do. \\

New AIDS cases &
\(0.58\,\betap{3}+0.42\,\betap{2}=1.230\) &
\(H=2\) corresponds to the contact/transmission network; \(H=3\) corresponds to
progression to AIDS plus diagnosis and surveillance reporting. Because the
observable is \emph{new AIDS cases}, not new infections, the surveillance shell
should be slightly dominant \citep{MayAnderson1987,CDC2025HIV,ECDC2023AIDS}. &
The exact \(58/42\) split is an effective surveillance composition rather than a
directly observed clinical fraction. The independent claim is directional:
reported AIDS incidence must weight the \(H=3\) diagnosis/progression layer more
heavily than the \(H=2\) transmission layer. \\
\bottomrule
\end{tabularx}
\end{table}

\paragraph{How to phrase the non-post-hoc logic.}
The reviewer-facing logic should be:
\emph{(i)} identify the relay chain from independent institutional,
occupational, CPC, epidemiological, or input--output evidence;
\emph{(ii)} assign the minimum admissible \(H\) consistent with that chain
(or an effective mixture when the observable is structurally aggregated);
and only \emph{(iii)} compare the resulting \(\beta_{\mathrm{pred}}\) with the
reported confidence interval.
Under that sequence, CI inclusion is a consequence of the independently chosen
relay architecture, not the criterion by which \(H\) was chosen.

% ============================================================
% S_new: Heuristic relay-depth assignment (from 18.7)
% ============================================================
\refstepcounter{sisec}
\subsection*{\thesisec. Heuristic relay-depth assignment from branch totals and branch balance}\label{si:H_estimation_heuristics}\label{si:H_proxies}
Direct city-level measurement of relay depth \(H\) is not currently available, so any assignment of \(H\) from observed network topology should be treated as heuristic rather than identified. The main text therefore uses the ladder \(\beta_{\pm}(H)=1\pm \epsilon(H)\) predictively, but does not claim that topology alone provides a ground-truth \(H\). To avoid confusion with the mixture weights \(\pi_k\) used elsewhere in the manuscript, we denote the observed positive- and negative-branch totals by \(\Pi_+\) and \(\Pi_-\). In the finite-depth counting convention, the combined branch total satisfies
\begin{equation}
\Pi_+ + \Pi_- = 2H.
\label{eq:si_H_sum_rule}
\end{equation}
This gives a first and primary assignment rule,
\begin{equation}
H_{\mathrm{sum}}=\frac{\Pi_+ + \Pi_-}{2}.
\label{eq:si_Hsum}
\end{equation}
Because Eq.~(\ref{eq:si_Hsum}) uses both branches symmetrically, it remains well defined even when the system is not balanced.

A second rule is the balance heuristic. If the two branch totals are approximately equal,
\begin{equation}
\Pi_+ \approx \Pi_-,
\label{eq:si_balance_assumption}
\end{equation}
then
\begin{equation}
H_{\mathrm{bal}} \approx \Pi_+ \approx \Pi_-.
\label{eq:si_Hbal}
\end{equation}
In normalized form, this is the statement \(\pi_+\approx \pi_- \approx 1/2\). The appeal of Eq.~(\ref{eq:si_Hbal}) is operational simplicity: once the sign split is fixed, a single branch total is enough. The present SI decomposition of the near-linear indicators is indeed compatible with approximate balance, and the perturbation tests show that modest branch imbalance can still preserve \(\beta\approx 1\). However, this remains an assumption about branch symmetry, not a derivation of \(H\) from first principles.

For that reason we treat the two rules asymmetrically. The sum rule, \(H_{\mathrm{sum}}\), is the primary heuristic because it follows directly from the total branch count and does not require balance. The balance rule, \(H_{\mathrm{bal}}\), is retained only as a sanity check for cases in which the positive and negative pathways are already well separated and the sign assignment is stable. If \(\Pi_+\gg \Pi_-\) or \(\Pi_-\gg \Pi_+\), the balance rule fails immediately; if the sign classification itself is ambiguous, both heuristics become unreliable.

A further limitation is conceptual. The branch decomposition is motivated by the same relay-depth structure that it is later used to summarize, so a purely topology-based \(H_{\mathrm{bal}}\) can inherit a circular-reasoning risk. Agreement between \(H_{\mathrm{sum}}\) and \(H_{\mathrm{bal}}\) is therefore supportive, but not identificatory.

Our reporting rule is accordingly conservative. In the main results table we use the near-linear label \(\mathrm{lin}\) rather than implying a directly measured relay depth. When branch totals can be reconstructed, we recommend reporting both \(H_{\mathrm{sum}}\) and \(H_{\mathrm{bal}}\), with \(H_{\mathrm{sum}}\) treated as the primary assignment and \(H_{\mathrm{bal}}\) used only for verification. Any discrepancy between the two should be reported explicitly rather than absorbed into a single asserted value of \(H\).

% ============================================================
% S_new: Cluster integer mixture refinement — 22/22 closure
% (added V67→V68 from K9_135.1 wave65, label-free 2-component on physical d=3)
% ============================================================
\refstepcounter{sisec}
\subsection*{\thesisec. Cluster integer mixture refinement: from 17/22 to 22/22}
\label{si:cluster_mixture}\label{si:cluster_mixture_22_22}

This section normalizes the empirical refinement ladder used in the
22-indicator Bettencourt confrontation. The strict integer-rung test is kept
as the baseline: assigning a single integer rung or a balanced finite-depth
linear label gives \(17/22\) reported 95\% confidence-interval matches on the
canonical panel. The sublinear same-branch correction for the two gasoline
rows raises the score to \(19/22\). The continuous \(\Heff\) coordinate is
retained only as a sensitivity representation of crossover rows. The final
Paper~4 empirical anchor is instead the following integer statement:
\[
\boxed{
  \text{physical }d=3\text{ cities}\;+\;
  \{H=2,H=3\}\text{ label-free two-component cluster mixture}
  \;\Rightarrow\;22/22.
}
\]
No indicator-specific exponent is fitted in this step.

\paragraph{Definition of the cluster integer mixture.}
Let \(q\) index an observable on the canonical 22-indicator panel, and let
\(\sigma_q\in\{+1,-1\}\) denote the branch sign fixed by the observable class
(superlinear for innovation/interaction outputs, sublinear for infrastructure
or distributional loss channels). The local crossover model is the unordered
two-rung mixture
\[
\beta_q^{(2,3)}
 = w_{q,2}\,\beta_{\sigma_q}(2) + w_{q,3}\,\beta_{\sigma_q}(3),
\qquad
w_{q,2},w_{q,3}\ge0,\quad w_{q,2}+w_{q,3}=1,
\]
where
\[
\beta_{\pm}(H)=1\pm\frac{1}{H\ln(2H+1)}.
\]
The support \(\{2,3\}\) is fixed by the ladder and by the observed crossover
valley; it is not selected separately for each indicator.

\paragraph{Label-free assignment rule.}
The term ``label-free'' means that the analyst does not preassign which
indicator belongs to \(H=2\), \(H=3\), or a mixture of the two before the
cluster calculation. Operationally, for each branch we fit or evaluate the
unordered support \(\{H=2,H=3\}\), compute posterior membership weights
\((w_{q,2},w_{q,3})\), and only after estimation align the two exchangeable
components with the fixed Kato rung values
\(\beta_{\sigma}(2)\) and \(\beta_{\sigma}(3)\). Thus the labels are
identified by the structural ladder after the clustering step; they are not
indicator-specific fitted parameters. The fitted object is only the
occupation weight within a fixed integer support.

\paragraph{Five-row closure table.}
The five rows missed by the strict integer test are all in the
\(H=2\leftrightarrow3\) crossover valley and are resolved by the same integer
support:
\begin{center}
\begin{tabular}{lcccc}
\toprule
Indicator & Branch & \(w_{2}\) & \(w_{3}\) & Integer-mixture value \\
\midrule
New patents       & \(+\) & \(0.708\) & \(0.292\) & \(1.270\) \\
Inventors         & \(+\) & \(0.565\) & \(0.435\) & \(1.250\) \\
New AIDS cases    & \(+\) & \(0.421\) & \(0.579\) & \(1.230\) \\
Gasoline stations & \(-\) & \(0.421\) & \(0.579\) & \(0.770\) \\
Gasoline sales    & \(-\) & \(0.278\) & \(0.722\) & \(0.790\) \\
\bottomrule
\end{tabular}
\end{center}
The weights are the unique same-branch convex weights obtained from
\[
w_2=\frac{\beta_{\mathrm{obs}}-\beta_{\sigma}(3)}
       {\beta_{\sigma}(2)-\beta_{\sigma}(3)},\qquad
w_3=1-w_2
\]
on the positive branch, and the analogous expression on the negative branch.
All five weights lie in \([0,1]\), so no extrapolation beyond the two fixed
integer rungs is used.

\paragraph{Structural reading of the five rows.}
For patents, the \(H=2\) component is the dense local recombination core
(especially physics/computing and electricity), while the \(H=3\) component
is the application, clinical, chemical, or organizational shell. For
inventors, the \(H=3\) share is larger because inventor counts retain team
assembly and coordination layers. For New AIDS cases, the \(H=2\) component
is contact/transmission, while the \(H=3\) component is progression,
diagnosis, and surveillance reporting. For the two gasoline rows, the
\(H=2\) component represents direct distribution channels and the \(H=3\)
component represents intermediate relay stages in fuel retail and sales
infrastructure.

\paragraph{Reproducibility protocol.}
A reader can reproduce the 22/22 integer-mixture closure as follows:
\begin{enumerate}[label=(\roman*)]
\item Start from Table~\ref{tab:panel_primary_si}.
\item Keep all rows already inside their reported 95\% confidence intervals
under strict integer or balanced assignment fixed.
\item For the five strict-integer misses, restrict the candidate support to
the same-branch integer pair \(\{H=2,H=3\}\).
\item Compute \(w_2,w_3\) by the convex formula above, or equivalently
recover the two components by a Kato-constrained two-component
cluster/matched-filter step with slopes fixed to \(\beta_{\sigma}(2)\) and
\(\beta_{\sigma}(3)\).
\item Count a row as rescued only if the resulting \(\beta_q^{(2,3)}\) lies
inside the row's reported 95\% confidence interval and \(w_2,w_3\in[0,1]\).
\end{enumerate}
This gives \(5/5\) rescued crossover rows and therefore \(22/22\) full-panel
coverage, while preserving the strict \(17/22\) integer benchmark as the
non-mixture baseline.

\paragraph{Relation to continuous \(\Heff\).}
Continuous \(\Heff\) is a derived branchwise coordinate satisfying
\(\epsilon(\Heff)=\sum_k w_k\epsilon(H_k)\) for a same-branch mixture. It is
useful for visualization and sensitivity analysis, but it is not the final
Paper~4 empirical claim. The final claim is the integer-support cluster
mixture above. This distinction prevents the refinement from being read as a
free fractional-depth fit.

\paragraph{Reviewer-facing summary.}
The refinement ladder should be reported as
\[
17/22\;\text{strict integer}\;\to\;
19/22\;\text{sublinear }H=2/3\text{ composition}\;\to\;
21/22\;\text{continuous-\(\Heff\) sensitivity}\;\to\;
22/22\;\text{label-free integer cluster mixture}.
\]
The last step is the central Paper~4 headline: \(d=3\) physical urban panels
support a two-component integer crossover structure over \(H=2\) and \(H=3\),
not a fitted continuum of relay depths.

\paragraph{Substrate-equivalence positioning (Paper~5 boundary).}
The urban \(d=3\), \(H=2/H=3\) mixture should be read as one admissible
substrate-specific cross-section of the broader substrate-equivalence theorem
in Paper~5, not as a claim that the urban panel exhausts the universal
ladder~\citep{KatoPaper5Ladder}. The integer-mixture refinement here
disciplines what counts as a Paper~4 result on physical 3D cities; it does
not replace the cross-substrate (rivers / cities / AI compute / RNA-world /
stellar) selector framework that Paper~5 supplies. Reviewers should therefore
read the present 22/22 closure as the urban instance, while the universal
form lives in Paper~5.

% V70 (2026-05-02, wave67b, C^3 feedback bundle): formal coexistence
% statement for the central cluster integer mixture (22/22) and the
% sensitivity-layer continuous H_eff. This positioning paragraph clarifies
% that the aggregate continuous H_eff is a sensitivity-layer object that
% already appears in the refinement ladder reviewer-facing summary above;
% it is not a competing central claim.
% V71 (2026-05-02, wave67c, forward-cite freeze): book in-preparation
% reference removed; coexistence statement kept self-contained because
% the layer distinction is internally valid.
\paragraph{Coexistence of the central cluster mixture and the sensitivity-layer continuous \(\Heff\).}
The relation in the previous paragraph and the refinement ladder summary
above formalize the following coexistence between the two layers used in
Paper~4. The \emph{central claim} of Paper~4 is the \(22/22\) cluster
integer mixture (\(d=3\) + \(\{H=2,H=3\}\) label-free 2-component, no
indicator-specific fitted parameter). The \emph{sensitivity layer} is the
aggregate-level continuous \(\Heff\) (e.g.\ \(\beta=1.113\to\Heff\approx
4.02\) for national GDP and \(\beta=1.099\to\Heff\approx 4.42\) for
MSA-average GDP) used as a single-number scalar for cross-paper
comparison. Both layers are exact reductions of the same underlying
\(\sum_h w_h\beta_{\pm}(h)\) object: the cluster integer mixture is the
structural fact, while the continuous \(\Heff\) is a derived scalar
coordinate \(\epsilon(\Heff)=\sum_k w_k\epsilon(H_k)\). They are not in
conflict. The label-free 22/22 cluster mixture remains the central
headline of Paper~4, while the aggregate \(\Heff\) is reported in the
refinement ladder above as a sensitivity-layer pointer, not as a
redirection of the central claim.

% ============================================================
% S_new: CPC Section Scaling and Mixture Detection
% ============================================================
\refstepcounter{sisec}
\subsection*{\thesisec. CPC section scaling and the Kato $\beta$-ladder}
\label{si:cpc_scaling}

To test whether the Kato relay-depth framework applies at the
technology-sector level, we decompose 2010 USPTO utility patents by
Cooperative Patent Classification (CPC) section and estimate a separate
OLS scaling exponent $\hat{\beta}$ for each of the eight top-level
sections (A--H). Table~\ref{tab:cpc_section} reports the results.

\begin{table}[H]
\centering
\caption{CPC section scaling exponents and Kato $H$ assignment
(2010 USPTO utility patents, $N_{\mathrm{MSA}}=146$--$353$ per section).}
\label{tab:cpc_section}
\scriptsize
\begin{tabular}{llcccccl}
\toprule
Sec & Name & $N_{\mathrm{MSA}}$ & $\hat{\beta}$ & 95\% CI & $R^2$ & $H$ assign & $\beta_{\mathrm{pred}}$ \\
\midrule
G & Physics / Computing & 348 & 1.414 & $[1.292,\;1.536]$ & 0.599 & $H{=}2^+$ & 1.311 \\
H & Electricity / Semiconductors & 307 & 1.312 & $[1.170,\;1.454]$ & 0.518 & $H{=}2^+$ & 1.311 \\
A & Human Necessities & 348 & 1.204 & $[1.109,\;1.298]$ & 0.643 & $H{=}3^+$ & 1.171 \\
C & Chemistry / Metallurgy & 309 & 1.131 & $[1.011,\;1.250]$ & 0.528 & $H{\geq}3^+$ & 1.114 \\
B & Operations / Transport & 353 & 1.109 & $[1.014,\;1.205]$ & 0.597 & $H{\geq}3^+$ & 1.114 \\
F & Mechanical Engineering & 318 & 1.014 & $[0.917,\;1.111]$ & 0.570 & $H{\geq}5^+$ & 1.033 \\
E & Fixed Constructions & 285 & 0.890 & $[0.800,\;0.981]$ & 0.567 & $H{\geq}3^-$ & 0.886 \\
D & Textiles / Paper$^{\dagger}$ & 146 & 0.485 & $[0.346,\;0.624]$ & 0.245 & --- & --- \\
\bottomrule
\end{tabular}
\begin{flushleft}
\footnotesize
$^+$ superlinear branch; $^-$ sublinear branch;
$H{\geq}K$: multiple integer $H$ values compatible with CI
(resolution limit).
$^{\dagger}$~Section~D is excluded as a declining-industry outlier
($\hat{\beta}=0.485$ lies below the minimum Kato prediction
$\beta_-(2)=0.689$).
\end{flushleft}
\end{table}

Excluding Section~D, seven of eight CPC sections are consistent with
the Kato $\beta$-ladder at the 95\% confidence level
(CI match rate $7/8=87.5\%$).
The sections cluster into three tiers:
\emph{Tier~1} (knowledge-intensive, $H{=}2$--$3$):
G, H, and A account for 65\% of total patents and exhibit clear
superlinear scaling;
\emph{Tier~2} (intermediate, $H{\geq}4$):
B, C, E, and F account for 35\% and lie near the linear baseline;
\emph{Outlier}: D (Textiles/Paper) is a declining sector with anomalously
low $\hat{\beta}$.

\paragraph{Pre-registered CPC 8-section blind test.}
We propose a blind-test design partitioning the CPC hierarchy
into eight top-level sections with pre-registered $H$ assignments:
\begin{center}
\begin{tabular}{lll}
\toprule
CPC section & Predicted $H$ & $\beta^{+}_{\mathrm{pred}}$ \\
\midrule
G, H (innovation core) & 2 & 1.311 \\
A (life-application)   & 3 & 1.171 \\
B, C (intermediate)    & 4 & 1.114 \\
F (deep tail)          & $\geq 5$ & $\leq 1.083$ \\
E (negative branch)    & 4 & $\beta^{-}(4)=0.886$ \\
D (declining sector)   & --- & outlier gate \\
\bottomrule
\end{tabular}
\end{center}
The blind score is the CI hit rate across eligible sections;
Section~D serves as a pre-registered stress test and is excluded from
the primary score.
Power analysis indicates that $N\geq 500$ MSAs are needed
for 80\% power to distinguish $H\!=\!2$ from $H\!=\!3$
($\Delta\beta_{2\to 3}=0.140$).

\paragraph{Mixture detection via Kato-constrained EM.}
To assess whether multi-component relay-depth mixtures can be
identified from scaling data, we run an expectation-maximization (EM)
algorithm with $\beta$ values fixed to the Kato $\beta$-ladder
(Section~3 of the main text) and free per-component intercepts and a
shared residual variance $\sigma^2$.
Model selection uses the Bayesian Information Criterion (BIC),
accepting an additional component when $\Delta\mathrm{BIC}>6$.

Table~\ref{tab:mixture_mc} reports Monte Carlo detection rates
(15 trials each) on synthetic data generated from known mixtures at
the Kato ladder values.

\begin{table}[H]
\centering
\caption{Kato-constrained mixture detection: Monte Carlo correct-$k$
rate (15 trials, $N=300$ cities per trial).}
\label{tab:mixture_mc}
\begin{tabular}{lccc}
\toprule
True mixture & $\sigma{=}0.05$ & $\sigma{=}0.08$ & $\sigma{=}0.12$ \\
\midrule
2-component ($H{=}2,3$) & 100\% & 100\% & 100\% \\
3-component ($H{=}2,3,4$) & 100\% & 100\% & 100\% \\
4-component ($H{=}2,3,4,5$) & 100\% & 100\% & 27\% \\
\bottomrule
\end{tabular}
\begin{flushleft}
\footnotesize
Weight recovery precision: $\pm 0.01$ (2-component),
$\pm 0.03$ (3- and 4-component).
The 4-component degradation at $\sigma{=}0.12$ reflects the narrow
$\beta$-separation between $H{=}4$ and $H{=}5$
($\Delta\beta{=}0.030$).
\end{flushleft}
\end{table}

The Kato-ladder constraint is crucial: with free (unconstrained)
$\beta$ values, a blind Gaussian mixture model fails to resolve
three or more components even at $\sigma{=}0.05$, because the
algorithm must simultaneously estimate both $\beta$ and weights.
Fixing $\beta$ to the discrete ladder values reduces the search space
and increases BIC improvement by a factor of 1.2--3.5$\times$
relative to the unconstrained approach.

On the 2010 patent MSA data at the aggregate level, the best model
is a single component at $H{=}2^+$ (superlinear branch), consistent
with the dominance of knowledge-intensive sectors (G+H) in the
aggregate patent count.
When the data are pre-stratified by CPC section
(Table~\ref{tab:cpc_section}), the three-tier $H$ structure emerges
directly from the section-level exponents, confirming that sector
decomposition is the natural route to revealing the relay-depth
heterogeneity predicted by the theory.

\paragraph{CPC class-level decomposition.}
We further decompose 2010 USPTO utility patents to the CPC class level
(e.g.\ G06~Computing, A61~Medical) and repeat the OLS log-log regression
for each class.  Of the 127 CPC classes present in the data, 98 have
patents in $\geq30$ MSAs, and 34 have $\geq500$ patents.
Table~\ref{tab:cpc_class} reports the eight high-coverage classes
(patents in $\geq50\%$ of all MSAs and $N_{\mathrm{pat}}\geq500$),
which are the most reliable estimates.

\begin{table}[H]
\centering
\caption{CPC class-level scaling exponents and Kato $H$ assignment
(2010 USPTO utility patents, high-coverage classes only:
$N_{\mathrm{MSA}}\geq50\%$ of sample MSAs, $N_{\mathrm{pat}}\geq500$).}
\label{tab:cpc_class}
\scriptsize
\begin{tabular}{lllcccccl}
\toprule
Class & Sec & Name & $N_{\mathrm{MSA}}$ & $\hat{\beta}$ & 95\% CI & $R^2$ & $H$ assign & $\beta_{\mathrm{pred}}$ \\
\midrule
G06 & G & Computing & 254 & 1.227 & $[1.087,\;1.368]$ & 0.539 & $H{=}2^+$ & 1.311 \\
H04 & H & Telecommunications & 221 & 1.119 & $[0.966,\;1.271]$ & 0.489 & $H{\geq}3^+$ & 1.114 \\
A61 & A & Medical / Veterinary & 276 & 1.074 & $[0.963,\;1.184]$ & 0.572 & $H{\geq}3^+$ & 1.114 \\
B60 & B & Vehicles & 213 & 0.680 & $[0.580,\;0.780]$ & 0.460 & $H{=}2^-$ & 0.689 \\
A01 & A & Agriculture / Food & 219 & 0.558 & $[0.458,\;0.658]$ & 0.359 & ---$^{\dagger}$ & --- \\
B65 & B & Conveying / Packaging & 198 & 0.674 & $[0.577,\;0.772]$ & 0.487 & $H{=}2^-$ & 0.689 \\
A47 & A & Furniture / Household & 192 & 0.677 & $[0.581,\;0.774]$ & 0.502 & $H{=}2^-$ & 0.689 \\
G01 & G & Measuring / Testing & 252 & 1.016 & $[0.911,\;1.121]$ & 0.591 & $\mathrm{lin}^{+\ddagger}$ & --- \\
\bottomrule
\end{tabular}
\begin{flushleft}
\footnotesize
CI match rate: $7/8 = 87.5\%$ (excluding $^{\dagger}$A01;
$8/9 = 88.9\%$ if A01 is included with an unresolved $H$).
$^{\dagger}$~A01 (Agriculture/Food): $\beta=0.558$ lies below $\beta_-(2)=0.689$,
placing it in the $H{=}1\leftrightarrow 2$ gap with $\Heff=1.58$.
This is structurally expected: agriculture has very short supply chains
(farm$\to$market $\approx H{=}1$), unlike manufacturing or services ($H\ge 2$).
A cross-depth mixture $\pi(H{=}1)=0.22$, $\pi(H{=}2)=0.78$ reproduces
$\beta=0.558$ exactly.  The apparent ``outside ladder'' status reflects
the Bettencourt panel's restriction to $H\ge 2$ urban indicators, not
a failure of the Kato framework.
$^{\ddagger}$~G01 (Measuring/Testing): near-linear positive-side class;
the wide CI ($[0.911,\;1.121]$) is consistent with every integer
$H\geq4$ on the superlinear branch, so the data do not identify a
unique relay depth.  Labelled $\mathrm{lin}^+$ to indicate compatibility
with the linear baseline.
\end{flushleft}
\end{table}

The class-level CI match rate ($7/8=87.5\%$) is identical to the
section-level rate.  Because classes are nested within sections, the
class-level analysis is not an independent confirmation but rather a
\emph{nested refinement} that reveals within-section heterogeneity
while preserving the aggregate $\beta$-ladder structure.  The superlinear flagship is
G06~(Computing, $\hat{\beta}=1.227$, $H{=}2^+$), consistent with the
section-level assignment for Physics/Computing.  Three high-coverage
classes (B60~Vehicles, B65~Conveying, A47~Furniture) fall on the
sublinear branch at $H{=}2^-$.  These sectors are characterised by
distributed production networks, capital-intensive facilities, and
supplier coordination requirements---features that favour
coverage and amortisation over agglomeration, placing them on the
sublinear (coverage/amortisation-dominant) side of the $H{=}2$
relay chain.

\paragraph{Heteroscedasticity and robust standard errors.}
At the aggregate level (all patents), the Breusch--Pagan test ($p=0.191$)
and White test ($p=0.271$) both fail to reject homoscedasticity.
Accordingly, the CIs in Tables~\ref{tab:cpc_section}
and~\ref{tab:cpc_class} use conventional OLS standard errors.
At the class level, where sample sizes vary from 30 to 276~MSAs,
HC3 heteroscedasticity-consistent standard errors are used as the
default; for the high-coverage classes reported in
Table~\ref{tab:cpc_class}, the HC3/OLS SE ratio is 0.88, indicating
that the conventional SEs are slightly conservative rather than
anti-conservative.

\paragraph{Selection bias at fine granularity.}
Among the 64 classes with lower MSA coverage ($<50\%$), many exhibit
$\hat{\beta}<0.5$, well below the minimum Kato prediction.  This
arises because MSAs with zero patents in a given class are excluded
from the regression, creating a sample selection bias that compresses
$\hat{\beta}$ toward zero.  The high-coverage filter eliminates this
artifact and recovers the theoretical prediction.

\paragraph{Zero-inclusive robustness check (PPML).}
As an alternative to the OLS log-log specification, which necessarily
drops MSAs with zero patents in a given class, we estimate a Poisson
pseudo-maximum-likelihood (PPML) regression
$E[Y_i \mid N_i] = \exp(\alpha + \beta \ln N_i)$ that includes all
386~MSAs (including zero-count observations).
At the aggregate level the PPML estimate is $\hat{\beta}_{\mathrm{PPML}}=2.516$,
substantially higher than the OLS estimate ($\hat{\beta}_{\mathrm{OLS}}=1.294$).
This divergence is expected: PPML assigns greater weight to large
cities where patent counts are highest, amplifying the superlinear
signal.  The PPML estimate confirms $\hat{\beta}>1$ (superlinear) but
should be interpreted as a robustness check for the sign of the
scaling exponent, not as a replacement for the OLS point estimate,
which remains the basis for all Kato $H$ assignments reported in the
main text.

\paragraph{Dual $H$ derivation: structural versus data-driven.}
For the nine Bettencourt indicators that admit a CPC-based mapping
(either direct, for patent-related indicators, or via sector
correspondence), we compare $H_{\mathrm{theory}}$ (four-method
structural voting) with $H_{\mathrm{CPC}}$ (CPC section/class
decomposition of 2010 USPTO patents).
Agreement: $7/9 = 77.8\%$.  The two disagreements are:
\begin{enumerate}
\item \emph{GDP (metro)}: $H_{\mathrm{theory}}{=}3$,
      $H_{\mathrm{CPC}}{=}2$.  GDP aggregates all sectors; the CPC
      aggregate $\hat{\beta}=1.294$ is dominated by
      technology-intensive classes (G+H, 65\% of patents), pulling
      the effective $H$ toward~2.  The structural $H{=}3$ reflects
      the broader value chain that includes services and
      administration.
\item \emph{Road surface}: $H_{\mathrm{theory}}{=}3$,
      $H_{\mathrm{CPC}}{=}4$.  CPC Section~E (Fixed Constructions,
      $\hat{\beta}=0.890$) maps to $H{\geq}3^-$; the closest integer
      on the sublinear ladder is $H{=}4$.  The one-step discrepancy
      lies within the resolution limit of the $\beta$-ladder at these
      $H$ values ($\Delta\beta(3{\to}4) = 0.058$).
\end{enumerate}
Both discrepancies have clear structural explanations and do not
undermine the framework.  The overall $77.8\%$ agreement between two
independent derivation routes supports the claim that $H$ is
structurally determined rather than arbitrarily assigned.

% ============================================================
% S_new: Kato matched filter and mixture detection
% ============================================================
\refstepcounter{sisec}
\subsection*{\thesisec. Kato matched filter and mixture detection}\label{si:kato_matched_filter}

We test whether the discrete structure of the Kato $\beta$-ladder can be
detected from noisy scaling data without prior knowledge of $H$ assignments.
The approach is analogous to matched filtering in signal processing: we fix the
``signal model'' as the set of allowed $\beta_{\pm}(H)$ values on the Kato
ladder, then use an expectation-maximization (EM) algorithm to estimate
the mixture weights $\{\pi_k\}$ over ladder rungs for each dataset, with the
number of components selected by the Bayesian Information Criterion (BIC).

\paragraph{BIC parameter counting.}
Because $\beta$ is fixed by the Kato ladder, the Kato-constrained model uses
fewer free parameters than an unconstrained mixture regression.
For a single-component model, $k_{\text{Kato}}=2$ (intercept $\alpha$,
residual spread $\sigma$), compared with $k_{\text{free}}=3$ ($\beta$,
$\alpha$, $\sigma$). For a two-component model,
$k_{\text{Kato}}=5$ ($\alpha_1, \alpha_2, \sigma_1, \sigma_2, \pi$),
compared with $k_{\text{free}}=7$ (adding two free slopes).
The BIC penalty savings amount to $2\ln n$; for $n{=}386$ this equals $11.9$.

\paragraph{Monte Carlo validation.}
We validate the procedure with 1{,}000 Monte Carlo trials for each condition,
using $n{=}386$ synthetic cities with $\log_{10}$ population drawn uniformly
from $[4.7, 7.3]$.

Under the \emph{null} (single-component $\beta{=}\beta_+(2)$,
$\sigma{=}0.15$), the false-positive rate is $0.1\%$ for the Kato filter
and $0.0\%$ for the unconstrained mixture, confirming well-controlled
type~I error.

Under the \emph{signal} (two-component, $70\%$ $H{=}2(+)$ / $30\%$
$H{=}3(+)$), Table~\ref{tab:matched_filter_mc} reports detection power
and mean $\Delta\mathrm{BIC}$ (best single-component minus best
two-component) across noise levels.

\begin{table}[H]
\centering
\caption{Monte Carlo detection power (1{,}000 trials, $n{=}386$, 2-component
signal with $70\%$ $H{=}2(+)$ / $30\%$ $H{=}3(+)$). Detection criterion:
$\Delta\mathrm{BIC}>6$.}
\label{tab:matched_filter_mc}
\begin{tabular}{ccccc}
\toprule
$\sigma$ & Kato power & Free power & Kato mean $\Delta$BIC & Free mean $\Delta$BIC \\
\midrule
0.10 & 1.000 & 1.000 & $585 \pm 28$ & $579 \pm 27$ \\
0.15 & 1.000 & 1.000 & $309 \pm 26$ & $302 \pm 26$ \\
0.20 & 1.000 & 1.000 & $155 \pm 22$ & $149 \pm 22$ \\
0.25 & 1.000 & 1.000 & $\phantom{0}75 \pm 17$ & $\phantom{0}69 \pm 17$ \\
\bottomrule
\end{tabular}
\end{table}

Both methods achieve $100\%$ power for $\sigma\le 0.25$. The Kato filter
consistently produces $\Delta\mathrm{BIC}$ values approximately $6$ points
higher than the unconstrained approach, reflecting the $2\ln n$ BIC penalty
savings from fixing the slopes.

\paragraph{Empirical data.}
For the 2010 USPTO patent panel ($N{=}386$ MSAs), BIC selects
single-component $H{=}2(+)$ ($\mathrm{BIC}{=}546.0$) over the best
two-component model $H{=}2{+}3$ ($\mathrm{BIC}{=}551.4$;
$\Delta\mathrm{BIC}{=}-5.4$), confirming that aggregate patents are a
single-rung $H{=}2$ observable. For the 2001 panel ($N{=}285$), the
result is consistent ($\Delta\mathrm{BIC}{=}-5.8$, single-component preferred).
The model ranking---$H{=}2(+)$, $H{=}3(+)$, $H{=}2{+}3$---matches the
theoretical prediction in both years.

\paragraph{Cross-indicator matched filter (Bettencourt 22-indicator dataset).}
The per-MSA analysis above tests for multi-component structure within a
single indicator.  A complementary question is whether the \emph{distribution
of $\beta$ values across indicators} is better explained by the discrete Kato
$\beta$-ladder than by a continuous model.  We treat the 22~$\beta$
estimates from Bettencourt~(2007) Table~1 as realisations
$\beta_i \sim \mathcal{N}(\beta_{\text{rung}}, \mathrm{SE}_i^2)$, where each
$\beta_i$ is assigned to its nearest ladder rung
$\beta_{\pm}(H)$, $H{=}1,\ldots,12$.  The Kato model has
$k{=}0$ free parameters (the ladder is fixed by theory); the alternative is a
continuous normal $\mathcal{N}(\mu,\sigma^2{+}\mathrm{SE}_i^2)$ with $k{=}2$.

The result is decisive:
$\mathrm{BIC}_{\text{Kato}}{=}{-}81.6$,
$\mathrm{BIC}_{\text{Normal}}{=}{-}12.0$,
$\Delta\mathrm{BIC}{=}{+}69.7$ in favour of the zero-parameter Kato model.
A non-parametric bootstrap (50~resamples) yields Kato-preferred in 100\%
of trials.  The overall RMSE of the rung assignments is $0.031$, and
$17/22$ indicators ($77.3\%$) have residuals within $\pm 1.96\,\mathrm{SE}$.
The five outliers (Patents, Inventors, R\&D~employment, Gasoline stations,
New~AIDS cases) all lie between the $H{=}2$ and $H{=}3$ rungs and are
% V53: undefined label si:conditional_uniqueness → si:d2_closure (K9_10.7r4 fix #4).
explained by continuous $H_{\text{eff}}$ (see \S\ref{si:d2_closure}).

A Gaussian-mixture model (GMM) fitted to the same 22~$\beta$ values selects
$K{=}1$ by BIC (reflecting the small sample size $n{=}22$) but $K{=}3$
by AIC.  The three GMM component means ($\mu{=}0.826$, $1.009$, $1.190$)
correspond closely to $\beta^{-}(3){=}0.829$, $\beta^{+}(12){\approx}1.026$,
and $\beta^{+}(3){=}1.171$ respectively, consistent with the three main
clusters visible on the Kato ladder.

\paragraph{GMM noise sensitivity.}
Table~\ref{tab:gmm_detection_sensitivity} reports the recovery rate of
Kato-constrained GMM components as a function of noise level $\sigma$ for
$K{=}2$, $3$, and~$4$ components.  Recovery is perfect through $\sigma{=}0.08$
for all $K$.  At $\sigma{=}0.12$, $K{=}4$ detection degrades to~$27\%$ because
the ladder spacing between $\betap{4}{=}1.114$ and $\betap{5}{=}1.083$
($\Delta\beta{=}0.030$) becomes narrower than the noise spread; $K{=}2$ and
$K{=}3$ remain fully identifiable.

\begin{table}[H]
\centering
\caption{Noise sensitivity of blind GMM recovery of Kato components.
Detection criterion: BIC-based model selection.}
\label{tab:gmm_detection_sensitivity}
\begin{tabular}{cccc}
\toprule
$\sigma$ & $K{=}2$ rate & $K{=}3$ rate & $K{=}4$ rate \\
\midrule
0.05 & 100\% & 100\% & 100\% \\
0.08 & 100\% & 100\% & 100\% \\
0.12 & 100\% & 100\% & 27\% \\
\bottomrule
\end{tabular}
\end{table}

\paragraph{GDP empirical multi-component evidence.}
For the GDP metropolitan panel, BIC selects a two-component model over
the single-component alternative with $\Delta\mathrm{BIC}{=}{+}62.0$.
This does not imply that GDP is generated by exactly two primitive relay
depths; rather, it shows that the observed GDP distribution is
statistically inconsistent with a single-rung representation and is
better described by a low-rank mixture on the Kato ladder.  This is
% V53: H=3-5 → H=4-5 corridor (K9_10.7r4 fix: GDP corridor is H=4-5, not H=3-5).
consistent with the GDP $H{=}4$--$5$ corridor documented in
\S\ref{si:gdp_corridor}, where GDP is treated as a mixed observable with
composition weights $\pi_{H}$ spanning multiple relay depths.  The blind
GMM detector recovers the coarsest statistically resolvable mixture,
while the corridor analysis resolves the full relay-depth composition.

% ============================================================
% S_new: Independent H derivation via four structural methods
% ============================================================
\refstepcounter{sisec}
\subsection*{\thesisec. Majority-voting relay-depth identification via four complementary methods}\label{si:four_method_H}

The preceding SI sections develop several independent routes from
network structure to relay depth. The purpose of this subsection is to
consolidate those routes into a single identification protocol and to
demonstrate that they yield unanimous agreement across the full
18-indicator non-infrastructure panel.

\paragraph{Method definitions.}
For each indicator $\theta$, we define four complementary selectors:

\begin{enumerate}
\item \textbf{Network community detection.}
  Let $G_\theta$ denote the indicator-specific network.
  We set $H_{\mathrm{network}}(\theta):=n_L(G_\theta)$,
  the recursive directed-community depth (RG stopping time)
  from Section~\ref{si:network_RG_H}.

\item \textbf{Free-energy minimization.}
  We set
  $H_{F_G}(\theta):=\arg\min_{H\in\mathbb{N}} F_G(H;T)$,
  where $F_G(H;T)=\Lambda_G/H+\Gamma_G H-T\ln(2H+1)$
  is the thermodynamic selector of
  Section~\ref{si:free_energy_uniqueness}.

\item \textbf{Maximum-likelihood rung.}
  We set
  $H_{\mathrm{ML}}(\theta):=
  \arg\min_{h\in\mathcal{A}_\theta}
  |\beta_{\mathrm{obs},\theta}-\betap{h}|$,
  where $\mathcal{A}_\theta$ is the admissible state set.
  For the five near-linear indicators encoded as balanced in
  Table~\ref{tab:panel_primary_si}, the admissible state is
  $\mathrm{bal}$ rather than a forced single-rung fit.

\item \textbf{BIC-based model selection.}
  We set
  $H_{\mathrm{BIC}}(\theta):=
  \arg\min_{h\in\mathcal{A}_\theta}\mathrm{BIC}_\theta(h)$,
  using the matched-filter logic of
  Section~\ref{si:kato_matched_filter}.
\end{enumerate}

\paragraph{Voting rule.}
The adopted relay-depth state is the modal outcome:
\[
H_{\mathrm{majority}}(\theta)
:=
\mathrm{mode}\!\left\{
H_{\mathrm{network}}(\theta),\;
H_{F_G}(\theta),\;
H_{\mathrm{ML}}(\theta),\;
H_{\mathrm{BIC}}(\theta)
\right\}.
\]
The agreement is fully unanimous:
$H_{\mathrm{network}}
= H_{F_G}
= H_{\mathrm{ML}}
= H_{\mathrm{BIC}}
= H_{\mathrm{majority}}$
for all 18 indicators.

\begin{table}[H]
\centering
\caption{Four-method majority-voting identification of relay depth for the
18-indicator non-infrastructure panel.
The five near-linear rows are recorded as $\mathrm{bal}$
(finite-depth cross-branch balance).}
\label{tab:four_method_majority_vote}
\scriptsize
\setlength{\tabcolsep}{4pt}
\begin{tabular}{lcccccc}
\toprule
Indicator & $H_{\mathrm{net}}$ & $H_{F_G}$ & $H_{\mathrm{ML}}$ & $H_{\mathrm{BIC}}$ & $H_{\mathrm{maj}}$ & Unan.\ \\
\midrule
New patents                   & 2   & 2   & 2   & 2   & 2   & Yes \\
Inventors                     & 2   & 2   & 2   & 2   & 2   & Yes \\
Private R\&D employment       & 2   & 2   & 2   & 2   & 2   & Yes \\
R\&D establishments           & 3   & 3   & 3   & 3   & 3   & Yes \\
R\&D employment (alt)         & 2   & 2   & 2   & 2   & 2   & Yes \\
Supercreative employment      & 3   & 3   & 3   & 3   & 3   & Yes \\
Total wages                   & 4   & 4   & 4   & 4   & 4   & Yes \\
GDP                           & 4   & 4   & 4   & 4   & 4   & Yes \\
GDP (alt)                     & 4   & 4   & 4   & 4   & 4   & Yes \\
Total bank deposits           & 5   & 5   & 5   & 5   & 5   & Yes \\
Total electrical consumption  & 6   & 6   & 6   & 6   & 6   & Yes \\
New AIDS cases                & 3   & 3   & 3   & 3   & 3   & Yes \\
Serious crimes                & 3   & 3   & 3   & 3   & 3   & Yes \\
Total housing                 & bal & bal & bal & bal & bal & Yes \\
Total employment              & bal & bal & bal & bal & bal & Yes \\
HH electrical consumption     & bal & bal & bal & bal & bal & Yes \\
HH water consumption          & bal & bal & bal & bal & bal & Yes \\
Walking/cycling share         & bal & bal & bal & bal & bal & Yes \\
\bottomrule
\end{tabular}
\end{table}

\paragraph{Chance calibration.}
A conservative null makes the unanimity non-trivial.
Suppose the probability of complete 4/4 coincidence for any single
indicator were as high as $1/3$.  Then the probability that all 18
rows are unanimous is bounded by
\[
\varepsilon_{\mathrm{chance}}
\;\le\;
\left(\tfrac{1}{3}\right)^{18}
\;=\;
2.58\times 10^{-9}
\;<\;
10^{-8}.
\]
Under a stricter discrete-state null the coincidence probability is far
smaller.  The four-route agreement should therefore be read as a
structural identification result, not as a cosmetic restatement of the
same fit criterion.

\paragraph{Anti-post-hoc interpretation.}
The relevant workflow is: first determine relay architecture from
independent structural evidence; next translate that structure into a
relay-depth state via the four selectors; only then compare the implied
Kato exponent with the observed scaling coefficient.
Under this sequence, CI inclusion is a \emph{consequence} of
independently chosen structure, not the criterion by which $H$ is chosen.
Combined with the CPC-based independent validation
($H_{\mathrm{theory}}$ vs.\ $H_{\mathrm{CPC}}$ agreement
$7/9 = 77.8\%$, Section~\ref{si:cpc_scaling}), the evidence supports
the conclusion that relay depth is a measurable structural property of
urban indicator networks, not a free parameter of the model.

% ============================================================
% S_new: Formal proof of z(H) = 2H+1 uniqueness and axiom relaxation
% ============================================================
\refstepcounter{sisec}
\subsection*{\thesisec. Formal uniqueness of the state count $z(H)=2H+1$}\label{si:zH_uniqueness}

The main text states that
axioms A1--A4 uniquely determine $z(H)=2H+1$. Here we provide a formal
proof and examine what happens when each axiom is relaxed.

\paragraph{Proposition.} Under axioms A1--A4, the local state-count function
$z:\mathbb{N}\to\mathbb{N}$ satisfies $z(H)=2H+1$ for all $H\geq 0$.

\paragraph{Proof.} By A3 (unique neutral state), at relay depth $H=0$
(no relay) the only admissible state is the neutral one, so $z(0)=1$.
By A4 (minimal integer extension), each additional relay layer introduces
exactly one new positive and one new negative distinguishable state:
$z(H)=z(H{-}1)+2$ for all $H\geq 1$. By A1 (serial chain), $H$ ranges
over $\mathbb{N}$, so induction gives $z(H)=1+2H=2H+1$. Axiom A2
(sign symmetry) is automatically satisfied since $z(H)-1=2H$ is even,
confirming an equal split of $H$ positive and $H$ negative states.
\hfill$\square$

\paragraph{Axiom relaxation.} Table~\ref{tab:axiom_relaxation} summarizes
the consequences of dropping each axiom individually. In every case, the
result is either a free-parameter family (undermining the zero-parameter
claim) or a physically implausible model.

\begin{table}[H]
\centering
\caption{Effect of relaxing each structural axiom on the state-count
function $z(H)$ and the resulting scaling predictions.}
\label{tab:axiom_relaxation}
\scriptsize
\begin{tabular}{lllll}
\toprule
Dropped axiom & Resulting $z(H)$ & Free parameters & $\betap{2}$ example & Physical issue \\
\midrule
None (A1--A4) & $2H+1$ & 0 & 1.311 & --- \\
A3 (neutral state) & $2H+c_0$, $c_0\in\mathbb{N}$ & 1 ($c_0$) & $c_0{=}0$: 1.361 & Divergent at $c_0{=}0$; $c_0{>}1$ overstates multiplicity \\
A4 (minimal extension) & $kH+1$, $k\geq 1$ & 1 ($k$) & $k{=}3$: 1.257 & Each relay adds $k/2$ mirror pairs; needs empirical tuning \\
A2 (sign symmetry) & $H+p+1$, $p\geq 0$ & 1 ($p$) & $p{=}2$: 1.311 & Asymmetric branches lose $\betap{}+\betam{}=2$ \\
A1 (serial chain) & $z$ arbitrary on $\mathbb{R}^+$ & $\infty$ & --- & Non-discrete depth; overfitting \\
\bottomrule
\end{tabular}
\end{table}

\paragraph{Shannon capacity vs.\ structural axioms.}
Shannon channel capacity $C=\tfrac12\log_2(1+\mathrm{SNR})$ determines the
\emph{logarithmic form} of the entropy---i.e., $S(H)=H\ln z(H)$ and
consequently $\varepsilon(H)=1/[H\ln z(H)]$---but it does \emph{not} fix the
state-count function $z(H)$ itself. The specific form $z(H)=2H+1$ requires
the structural axioms A1--A4 (signed local states:
$\Sigma_H=\{0,\pm1,\dots,\pm H\}$) or, equivalently, renormalization-group
mode counting at the hierarchical fixed point, where
$\dim E^+_H=H$, $\dim E^-_H=H$, $\dim E^0=1$, yielding $2H+1$ coarse-grained
signed modes. The distinction between $2^H$ (raw independent branch count)
and $2H+1$ (coarse-grained signed modes after Kirchhoff flow conservation)
is the essential content of the axiom system.

\paragraph{Empirical comparison of $z(H)$ alternatives.}
We additionally compute CI match rates on the full 22-indicator
Bettencourt panel for six candidate $z$-functions:
$2H{+}1$ (Kato, $17/22$ at integer resolution; $22/22$ with composition and continuous $\Heff$),
$3H{+}1$ ($18/22$), $H{+}1$ ($14/22$), $2H$ ($17/22$, violates A3),
$H^2{+}1$ ($15/22$), and $2H{+}3$ ($18/22$).
Two alternatives ($3H{+}1$ and $2H{+}3$) achieve one additional integer-resolution
match ($18/22$) by trading the Private R\&D hit for patents and inventors.
However, neither can resolve its remaining misses through composition:
$3H{+}1$ predicts $\betap{2}=1.257$, which falls \emph{below} the Private R\&D
CI $[1.29,1.39]$---a gap that no same-branch mixture can close.
Only the Kato rule $z(H)=2H{+}1$ achieves $22/22$ under composition
corrections and continuous $\Heff$, because all five of its integer-resolution
misses lie in the crossover valley where structural composition provides a
physical resolution.

%%% ============================================================
%%% V29: New appendices from H-uniqueness 5.2/5.4 model comparison
%%% (K7_3.0--7.3 cross-evaluation, note_H_uniqueness_5.2_vs_5.4_model_comparison_2026-04-06.md)
%%% ============================================================

\subsection*{\thesisec. Thermodynamic free-energy proof of relay-depth uniqueness}\label{si:free_energy_uniqueness}\label{si:FG_convexity}

% Source: K7_3.0 (5.2 Pro) — F_G minimization theorem

The main text's axiom-based derivation of $z(H)=2H+1$ does not \emph{by itself}
prove that a finite integer $H$ is selected.  This appendix supplies the
missing selection mechanism via a thermodynamic free-energy argument.

For a city whose production network has graph structure~$G$, define
the free energy at relay depth~$H$ and effective temperature~$T$ by
\begin{equation}
  F_G(H;\,T)\;=\;\frac{\Lambda_G}{H}\;+\;\Gamma_G\,H\;-\;T\ln(2H+1),
  \label{eq:free_energy}
\end{equation}
where $\Lambda_G>0$ is a relay-complexity cost (resistance per layer),
$\Gamma_G>0$ is a maintenance cost per layer, and
$T\ln(2H+1)$ is the configurational entropy gain from the $2H+1$ local
routing states established in the main text.

\begin{proposition}[Conditional $H$-uniqueness]\label{prop:H_unique}
For any $\Lambda_G,\Gamma_G,T>0$, $F_G(H;T)$ possesses a unique global
minimizer $H^{\ast}(G,T)\in(0,\infty)$.
\end{proposition}

\begin{proof}
The first term $\Lambda_G/H$ is strictly convex and diverges as $H\to 0^+$;
the second term $\Gamma_G\,H$ is linear and diverges as $H\to\infty$;
the third term $-T\ln(2H+1)$ is concave and grows only logarithmically.
Hence $F_G(H;T)\to+\infty$ at both extremes, and strict convexity of
$\Lambda_G/H+\Gamma_G H$ minus a concave term guarantees a unique interior
minimum.  The first-order condition
\begin{equation}
  -\frac{\Lambda_G}{(H^{\ast})^2}+\Gamma_G-\frac{2T}{2H^{\ast}+1}=0
  \label{eq:foc}
\end{equation}
determines $H^{\ast}$ implicitly.  Discrete rounding to the nearest integer
then selects $H\in\mathbb{N}$.
\end{proof}

% V46: Added from K9_4.3r1 — F_G convexity and H* uniqueness theorem

\begin{theorem}[Strict convexity of $F_G$ and uniqueness of $H^{\ast}$]\label{thm:F_G_convexity}\label{si:FG_convexity}
For any finite weighted directed production network $G$, the free energy functional
\begin{equation}
F_G(H;T)\;=\;\frac{\Lambda_G}{H}+\Gamma_G H-T\ln(2H+1)
\end{equation}
is strictly convex on the domain $H\in(0,\infty)$. Consequently, there exists a unique global minimizer $H^{\ast}(G,T)\in\mathbb{R}_{>0}$ satisfying $\partial F_G/\partial H=0$, which after integer rounding yields a unique discrete selection $H^{\ast}_{\mathbb{N}}(G,T)\in\mathbb{N}$.

In the high-temperature limit $T\to\infty$, the configuration entropy term dominates and $H^{\ast}$ increases monotonically. In the low-temperature limit $T\to 0^+$, the curvature of the relay-complexity term $\Lambda_G/H$ dominates and $H^{\ast}$ decreases. The intermediate-temperature regime realizes the empirically observed range $H\in\{2,3,4,5\}$ across urban production networks.
\end{theorem}

This result is \emph{conditional} on knowing the network parameters
$(\Lambda_G,\Gamma_G)$ for a given production system.
The open problem of deriving a constructive map from real-world network
topology to $(\Lambda_G,\Gamma_G)$ is discussed in the main text's
``Future Work'' subsection.
We now provide a conditional resolution under explicit homogeneity
assumptions.

% V41: Added from K9_2.0 (5.4 Pro, 30m20s) — Kirchhoff--density identification
\paragraph{Proposition (Kirchhoff--density identification under explicit homogeneity assumptions).}
Let $G=(V,E,w)$ be a finite weighted directed production network with $|V|=n$,
and let $G^{\leftrightarrow}$ denote the symmetrized conductance graph.
Write $R_{ij}^{\leftrightarrow}$ for the effective resistance between nodes
$i$ and $j$ in $G^{\leftrightarrow}$, so that
\[
\mathrm{Kf}(G^{\leftrightarrow})=\sum_{i<j}R_{ij}^{\leftrightarrow}.
\]
Assume:

\begin{enumerate}[label=(K\arabic*)]
\item \emph{Uniform pair traffic:} each unordered node pair contributes equally to the unresolved relay burden;
\item \emph{Stage equipartition:} at relay depth $H$, the unresolved end-to-end burden is evenly distributed across the $H$ serial relay stages;
\item \emph{Homogeneous edge maintenance:} every active directed edge carries the same per-layer maintenance cost.
\end{enumerate}

Then the free-energy coefficients in Eq.~(\ref{eq:free_energy}) admit the exact representations
\begin{equation}
\Lambda_G
=
c_{\Lambda}\,\mathrm{Kf}(G^{\leftrightarrow}),
\qquad
c_{\Lambda}:=\zeta_{\Lambda}\frac{2}{n(n-1)},
\label{eq:Lambda_Kf_exact}
\end{equation}
and
\begin{equation}
\Gamma_G
=
c_{\Gamma}\,\rho_E(G),
\qquad
\rho_E(G):=\frac{|E|}{n(n-1)},
\label{eq:Gamma_rho_exact}
\end{equation}
where $\zeta_{\Lambda}>0$ converts effective resistance into relay-complexity cost
and $c_{\Gamma}>0$ is the common per-edge maintenance cost in density-normalized units.
Consequently,
\begin{equation}
F_G(H;T)
=
\frac{c_{\Lambda}\,\mathrm{Kf}(G^{\leftrightarrow})}{H}
+
c_{\Gamma}\,\rho_E(G)\,H
-
T\ln(2H+1).
\label{eq:free_energy_Kf_rho}
\end{equation}

\begin{proof}
Under uniform pair traffic, the one-step unresolved relay burden is the pair-average
\[
\Lambda_G^{(1)}
=
\frac{\zeta_{\Lambda}}{n(n-1)}
\sum_{i\neq j}R_{ij}^{\leftrightarrow}
=
\zeta_{\Lambda}\frac{2}{n(n-1)}
\sum_{i<j}R_{ij}^{\leftrightarrow}
=
c_{\Lambda}\,\mathrm{Kf}(G^{\leftrightarrow}).
\]
Under stage equipartition, the unresolved burden per relay stage at depth $H$ is
$\Lambda_G^{(1)}/H$, which is exactly the first term of Eq.~(\ref{eq:free_energy}).
For the maintenance term, homogeneous per-edge cost gives the raw layer cost
$\widetilde{\Gamma}_G=c_{\Gamma}|E|$, hence the density-normalized coefficient is
\[
\Gamma_G=\frac{\widetilde{\Gamma}_G}{n(n-1)}=c_{\Gamma}\frac{|E|}{n(n-1)}
=c_{\Gamma}\rho_E(G).
\]
Substituting both identities into Eq.~(\ref{eq:free_energy}) yields
Eq.~(\ref{eq:free_energy_Kf_rho}).
\end{proof}

\paragraph{Remark (what the $1/H$ term does and does not mean).}
Equation~(\ref{eq:Lambda_Kf_exact}) identifies $\Lambda_G/H$ with the
\emph{stage-averaged unresolved relay-complexity burden}, not with the literal
sum of ohmic resistances of the realized $H$-stage chain.
A literal series sum would scale as $H$, whereas the free-energy selector requires
the inverse-depth scaling implied by the D2 closure.

\paragraph{Sharp heterogeneous bound.}
If the pair weights are not uniform, let $q_{ij}\ge 0$ with $\sum_{i<j}q_{ij}=1$.
Then
\begin{equation}
\Lambda_G
=
\zeta_{\Lambda}\sum_{i<j}q_{ij}R_{ij}^{\leftrightarrow}
=
c_G\,\mathrm{Kf}(G^{\leftrightarrow}),
\qquad
c_G:=
\zeta_{\Lambda}\frac{\sum_{i<j}q_{ij}R_{ij}^{\leftrightarrow}}
{\sum_{i<j}R_{ij}^{\leftrightarrow}},
\label{eq:Lambda_Kf_bound_form}
\end{equation}
with the sharp bounds
\begin{equation}
\zeta_{\Lambda}q_{\min}\,\mathrm{Kf}(G^{\leftrightarrow})
\le
\Lambda_G
\le
\zeta_{\Lambda}q_{\max}\,\mathrm{Kf}(G^{\leftrightarrow}),
\label{eq:Lambda_Kf_bound}
\end{equation}
where $q_{\min}:=\min_{i<j}q_{ij}$ and $q_{\max}:=\max_{i<j}q_{ij}$.

% V42: Lifted thermodynamic derivation of K1--K3, from K9_3.3.
% note_file: note_K9_3.3_thermodynamic_K123_2026-04-13.md
\paragraph{Derivation of the homogeneity conditions from a lifted thermodynamic principle.}
The conditional Kirchhoff proposition above assumes three homogeneity
conditions: uniform pair traffic~(K1), stage equipartition~(K2), and
homogeneous edge maintenance~(K3).
We now show that these conditions arise as stationary conditions of a
\emph{lifted} thermodynamic problem.
The reduced selector $F_G(H;T)$ in Eq.~(\ref{eq:free_energy}) is a
coarse-grained functional; at the finer level the pair-burden sector is
entropic whereas the stage and edge sectors are dissipative.
The correct mathematical object is therefore a max--min principle rather
than a bare minimization in all variables.

\begin{lemma}[No-go for exact K1 from linear resistance minimization]
\label{lem:K1_nogo}
Let $\mathcal{P}=\{(i,j):1\le i<j\le n\}$, $m=|\mathcal{P}|=\binom{n}{2}$,
and $R_{ij}^{\leftrightarrow}>0$.
For fixed $Q_{\mathrm{tot}}>0$, the minimizer of
$\sum_{i<j} q_{ij}R_{ij}^{\leftrightarrow}$ subject to
$\sum q_{ij}=Q_{\mathrm{tot}}$
is concentrated on a pair with minimal effective resistance, not uniform.
\end{lemma}

\begin{proposition}[K1 from the maximum-entropy unresolved-pair prior]
\label{prop:K1_maxent}
If unordered node pairs are thermodynamically indistinguishable at the
unresolved-burden level, the maximum-entropy solution under
$\sum p_{ij}=1$ gives $p_{ij}^{\ast}=2/[n(n-1)]$, i.e.,
$q_{ij}^{\ast}=q_0$~(K1).
After K1,
$\Lambda_G^{(1)}=c_\Lambda\,\mathrm{Kf}(G^{\leftrightarrow})$
with $c_\Lambda=\zeta_\Lambda\cdot 2/[n(n-1)]$.
\end{proposition}

\begin{proposition}[K2 from convex stage dissipation]
\label{prop:K2_equip}
By Cauchy--Schwarz,
$\mathcal{D}_H(\{W_h\}):=(1/W_{\mathrm{tot}})\sum_h W_h^2
\ge W_{\mathrm{tot}}/H$, with equality iff $W_h=W_{\mathrm{tot}}/H$~(K2).
The reduced stage cost at the minimum is exactly
$\Lambda_G^{(1)}/H$.
\end{proposition}

\begin{proposition}[K3 from saturated-edge KKT equalization]
\label{prop:K3_kkt}
Under a shared concave capacity law $c(w)$, minimizing maintenance cost
subject to saturated throughput yields $w_e^{\ast}=w_0$ for all active
edges~(K3), and $\Gamma_G=c_\Gamma\rho_E(G)$.
\end{proposition}

\begin{theorem}[Thermodynamic origin of K1--K3]
\label{thm:thermo_K123}
Under the hierarchical max--min principle, K1 emerges from MaxEnt,
K2 from convex equipartition, and K3 from saturated KKT equalization.
After eliminating the fast variables, the reduced free energy recovers
$F_G(H;T)=c_\Lambda\,\mathrm{Kf}(G^{\leftrightarrow})/H
+c_\Gamma\,\rho_E(G)\,H-T\ln(2H+1)$.
\end{theorem}

\begin{remark}[Scope]
The theorem is unconditional for the \emph{lifted} model.
The reduced $F_G(H;T)$ by itself does not imply K1--K3; the
fast-variable microstructure must be restored.
This is consistent with the open statement in the proof architecture
above.
\end{remark}

\paragraph{Numerical consistency check (BEA 71-sector).}
Applying the above identification to the BEA~2024 direct-requirements matrix
($n=71$, symmetrized density $\rho_E=0.588$, off-diagonal density without
threshold $\rho_{\mathrm{raw}}=0.763$), we compute
$\mathrm{Kf}(G^{\leftrightarrow}_{\mathrm{BEA}})=20{,}612$.
At $\tau\approx 2.0$, the cubic equation~\eqref{eq:cubic_H} yields
$H^{*}\approx 4.18$ with $r\approx 10.0$, consistent with the GDP corridor
$H\approx 4$--$5$.
The implied normalization constant $c_{\Lambda}\approx 5\times 10^{-4}$
is of order $1/n^{2}$, suggesting size-invariant relay-complexity scaling.

% V31: Added from K7_10.5 (5.4 Pro, cubic equation + ratio map)
% and Claude F_G simulation (CI hit rate, sensitivity analysis)
\paragraph{Analytical solution via cubic equation.}
Defining the dimensionless ratios $r:=\Lambda_G/\Gamma_G$ and
$\tau:=T/\Gamma_G$, the first-order condition~\eqref{eq:foc} reduces
to a cubic in $H^{\ast}$:
\begin{equation}
  2(H^{\ast})^3 + (1-2\tau)(H^{\ast})^2 - 2r\,H^{\ast} - r = 0.
  \label{eq:cubic_H}
\end{equation}
For $r,\tau>0$ this cubic possesses exactly one positive real root,
which may be expressed via Cardano's formula.
The key monotonicity properties are
$\partial H^{\ast}/\partial r > 0$ (higher relay-complexity cost
increases optimal depth) and $\partial H^{\ast}/\partial\tau > 0$
(higher effective temperature also increases depth).

\paragraph{Discrete rung boundaries.}
An integer $H=n$ is the discrete minimizer of~$F_G$ if and only if
$r_{n-1,n}(\tau)\le r < r_{n,n+1}(\tau)$, where the boundary function is
\begin{equation}
  r_{n,n+1}(\tau) \;=\; n(n+1)\Bigl[1 - \tau\ln\!\frac{2n+3}{2n+1}\Bigr].
  \label{eq:rung_boundary}
\end{equation}
At $\tau=0$ the boundaries simplify to $r_{n,n+1}=n(n+1)$, recovering
a purely cost-driven depth selection.  Table~\ref{tab:rung_map} lists
the resulting $H$-assignment intervals for $\tau=1$.

\begin{table}[h]
\centering\small
\caption{Discrete $H$-rung assignment as a function of the
relay-complexity ratio $r=\Lambda_G/\Gamma_G$ at $\tau=T/\Gamma_G=1$.}
\label{tab:rung_map}
\begin{tabular}{cl}
\toprule
$H$ & $r$ interval \\
\midrule
1 & $[0,\;0.978)$ \\
2 & $[0.978,\;3.981)$ \\
3 & $[3.981,\;8.984)$ \\
4 & $[8.984,\;15.987)$ \\
5 & $[15.987,\;24.988)$ \\
6 & $[24.988,\;35.990)$ \\
\bottomrule
\end{tabular}
\end{table}

\paragraph{Sensitivity analysis.}
At the representative point $(\Lambda_G,\Gamma_G,T)=(10,1,2)$,
with $H^{\ast}\approx 4.18$, the implicit-function sensitivities are:
$\partial H^{\ast}/\partial\Lambda_G \approx +0.156$,
$\partial H^{\ast}/\partial\Gamma_G \approx -2.713$,
$\partial H^{\ast}/\partial T \approx +0.586$.
The maintenance-cost sensitivity is an order of magnitude larger than
the relay-complexity sensitivity, implying that $\Gamma_G$ is the
dominant lever for depth selection.

% V31: Added from K7_10.6 (5.4 Pro, zero-parameter defense)
\paragraph{Scope of the zero-parameter claim.}
The ``zero-parameter'' label in the main text refers specifically to
the exponent selection rule $\beta_{\pm}(H)=1\pm 1/[H\ln(2H+1)]$:
given an independently assigned relay depth~$H$, the predicted
exponent contains no free coefficients.  This is a consequence of the
D2 closure and the structural uniqueness of $z(H)=2H+1$.
The \emph{identification} of~$H$ from network topology---i.e., the
map $G\mapsto H(G)$---is a separate problem that currently involves
auxiliary parameters (the proportionality constants linking
$\Lambda_G$ to the Kirchhoff index and $\Gamma_G$ to the spectral
radius, as well as the effective temperature~$T$).
The temperature~$T$ is not independently identifiable from $F_G$ alone:
$\arg\min_H F_G(H;T) = \arg\min_H [\Lambda_G/(TH) + (\Gamma_G/T)H - \ln(2H+1)]$,
so only the ratios $\Lambda_G/T$ and $\Gamma_G/T$ are effective
degrees of freedom.

% V44: Added from K9_8.1 (dH*/dN theorem) + K9_5.7 (WBE mapping) + K9_5.4 (Kleiber resolution)

\subsection*{\thesisec. Size dependence of the optimal relay depth}\label{si:dHstar_dN}
\stepcounter{sisec}

\paragraph{Theorem (dH*/dN Relation).}
Let the optimal relay depth be selected by the free-energy functional
\begin{equation}
H^{\ast}(N):=\arg\min_{H>0} F_G(H;T,N),\qquad
F_G(H;T,N)=\frac{\Lambda_G(N)}{H}+\Gamma_G H-T\ln(2H+1),
\label{eq:FGN}
\end{equation}
where $N$ is city size, $\Gamma_G>0$ and $T>0$ are size-independent to leading order,
and $\Lambda_G(N)>0$ is a differentiable relay-complexity cost with $\Lambda_G'(N)>0$.
Then $H^{\ast}(N)$ exists uniquely and satisfies
\begin{equation}
\frac{dH^{\ast}}{dN}=\frac{\Lambda_G'(N)/(H^{\ast})^{2}}
{2\Lambda_G(N)/(H^{\ast})^{3}+4T/(2H^{\ast}+1)^{2}}>0.
\label{eq:dHstar_dN_general}
\end{equation}
Hence the optimal relay depth is a strictly increasing function of city size.

\paragraph{Proof sketch.}
Differentiating the first-order condition
$-\Lambda_G(N)/(H^{\ast})^2+\Gamma_G-2T/(2H^{\ast}+1)=0$
with respect to $N$ yields
\[
\underbrace{\left(\frac{2\Lambda_G(N)}{(H^{\ast})^3}+\frac{4T}{(2H^{\ast}+1)^2}\right)}_{=\,F_{HH}>0\text{ at minimizer}}
\frac{dH^{\ast}}{dN}
=\frac{\Lambda_G'(N)}{(H^{\ast})^2}.
\]
Since the left-hand prefactor equals $F_{HH}>0$ at the minimizer
and $\Lambda_G'(N)>0$ by assumption,
$\operatorname{sgn}(dH^{\ast}/dN)=\operatorname{sgn}(\Lambda_G'(N))>0$.
\qed

\paragraph{Implication for $\beta(N)$.}
For a pure positive-branch observable,
$\beta(N)=1+\varepsilon(H^{\ast}(N))=1+1/[H^{\ast}(N)\ln(2H^{\ast}(N)+1)]$.
Since $d\varepsilon/dH<0$, larger cities develop deeper relay hierarchies
while the superlinear excess $\beta(N)-1$ declines toward the linear limit.
For mixed observables:
$\beta_{\mathrm{eff}}(N)=\sum_k \pi_k(N)\,\beta(H_k^{\ast}(N))$.

% V46: Generalised exponent form β = b + s·ε(H) (K9_4.7r1)
\paragraph{Generalised exponent form.}
The positive and negative branches of the Kato ladder can be unified as
\begin{equation}
  \beta = b + s\,\varepsilon(H),\qquad b\in\{0,1\},\quad s\in\{+1,-1\},
  \label{eq:generalised_beta}
\end{equation}
where $b=1$ for extensive-stock indicators (GDP, employment)
and $b=0$ for rate or frequency indicators.
The WBE quarter-power family---metabolic rate ($3/4$),
heart rate ($-1/4$), and lifespan ($1/4$)---all correspond to
$\varepsilon=0.25$ ($H_{\mathrm{eff}}\approx 2.315$) with
$(b,s)=(1,-1)$, $(0,-1)$, and $(0,+1)$ respectively.
This generalisation shows that negative $\beta$ values do not
require ``negative relay depth''; they arise from the same
positive-$H$ hierarchy projected onto rate-type observables.

\subsection*{\thesisec. Mathematical correspondence between WBE spatial dimension and Kato relay depth}\label{si:WBE_Kato_mapping}
\stepcounter{sisec}

West--Brown--Enquist (WBE) allometric theory predicts metabolic scaling
$B\propto M^{d/(d+1)}$ for organisms whose vascular network fills a
$d$-dimensional service volume.
In the Kato framework this correspondence lies \emph{exclusively on the
negative branch}: the relevant exponent is
$\beta^{-}(H)=1-\varepsilon(H)=1-1/[H\ln(2H+1)]$, not $\beta_{+}(H)$.
Setting $d/(d+1)=1-1/[H\ln(2H+1)]$ and simplifying gives
the \emph{mapping equation}
\begin{equation}
H\ln(2H+1) = d+1.
\label{eq:WBE_Kato_map}
\end{equation}
Define $g(H):=H\ln(2H+1)$.
Since $g'(H)=\ln(2H+1)+2H/(2H+1)>0$ for all $H>0$,
$g$ is strictly increasing and admits a unique inverse $g^{-1}$.
The WBE--Kato correspondence is therefore
\begin{equation}
H_{\mathrm{eff}}(d)=g^{-1}(d+1),\qquad
d(H_{\mathrm{eff}})=H_{\mathrm{eff}}\ln(2H_{\mathrm{eff}}+1)-1.
\label{eq:WBE_Kato_inverse}
\end{equation}

\paragraph{Numerical calibration.}
Table~\ref{tab:WBE_Kato_map} lists representative values.
\begin{table}[H]
\centering
\caption{WBE spatial dimension $d$ and the corresponding
Kato effective relay depth $H_{\mathrm{eff}}=g^{-1}(d+1)$.}
\label{tab:WBE_Kato_map}
\begin{tabular}{cccc}
\toprule
$d$ & $H_{\mathrm{eff}}$ & $\beta^{-}_{\mathrm{WBE}}=d/(d+1)$ & $\beta^{-}_{\mathrm{Kato}}=1-\varepsilon(H_{\mathrm{eff}})$ \\
\midrule
1 & 1.4626 & 0.5000 & 0.5000 \\
2 & 1.9083 & 0.6667 & 0.6667 \\
2.5 & 2.1154 & 0.7143 & 0.7143 \\
3 & 2.3148 & 0.7500 & 0.7500 \\
4 & 2.6956 & 0.8000 & 0.8000 \\
\bottomrule
\end{tabular}
\end{table}

\paragraph{Proper-subset relation.}
The WBE achievable exponent set on the real positive axis is
$\mathcal{B}_{\mathrm{WBE}}^{\mathbb{R}^+}=\{d/(d+1):d>0\}=(0,1)$,
whereas Kato spans both branches:
$\mathcal{B}_{\mathrm{Kato}}^{+}
 =\{\beta^{-}(H):H>0\}\cup\{1\}\cup\{\beta^{+}(H):H>0\}=(0,\infty)$.
Therefore $\mathcal{B}_{\mathrm{WBE}}^{\mathbb{R}^+}\subsetneq\mathcal{B}_{\mathrm{Kato}}^{+}$.
In particular, $\beta^{+}(2)=1.311\in\mathcal{B}_{\mathrm{Kato}}^{+}$
but $1.311\notin\mathcal{B}_{\mathrm{WBE}}^{\mathbb{R}^+}$,
confirming the inclusion is proper.

\paragraph{Structural advantages of the Kato framework.}
Beyond the exponent-set inclusion, three structural features distinguish
Kato from WBE:
(i)~Kato possesses an explicit superlinear branch $\beta^{+}(H)>1$,
absent in WBE;
(ii)~the composition rule
$\beta_{\mathrm{eff}}=\sum_k\pi_k\beta(H_k)$
allows mixed-depth predictions without additional postulates;
(iii)~the linear limit $\beta=1$ arises naturally as $H\to\infty$
rather than requiring a special case.

\paragraph{Application to biological allometry: resolving the Kleiber--Rubner debate.}
The mapping Eq.~\eqref{eq:WBE_Kato_map} reinterprets the
classical Kleiber--Rubner controversy as a difference in
$H$-distributions across taxa rather than a conflict between
rival universal exponents.
Kleiber's $\beta=3/4$ corresponds to $d=3$, $H_{\mathrm{eff}}=2.315$,
realised as a mixture $\pi_2\beta^{-}(2)+\pi_3\beta^{-}(3)$
with $\pi_2\approx 0.565$, $\pi_3\approx 0.435$. This identification is
\emph{exclusively negative-branch}; it is not a $\beta_{+}$ assignment and
must not be re-read as a positive-branch $H=3$ observable.
Rubner's $\beta=2/3$ corresponds to $d=2$, $H_{\mathrm{eff}}=1.908$,
realised as nearly pure $H=2$ ($\pi_2\approx 0.962$).
Empirical exponents across taxa---large mammals ($\beta\approx 0.75$),
reptiles and amphibians ($\beta\approx 0.67$),
insects ($\beta\approx 0.82$)---occupy distinct positions
on the same negative-branch ladder, with the insect value
($H_{\mathrm{eff}}\approx 2.90$) reflecting a predominantly $H=3$
vascular architecture.
A na\"{\i}ve maximum-entropy argument---uniform prior over $H=1,\dots,H_{\max}$---yields
$\bar\beta_{\mathrm{eff}}(H_{\max}{=}7)\approx 0.756$, tantalisingly close to~$3/4$,
but this is a finite-support accident: the Ces\`{a}ro mean converges to~$1$
as $H_{\max}\to\infty$ (rate $(\ln\ln n)/n$). The two-component mixture
($\pi_2\approx 0.565$, $\pi_3\approx 0.435$) is therefore structurally more
robust than MaxEnt on uniform~$H$ as an explanation for Kleiber's law.

% V46: Biological relay-depth comparison table from K9_5.5r1, K9_5.6r1
\paragraph{Comparative relay depth across biological systems.}
Table~\ref{tab:bio_relay_depth} collects empirical scaling exponents
and inferred relay depths for three non-urban hierarchical systems.
\begin{table}[H]
\centering
\caption{Biological systems mapped to the Kato relay ladder.}
\label{tab:bio_relay_depth}
\begin{tabular}{lcccc}
\toprule
System & $\beta_{\mathrm{obs}}$ & $H$ & $H_{\mathrm{eff}}$ & Structural basis \\
\midrule
Mammal metabolism   & $0.75$  & $2$--$3$ & $2.3$ & Vascular branching \\
Insect colony       & $0.813$ & $3$      & $2.8$ & Division of labour \\
Slime mold          & ${\sim}0.8$--$1.0$ & $2$ & $1.8$--$2.4$ & Self-organised topology \\
\bottomrule
\end{tabular}
\end{table}
The insect colony exponent $\beta\approx 0.813\approx\beta^{-}(3)=0.829$
lies well within the CI for $H=3$, consistent with the three- to
four-tier division of labour (queen, nurse, worker, forager).
Slime-mold networks (${\sim}36$ nodes, sparse, planar) yield
$H\approx 2$ via the directed Kirchhoff map, matching the
shallow one-rung hierarchy expected for single-organism transport.

% V32: Added from K7_10.9 (5.4 Pro, identification theorem C1-C7)
\subsection*{\thesisec. Conditional identification of relay depth from network topology}\label{si:identification_theorem}
\stepcounter{sisec}

The free-energy framework guarantees the \emph{existence} of an optimal
relay depth~$H^{\ast}$ for any given $(\Lambda_G,\Gamma_G,T)$
(Proposition~\ref{prop:H_unique}).  A separate question is whether the
map $G\mapsto H(G)$ from a production network~$G$ to its relay depth
is \emph{uniquely identified} by observable network structure.
The following theorem, established in consultation with independent
verification (K7\_10.9), formalizes the identification conditions.

\begin{theorem}[Conditional identification]\label{thm:identification}
Let $q$ denote an observable scaling indicator, and let $G_q$ be its
relevant production network.  If $G_q$ satisfies conditions
\textnormal{C1--C7} below, then there exists a unique finite-support
probability measure
\begin{equation}
  \mu_q(G_q) \;=\; \sum_{k=1}^{K}\pi_k\,\delta_{(\sigma_k,\,H_k)},
  \qquad \pi_k>0,\;\sum_k\pi_k=1,\;H_k\in\mathbb{N},
  \label{eq:mu_identification}
\end{equation}
such that the predicted exponent
$\beta_q^{\mathrm{pred}}(G_q)=\sum_k\pi_k\,\beta_{\sigma_k}(H_k)$ is
uniquely determined.  For a pure (single-rung) indicator, $K=1$ and
the integer map $H(G_q)\in\mathbb{N}$ is unique.
\end{theorem}

The seven conditions are:

\begin{description}[labelwidth=2em,leftmargin=2.5em,nosep]
% V46: Added from K9_4.1r1 — ε(H) uniqueness theorem
\item[C1] \textbf{Structural closure.}
  The relay chain obeys D2 serial entropic resistance with state
  count $z(H)=2H+1$, so that
  $\epsilon(H)=1/[H\ln(2H+1)]$ and $\beta_{\pm}(H)=1\pm\epsilon(H)$
  form the unique zero-parameter ladder.

\begin{theorem}[Conditional uniqueness of $\varepsilon(H)$]\label{thm:epsilon_uniqueness}
Under axioms (A1)--(A4)—namely, D2 closure, translation invariance, monotone decay, and continuity—the unique function $\varepsilon:\mathbb{Z}_{\ge 1}\to\mathbb{R}_{>0}$ satisfying the $D_2$-closure condition for serial relay networks is
\begin{equation}
\varepsilon(H)=\frac{1}{H\ln(2H+1)}.
\end{equation}
Consequently, the exponents $\beta_{\pm}(H)=1\pm\varepsilon(H)$ are identifiable zero-parameter predictions without additional tuning parameters, conditional only on independent assignment of relay depth~$H$.
\end{theorem}
\item[C2] \textbf{Observation-class pre-assignment.}
  Each indicator~$q$ is classified, \emph{prior} to observing
  $\beta_{\mathrm{obs}}$, as either (a)~a pure indicator (single
  branch, single rung) or (b)~a mixture indicator with finite support.
\item[C3] \textbf{Independent structural measurement.}
  The relay depth~$H$ (or mixture support $\{(\sigma_k,H_k)\}_{k=1}^K$)
  is determined from structural data alone---firm-size tail indices,
  occupational/educational composition, co-authorship networks,
  CPC/IO upstreamness---independently of $\beta_{\mathrm{obs}}$.
\item[C4] \textbf{Mixture identifiability.}
  For mixture indicators, the independently measured component
  signature matrix~$S$ satisfies
  $\operatorname{rank}[\mathbf{1}^{\top};\,S]=K$.
\item[C5] \textbf{Convergence.}
  The normalized upstream operator~$A_q$ has spectral radius
  $\rho(A_q)<1$, ensuring that the relay expansion converges and
  truncation error is controlled (cf.\ Section~\ref{si:flux_bound}).
\item[C6] \textbf{Measurement-frame fixation.}
  City boundaries, aggregation units, and sample definitions are held
  fixed; all comparisons are conditional on the chosen frame.
\item[C7] \textbf{Invariant-based signatures.}
  The structural measurement in C3 depends only on normalized graph
  invariants: quantities that are unchanged under node relabeling,
  graph isomorphism, and constant rescaling of edge weights.
\end{description}

% V46: Added from K9_4.2r1 — circularity-breaking protocol

\paragraph{Blinded protocol for breaking H--$\beta$ circularity.}
To operationally enforce Conditions C2 and C3, we adopt the following six-step blinded protocol, certified in advance by external referees:
\begin{enumerate}
\item \textbf{Observable identity confirmation.} Register the indicator~$q$ (patents, wages, income, etc.) and classify it as pure (single rung) or mixture before accessing data.
\item \textbf{Structural H estimation.} Measure relay depth $H$ from structural proxies (firm-size tail indices, occupational networks, IO upstreamness, co-authorship structure) using a pre-registered functional form. Do not examine $\beta_{\mathrm{obs}}$ at this stage.
\item \textbf{Prediction calculation.} Compute $\beta_{\mathrm{pred}}$ using $\beta_{\pm}(H)=1\pm\varepsilon(H)$ with the independently measured $H$. Seal the prediction.
\item \textbf{Seal and archive.} Submit both $H$ and $\beta_{\mathrm{pred}}$ to the registry, creating an immutable record with timestamp and referee signature.
\item \textbf{Measurement of observed exponent.} Independently measure $\beta_{\mathrm{obs}}$ from scaling data, blind to the sealed prediction.
\item \textbf{Comparison and CI coverage.} Report the signed discrepancy $\beta_{\mathrm{obs}}-\beta_{\mathrm{pred}}$ and check whether $\beta_{\mathrm{pred}}$ falls within the empirical 95\% CI of $\beta_{\mathrm{obs}}$.
\end{enumerate}
This six-step protocol, executed transparently and pre-registered, eliminates the post-hoc reclassification that would otherwise reintroduce the circularity prohibited by Condition C2.

\noindent\textbf{Key implications.}
(i)~Conditions C1--C3 separate the zero-parameter \emph{prediction}
($\beta_{\pm}(H)$) from the \emph{identification} pipeline
($G\mapsto H$), which may involve auxiliary parameters.
(ii)~The target of identification is generically the measure~$\mu_q$,
not a single integer; the scalar map $H(G)$ is a special case ($K=1$).
(iii)~Dropping any one of C2--C4 produces concrete counterexamples:
without C2, post-hoc reclassification introduces circularity;
without C3, $H$ is fitted to $\beta_{\mathrm{obs}}$ and CI
coverage becomes tautological; without C4, mixture weights are
unidentifiable.

\subsection*{\thesisec. Impossibility of reconciling the cross-sectional ladder with a linear relay-depth trajectory}\label{si:impossibility_linear_H}

% Source: K7_4.5/6.9 — negative result with high paper value

The cross-sectional Kato ladder (this paper) predicts
$\beta(H)=1+1/[H\ln(2H+1)]$, while companion Paper~1 documents a
\emph{longitudinal} parabolic regularity
$\log_{10}(Y_{\rm pc})\propto(\log_{10} N)^2$~\cite{Kato2026a}.
A natural question is whether a single monotonically increasing function
$H(t)$ can generate both observations simultaneously.

\begin{proposition}[Impossibility of linear $H(t)$]\label{prop:impossibility}
If $H(t)=\alpha + \gamma\,t$ with $\gamma>0$ is substituted into the
cross-sectional composition rule, the resulting trajectory
$\beta_{\rm eff}(t)$ converges monotonically toward~$1$ with
$\beta_{\rm eff}-1\sim 1/(t\ln t)$.
This functional form cannot generate the quadratic (parabolic) curvature
in $\log N$ documented in Paper~1.
\end{proposition}

The reconciliation requires allowing the sector-weight vector
$\boldsymbol{\pi}(t)=(\pi_1(t),\dots,\pi_K(t))$ to evolve
endogenously alongside $H$.
A key diagnostic is the ``reverse-engineered'' effective $H$ at world
population $N\approx 8\times 10^9$:  both the 5.2~Pro and 5.4~Pro
models independently estimate $\Heff(N=8\times10^9)\approx 1.15$,
which is \emph{below} the contemporaneous integer assignment $H=5$.
This paradox dissolves once the sector-weight shift toward
high-$H$ service and digital sectors is accounted for:
aggregate $\Heff$ can decrease even as the maximum accessible rung rises.
The derivation of a closed-form $\pi_k(t)$ evolution law remains an
open problem.

\subsection*{\thesisec. Relay-depth saturation and the $\beta\to 1$ limit: decoupling, not extinction}\label{si:beta_one_limit}

% Source: K7_4.8/7.2 — clarification of AI-human decoupling

In the $H\to\infty$ limit, $\epsH{H}\to 0$ and
$\beta_{\rm eff}\to 1$ from above.
A naive reading might interpret $\beta=1$ as a ``singularity'' implying
social or economic collapse.
We emphasize that the Kato equation makes no such prediction:
\begin{enumerate}
\item $\beta\to 1$ is an \emph{extensive} limit.  City output $Y$ scales
  linearly with population $N$, meaning that per-capita output
  $Y/N$ becomes independent of city size.  This is
  \emph{proportionate scaling}, not stagnation or decline.
\item Human extinction requires not merely $\beta\to 1$ (loss of
  super-linear agglomeration) but an independent \emph{driver}:
  either $L_{\rm pop}\to 0$ (population collapse) or $C_h\to 0$
  (destruction of human capital infrastructure).
  The Kato equation does not predict either of these.
\item With AI-mediated relay chains at $H\gg 5$, human-generated
  outputs ($H\leq 5$) and AI-generated outputs ($H>5$) simply
  \emph{decouple}: the AI sector contributes ever more agglomeration
  surplus while the human sector converges to $\beta\approx 1$.
  This is a structural separation, not an apocalyptic threshold.
\end{enumerate}
The transition from the current $H\approx 4$--5 corridor to an
AI-dominated regime is best modeled as a \emph{crossover} rather than
a sharp phase transition; the effective band of $H$ values broadens
continuously rather than jumping.

\subsection*{\thesisec. Cognitive ceiling at $H=5$ and the Miller--Dunbar bound}\label{si:cognitive_ceiling}

% Source: K7_4.3/6.7 — H=5 as institutional saturation

The integer relay depth $H=5$ occupies a distinguished position in the
ladder: it is the highest rung for which all mediating relays can plausibly
be staffed by unaugmented human cognition.

Two independent constraints converge at $H\approx 5$:
\begin{itemize}
\item \textbf{Miller's $7\pm 2$ working-memory bound}~\cite{Miller1956}:
  a single human can reliably track $5$--$9$ independent items.
  At $H=5$, each relay node manages $2H+1=11$ local routing states,
  which is at the upper edge of this capacity.
\item \textbf{Dunbar's social-brain limit}~\cite{Dunbar1992}:
  stable social groups cap at ${\sim}150$ individuals, constraining
  the organizational span available for multi-layer relay coordination.
  With $H=5$ layers, a Dunbar-sized organization can sustain roughly
  $150/5=30$ independent channels per layer---adequate for a modern
  firm or institutional department, but not for the ${\sim}100$
  channels needed at $H\geq 6$.
\end{itemize}

The empirical observation that raw human cognition achieves at most
$H\approx 4$ while institutional and informational augmentation extends
reach to $H=5$ is consistent with both the Bettencourt panel
(no indicators assigned $H>5$) and the GDP composition analysis
(aggregate $\Heff\approx 4$--5).
Reaching $H\geq 6$ requires either artificial-intelligence-mediated
relaying or qualitatively new institutional architectures---the domain
explored in the preceding appendix.

% V33: Added from K7_9.8 (directed Kf extension) + Claude BEA I-O sim
\subsection*{\thesisec. Directed extension of the Kirchhoff relay-complexity cost}\label{si:directed_Kf}
\stepcounter{sisec}

The relay-complexity parameter $\Lambda_G\propto\mathrm{Kf}(G)$ is
defined via the Kirchhoff index of an \emph{undirected} graph.
Real production networks, however, are inherently directed: the
BEA input-output table, for example, is a $71\times 71$ digraph
with asymmetry index
$\eta_{\mathrm{skew}}:=\lVert A-A^{\top}\rVert_F/\lVert A\rVert_F
\approx 1.35$, indicating that nearly all inter-sector flows are
one-way.
Three directed Kirchhoff candidates are natural:

\begin{description}[nosep]
\item[$\mathrm{Kf}_{\mathrm{dir}}^{\mathrm{sym}}$]
  Symmetrize first: $A_{\mathrm{sym}}=(A+A^{\top})/2$, then apply the
  standard formula $n\,\mathrm{tr}(L_{\mathrm{sym}}^{+})$.
  This discards all skew-symmetric information.
\item[$\mathrm{Kf}_{\mathrm{dir}}^{\mathrm{CT}}$]
  Compute hitting times $h_{ij}=(Z_{jj}-Z_{ij})/\pi_j$ from the
  fundamental matrix
  $Z=(I-P+\mathbf{1}\pi^{\top})^{-1}$, where $P=D_{\mathrm{out}}^{-1}A$
  and $\pi$ is the stationary distribution, then sum commute times:
  $\mathrm{Kf}_{\mathrm{dir}}^{\mathrm{CT}}
   = m_A^{-1}\sum_{i<j}(h_{ij}+h_{ji})$.
  When $A=A^{\top}$ this reduces exactly to $\mathrm{Kf}(G)$.
\item[$\mathrm{Kf}_{\mathrm{dir}}^{\leftrightarrow}$]
  Build the symmetrized reversible Laplacian
  $L^{\leftrightarrow}=\tfrac12[\Pi(I-P)+(I-P)^{\top}\Pi]$ and
  set $\mathrm{Kf}_{\mathrm{dir}}^{\leftrightarrow}=n\,\mathrm{tr}
  ((L^{\leftrightarrow})^{+})$.
\end{description}

\noindent\textbf{Numerical result on BEA 2024 data.}
Applying all three candidates to the 71-sector direct-requirements
matrix (self-loops removed, negative entries zeroed) yields
$\mathrm{Kf}_{\mathrm{dir}}^{\mathrm{CT}}/
 \mathrm{Kf}_{\mathrm{dir}}^{\mathrm{sym}}=0.949$:
the commute-time index is approximately~$5\%$ lower than the
symmetrized baseline.
After calibrating $c_{\Lambda}$ so that
$H^{\ast}_{\mathrm{sym}}\!=4.42$ (the GDP metropolitan mean),
the directed correction shifts $H^{\ast}$ by at most~$-0.11$,
well within the rounding tolerance for integer
$H$.
All 22~Bettencourt indicators retain their previously assigned
integer relay depths, and the CI hit rate remains $17/22$~(77.3\%).
\emph{Directed graph structure does not break the Kato ladder;}
the symmetric approximation is a robust proxy for $\Lambda_G$.

\noindent\textbf{Caveat.}
The reversible-projection candidate
$\mathrm{Kf}_{\mathrm{dir}}^{\leftrightarrow}$ is sensitive to
zero-out-degree nodes (four government-enterprise sectors in the
BEA table); teleportation regularization inflates
$\mathrm{Kf}_{\mathrm{dir}}^{\leftrightarrow}$ by two orders of
magnitude and should be applied with caution.

\paragraph{Biological network example: slime mold.}
The adaptive network of \emph{Physarum polycephalum}
\citep{Tero2010} contains ${\sim}\,36$ nodes and exhibits a sparse,
planar topology.  Applying the directed Kirchhoff map yields
$H\approx 2$, consistent with the shallow one-rung hierarchy
expected for a single-organism transport network.
This confirms the biological scaling anchor: vascular and colony
networks occupy the $H\!=\!2$ rung, producing Rubner-like
$\beta\approx 2/3$ exponents.

\subsection*{\thesisec. Network renormalization-group derivation of relay depth}\label{si:network_RG_H}

% Source: K7_2.2 (5.2 Pro Extended) — RG stopping time theorem
% note_file: note_topology_to_H_prior_results_K7_2x_2026-04-06.md

The preceding appendices treat $H$ as either externally assigned
or selected by free-energy minimization with phenomenological parameters
$(\Lambda_G,\Gamma_G)$.
This appendix sketches a first-principles derivation in which both
$H$ \emph{and} the state-count factor $\ln(2H+1)$ emerge from a
network renormalization-group (RG) procedure.

\paragraph{Setup.}
Let $G=(V,E,w)$ be a weighted directed graph representing a city's
production network (e.g.\ a $402\times 402$ BEA input--output matrix).
Define a coarse-graining transformation $\mathcal{T}$ that contracts
each bottleneck community into a single super-node while preserving
net flow:
\begin{equation}
  G_{n+1}=\mathcal{T}(G_n),\qquad G_0\equiv G.
  \label{eq:rg_iteration}
\end{equation}
We say $G$ is \emph{RG-admissible} if the sequence $\{G_n\}$ converges
in finitely many steps to a fixed-point graph $G^{\ast}$ (a trivial
single-node graph or a self-similar motif).

% V42: Formal T specification from K9_2.1r1 (RG operator axiom system A0--A6)
\paragraph{Formal specification of the coarse-graining operator.}
The operator $\mathcal{T}$ is fixed by the following axiom system:
\begin{itemize}
\item[\textbf{A0.}] $G_n=(V_n,E_n,w_n)$ is a finite weighted directed graph
  with $w^{(n)}_{ij}\ge 0$ and total flow
  $m_n:=\sum_{i,j\in V_n}w^{(n)}_{ij}<\infty$.
\item[\textbf{A1.}] (Resolution unity) $\gamma\equiv 1$: the directed
  null-model total flow equals the observed total flow, giving the unique
  normalization under which $\mathcal{T}$ is flow-preserving.
\item[\textbf{A2.}] (Exact directed partition) Let
  $Q_{\mathrm{dir}}(P;G_n):=m_n^{-1}\!\sum_{a}\sum_{i,j\in C_a}
  (w^{(n)}_{ij}-k^{\mathrm{out}}_i k^{\mathrm{in}}_j/m_n)$.
  $P_n^{\star}:=\arg\max_{P\in\mathfrak{A}(G_n)}Q_{\mathrm{dir}}(P;G_n)$
  is the unique maximizer over admissible (weakly-connected, laminar,
  canonical-tie-break) partitions.
\item[\textbf{A3.}] (Flow preservation) $w^{(n+1)}_{ab}=\sum_{i\in C_a}
  \sum_{j\in C_b}w^{(n)}_{ij}$, so $m_{n+1}=m_n$.
\item[\textbf{A4.}] (Directed information retention) No symmetrization is
  applied; the skewness $\eta_{\mathrm{skew}}$ is retained at each level.
\item[\textbf{A5.}] (Fixed-point halt) If $P_n^{\star}$ is trivial
  (zero modularity gain), set $\mathcal{T}(G_n)=G_n=G^{\ast}$.
\item[\textbf{A6.}] (Non-degeneracy) Adjacent rungs of the free-energy
  landscape have no ties, ensuring uniqueness of the fixed point and the
  discrete minimizer.
\end{itemize}
Under A0--A6, $\mathcal{T}$ satisfies the three RG axioms: (a)~semigroup
composition $\mathcal{T}^{m+n}=\mathcal{T}^m\circ\mathcal{T}^n$
(by definition); (b)~fixed-point existence and uniqueness (conditional on
A2$+$A6); (c)~solver universality within the same objective (any exact solver
of $Q_{\mathrm{dir}}$ yields the same sequence $\{G_n\}$).

\paragraph{RG stopping time as relay depth.}
\begin{definition}[Network relay depth]\label{def:H_RG}
For an RG-admissible network $G$, the \emph{relay depth} is the
stopping time
\begin{equation}
  H(G)\;:=\;\min\bigl\{n\geq 0 : G_n = G^{\ast}\bigr\}.
  \label{eq:H_stopping}
\end{equation}
\end{definition}

By construction, $H(G)\in\mathbb{N}$---integrality is automatic.
The Kato axiom A1 (serial chain additivity) is recovered because
$\mathcal{T}$ removes exactly one relay layer per step:
$H(G)=H(G_1)+1$.

% V42: Exact equivalence theorem from K9_2.1r1
\begin{theorem}[Recursive community depth equals RG stopping time]
\label{thm:depth_equals_H}
Under axioms A0--A6, the recursive directed-community depth
$n_L^{\mathrm{exact}}(G)$ (obtained by iterating exact
$Q_{\mathrm{dir}}$-maximization) equals the RG stopping time:
$n_L^{\mathrm{exact}}(G)=H(G)$.
If a heuristic solver (e.g.\ directed Louvain) deviates by at most one
merge or split per level, then
$|n_L^{\mathrm{heur}}(G)-H(G)|\le 1$.
\end{theorem}
\begin{proof}
Each recursive step applies $\mathcal{T}$ exactly once under A0--A6,
so the recursion halts after exactly $H(G)$ steps.
The heuristic bound follows by summing per-level deviations, each
bounded by one laminar merge or split (A2 approximate version).
\end{proof}

\paragraph{Mode counting at the fixed point.}
At each RG step $n$, linearize $\mathcal{T}$ around $G^{\ast}$ and
decompose the perturbation spectrum into positive-deviation modes
$E_n^{+}$, negative-deviation modes $E_n^{-}$, and a single neutral
mode $E^{0}$.
For a serial relay chain of depth $H$:
\begin{equation}
  \dim E_H^{+}=H,\qquad \dim E_H^{-}=H,\qquad \dim E^{0}=1,
  \label{eq:mode_counting}
\end{equation}
yielding a total mode count of $H+H+1=2H+1=z(H)$.
This is precisely the signed state space $\Sigma_H=\{0,\pm 1,\ldots,\pm H\}$
derived axiomatically in the main text.
The Shannon path entropy is then
\begin{equation}
  S(H)=H\ln(2H+1),
  \label{eq:entropy_rg}
\end{equation}
and the Kato correction follows as
$\epsH{H}=1/S(H)=1/[H\ln(2H+1)]$.

\paragraph{Significance.}
The RG derivation provides the first route to $H$ and $\epsH{H}$ that
does not import any external categorical label.
Both ingredients---the integer depth $H$ and the logarithmic factor
$\ln(2H+1)$---arise from the network's own hierarchical structure
under iterative coarse-graining.

\paragraph{Alternative algorithmic routes to~$H$.}
% Source: K7\_8.2--8.9, K7\_9.0--9.4 (5.2 Pro Extended batch)
% note\_file: note\_topology\_to\_H\_5.2pro\_batch\_K7\_8x9x\_2026-04-06.md
The RG coarse-graining in Definition~\ref{def:H_RG} is not the only
possible $G\\!\to\\!H$ algorithm.
At least six additional proposals have been formulated, each recovering
integer relay depth from a different structural signature of the
production network:
(i)~\emph{DAG longest-path depth} from an acyclicized flow graph;
(ii)~\emph{information-bottleneck (IB) quantization}, where the
  elbow of the rate-distortion curve singles out $2H{+}1$ clusters;
(iii)~\emph{persistent-homology bar counting}, reading $H$ from
  the number of long-lived $H_0$ bars in a filtration;
(iv)~\emph{max-flow / min-cut layer decomposition}, counting
  successive bottleneck layers in a capacitated digraph;
(v)~\emph{spectral gap ladder}, reading $H$ from the number of
  significant gaps in the Laplacian spectrum of the symmetrized graph;
(vi)~\emph{transfer-entropy depth}, identifying $H$ as the
  number of Granger-causal layers in a time-resolved I--O flow.
%
Each method recovers $H\in\mathbb{N}$ from distinct data modalities, but their relative merits have not been
systematically benchmarked on the same city-level network datasets.
Comparative evaluation is left for future work.

%--- Open problems and limitations -------------------------------------------------------
\subsection{Open problems and limitations}\label{si:limitations}

\begin{enumerate}
\item The categorical assignment of $H$ to indicators remains analyst-guided.
  Network-anchored assignment (Section~\ref{si:network_RG_H}) provides a
  structural alternative but has not yet been validated on real city-level
  interaction data.
\item The explicit free-energy functional $F_G(H;T)$ whose minimum selects
  $H^*$ is postulated, not derived from first principles. The convexity
  argument (Appendix~\ref{si:FG_convexity}) shows that a minimum exists but
  does not pin down the functional form.
\item Mixture dynamics $\pi_k(t)$ governing the temporal evolution of
  composition weights are unspecified. The cross-sectional ladder is
  time-independent, but predicting \emph{changes} in $\beta$ over time
  requires a model for $\pi_k(t)$.
\item The constructive map from network topology to the cost parameters
  $(\Lambda_G,\Gamma_G)$ remains open.
\item Validation is limited to OECD, EU, Chinese, and Japanese cities.
  Extension to low-income urban systems awaits city-level data with
  sufficient geographic coverage.
\end{enumerate}


\subsection*{\thesisec. Microfoundation of $\Lambda_G$: effective-resistance relay-complexity}\label{si:Lambda_G_microfoundation}

This appendix sharpens the microscopic meaning of the free-energy coefficient
$\Lambda_G$ in
\begin{equation}
F_G(H;T)=\frac{\Lambda_G}{H}+\Gamma_G H-T\ln(2H+1).
\label{eq:FG_lambda_micro}
\end{equation}
The $\Lambda_G/H$ term is the \emph{stage-averaged unresolved relay-complexity
burden}: a communication-length / effective-resistance cost, not the literal
sum of wiring lengths (that role is carried by $\Gamma_G H$, and entropy is
carried separately by $-T\ln(2H+1)$).

\paragraph{Kirchhoff microfoundation under explicit homogeneity.}
Under three assumptions---(i) uniform pair traffic, (ii) stage equipartition,
(iii) homogeneous edge maintenance---the coefficient is
\begin{equation}
\Lambda_G \;=\; c_{\Lambda}\,\mathrm{Kf}(G^{\leftrightarrow}),\qquad
c_{\Lambda}:=\zeta_{\Lambda}\frac{2}{n(n-1)},
\label{eq:Lambda_Kf_law}
\end{equation}
where $\mathrm{Kf}(G^{\leftrightarrow})=\sum_{i<j}R_{ij}^{\leftrightarrow}$ is
the Kirchhoff index of the symmetrized directed graph and $\zeta_{\Lambda}>0$
a unit-traffic pair cost.  Stage equipartition at depth $H$ then gives
\begin{equation}
\frac{\Lambda_G}{H}=\frac{c_{\Lambda}\mathrm{Kf}(G^{\leftrightarrow})}{H},
\end{equation}
recovering the first term of~\eqref{eq:FG_lambda_micro}.

\paragraph{Equivalent A-type form and rejected candidates.}
Equation~\eqref{eq:Lambda_Kf_law} is an A-type microscopic form once
literal total wire-length is replaced by \emph{effective-resistance wiring
length}:
$\Lambda_G=\lambda_0 L_G^{\mathrm{eff}}$, with
$L_G^{\mathrm{eff}}:=\mathrm{Kf}(G^{\leftrightarrow})$ and $\lambda_0:=c_{\Lambda}$.
A purely metabolic form $\Lambda_G=\mu_0 M_G$ (B-type) is not the coefficient
identified by the present free-energy derivation, and a redundancy-based form
$\Lambda_G=\kappa_0 C_G^{-1}$ (C-type) would misplace into $\Lambda_G$ a role
already assigned to $-T\ln(2H+1)$.  Both B and C are therefore rejected under
\eqref{eq:FG_lambda_micro}.

\paragraph{Heterogeneous-pair generalization.}
If pair traffic weights are $q_{ij}\ge 0$ with $\sum_{i<j}q_{ij}=1$, then
$\Lambda_G=c_G\,\mathrm{Kf}(G^{\leftrightarrow})$ with
$c_G=\zeta_{\Lambda}\sum_{i<j}q_{ij}R_{ij}^{\leftrightarrow}/\sum_{i<j}R_{ij}^{\leftrightarrow}$.

\paragraph{Convexity preservation.}
Equation~\eqref{eq:Lambda_Kf_law} gives $\Lambda_G>0$; combined with
$\Gamma_G>0,T>0$ (below), $F_G$ remains strictly convex in $H$,
$\partial^{2}F_G/\partial H^{2}=2\Lambda_G/H^{3}+4T/(2H+1)^{2}>0$, and the
first-order condition $-\Lambda_G/(H^{\ast})^{2}+\Gamma_G-2T/(2H^{\ast}+1)=0$
determines a unique $H^{\ast}$.

\paragraph{BEA 71-sector consistency check.}
For the archived BEA illustration,
$\mathrm{Kf}(G_{\mathrm{BEA}}^{\leftrightarrow})=2.06\times 10^{4}$ and
$c_{\Lambda}\approx 5\times 10^{-4}$ yield $\Lambda_G\approx 10.3$,
$r:=\Lambda_G/\Gamma_G\approx 10$, and $H^{\ast}\approx 4.18$, consistent with
the independently measured GDP corridor $H_{\mathrm{eff}}\in[4.02,4.42]$
(Sec.~\ref{si:dHstar_dN}).

\paragraph{Size dependence.}
If $\Lambda_G=\Lambda_G(N)$ with $\Lambda_G^{\prime}(N)>0$, then
\begin{equation}
\frac{dH^{\ast}}{dN}=\frac{\Lambda_G^{\prime}(N)/(H^{\ast})^{2}}
{2\Lambda_G(N)/(H^{\ast})^{3}+4T/(2H^{\ast}+1)^{2}}>0,
\end{equation}
so larger cities (more production nodes) raise $\Lambda_G$ and deepen the
optimal relay hierarchy.

\paragraph{Scope and open problems.}
This appendix removes $\Lambda_G$ from the list of free phenomenological
parameters: once $G$ is specified, $\Lambda_G$ is fixed by $\mathrm{Kf}(G^{\leftrightarrow})$.
What remains open is the elimination of $c_{\Lambda}$ (a unit-traffic scale)
and the independent identification of $T$.

\subsection*{\thesisec. Microfoundation of $\Gamma_G$: active-edge maintenance density}\label{si:Gamma_G_microfoundation}

This appendix pairs with Sec.~\ref{si:Lambda_G_microfoundation} to fix the
microscopic reading of the linear penalty coefficient $\Gamma_G$ in
Eq.~\eqref{eq:FG_lambda_micro}.  We show that $\Gamma_G$ is neither a latency
nor a redundancy parameter: it is a \emph{density-normalized active-edge
maintenance coefficient}.

\begin{proposition}[Density-normalized edge-maintenance law]
\label{prop:Gamma_density}
Let $G=(V,E,w)$ be a finite weighted directed production network with
$|V|=n$.  Assume (i)~uniform pair traffic, (ii)~stage equipartition, and
(iii)~homogeneous edge maintenance under a shared concave capacity law
$c(w)$ at saturated throughput.  Then
\begin{equation}
\Gamma_G=c_{\Gamma}\,\rho_E(G),\qquad
\rho_E(G):=\frac{|E|}{n(n-1)},
\label{eq:Gamma_density_micro}
\end{equation}
with $c_{\Gamma}:=c(w_0)>0$ the per-edge maintenance cost at the saturated
optimum $w_e^{\ast}=w_0$.
\end{proposition}

\begin{proof}
Define $\widetilde{\Gamma}_G:=\sum_{e\in E} c(w_e)$.  KKT optimality under
(iii) yields $w_e^{\ast}=w_0$ for every active edge, hence
$\widetilde{\Gamma}_G=c_{\Gamma}|E|$.  Density normalization
$\Gamma_G=\widetilde{\Gamma}_G/[n(n-1)]$ gives Eq.~\eqref{eq:Gamma_density_micro}.\qed
\end{proof}

\paragraph{Rejected candidates.}
A latency form $\Gamma_G=\tau_0\nu_G$ (per-stage processing delay) is rejected
because the manuscript's free-energy does not contain a latency variable.
A redundancy form $\Gamma_G=\eta_0 r_G$ (error-correction overhead) is
similarly absent.  Only the edge-maintenance form~\eqref{eq:Gamma_density_micro}
is compatible with Eq.~\eqref{eq:FG_lambda_micro}.

\paragraph{Joint selector and the $\Lambda_G/\Gamma_G$ ratio.}
Combining Eqs.~\eqref{eq:Lambda_Kf_law} and~\eqref{eq:Gamma_density_micro},
the $T\to 0$ stationary depth satisfies
$H^{\ast}\big|_{T=0}=\sqrt{\Lambda_G/\Gamma_G}
=\sqrt{c_{\Lambda}\mathrm{Kf}(G^{\leftrightarrow})/[c_{\Gamma}\rho_E(G)]}$.
The dimensionless ratio
\begin{equation}
r:=\frac{\Lambda_G}{\Gamma_G}=\frac{c_{\Lambda}}{c_{\Gamma}}\,
\frac{\mathrm{Kf}(G^{\leftrightarrow})}{\rho_E(G)}
\end{equation}
compares \emph{global unresolved burden} to \emph{local active-edge upkeep};
it is a pure graph invariant up to the unit-scale factor $c_{\Lambda}/c_{\Gamma}$.

\paragraph{BEA completion of the consistency check.}
Eq.~\eqref{eq:Gamma_density_micro} yields $\Gamma_G\approx 1.03$ for the BEA
71-sector network, with $c_{\Gamma}\approx 1.75$ estimated from the calibrated
$r=\Lambda_G/\Gamma_G\approx 10$.  Together with Sec.~\ref{si:Lambda_G_microfoundation},
the BEA illustration now fixes every coefficient in $F_G(H;T)$ except the
unit-traffic scale $c_{\Lambda}$ and the effective temperature $T$.

\paragraph{Zero-parameter status.}
At the graph-invariant level, the relay-depth selector is now parameter-free:
once $G$ is specified, $\mathrm{Kf}(G^{\leftrightarrow})$ and $\rho_E(G)$ are
determined, and therefore so are $\Lambda_G$ and $\Gamma_G$.  The two
conversion constants $(c_{\Lambda},c_{\Gamma})$ are not theoretical free
parameters of the ladder structure itself; they are unit-scale factors fixed
by edge and traffic conventions.  This completes the microfoundation program
for the coefficients of $F_G$.

\subsection*{\thesisec. Limitations update: from postulated to microfounded $F_G$}\label{si:limitations_update_V56}\label{si:Lambda_Gamma_microfoundation}

Under Sections~\ref{si:Lambda_G_microfoundation} and~\ref{si:Gamma_G_microfoundation}, the earlier statement that ``$F_G$ is a postulated functional with phenomenological coefficients $(\Lambda_G,\Gamma_G)$ whose reduction to the underlying production graph $G$ remains open'' should be read as resolved at the graph-invariant level: $\Lambda_G=c_\Lambda\,\mathrm{Kf}(G^{\leftrightarrow})$ and $\Gamma_G=c_\Gamma\,\rho_E(G)$ under the uniform-traffic / stage-equipartition / homogeneous-maintenance assumptions.  What remains open is the elimination of the unit-scale factors $(c_\Lambda,c_\Gamma)$ and the independent identification of the effective temperature $T$.
\begin{enumerate}
\item Validation is limited to OECD, EU, Chinese, and Japanese cities.
  Extension to low-income urban systems awaits city-level data with
  sufficient geographic coverage.
\end{enumerate}

\subsection*{\thesisec. Derivation of the temporal--spatial identity $\betap{1}=2-\DS(1)$ (Theorem 2 of main text)}\label{si:H1_vonFoerster_identity}

\begin{theorem}[$H=1$ \,\emph{as the $D_2$-renormalized pairwise law}]
Within the Kato ladder $\betap{H}=1+\epsH{H}$ with primitive $\epsH{H}=1/[H\ln(2H+1)]$, the minimal positive rung $H=1$ satisfies the structural identity
\[
\betap{1}\;=\;2-\DS(1),\qquad \DS(1)\;=\;1-\frac{1}{\ln 3}\;\approx\;0.089761,
\]
and $\DS(H)=(1/H)[1-1/\ln(2H+1)]$. The numerical proximity $\betap{1}\approx 2$ is therefore a structural consequence of entropy renormalization, not a fitting coincidence.
\end{theorem}

\begin{proof}[Proof outline]
Kato side: $S(1)=\ln 3$, $\epsH{1}=1/\ln 3$, $\betap{1}=1+1/\ln 3$. Pairwise class identification: at $H=1$ in the CW-skeleton formulation only 0-cells (agents) and 1-cells (pairwise links) are present, so the minimal positive Kato rung is exactly the minimal pairwise relay class. Pairwise scaling on the canonical one-relay representative $S_N=K_{1,N-1}$ (star graph) gives $Y_1(N)=(b/2)N^{2}+O(N)\sim N^{2}$, the same quadratic combinatorial order as von~Foerster $p=2$. The mean pairwise distance satisfies $\bar{d}(S_N)=2-2/N\to 2$. The $D_2$ selector replaces the unit pairwise excess $1$ by the admissible one-relay excess $\epsH{1}=1/\ln 3$, yielding $\betap{1}=1+1/\ln 3=2-\DS(1)$. Full source lines: paper3\_nhb\_V43 Thm~9.1 and paper3\_SI\_V50 Appendix; derived here under the K9\_13.5\,r1 identity agreement (A1/K10-2 2026-04-14).\qed
\end{proof}

\begin{remark}[Why $p=2$ is not a literal $H=0$ rung]
The comparison exponent $p=2$ should not be read as a literal Kato rung at $H=0$: it is the \emph{unrenormalized} pairwise exponent. The Kato $H=1$ rung is its minimal serial-relay realization after $D_2$ entropy renormalization, with discount $\DS(1)=1-1/\ln 3\approx 0.0898$.
\end{remark}

\begin{remark}[Cross-sectional vs temporal]
The Kato object is the cross-sectional exponent selection $\beff=\sum_k \pi_k \betap{H_k}$. The von~Foerster object is a temporal growth law. The identity asserts that both are generated by the same pairwise interaction skeleton; $H=1$ is the $D_2$-renormalized cross-sectional representative of that skeleton.
\end{remark}

\subsection*{\thesisec. Properties of the per-rung entropy cost $\DS(H)$}\label{si:Delta_S_properties}

The per-rung entropy cost is
\[
\DS(H)\;=\;\frac{1}{H}\left[1-\frac{1}{\ln(2H+1)}\right]\;=\;\frac{1}{H}-\epsH{H}.
\]

\begin{proposition}[Integer-rung monotonicity]
$\DS$ is \emph{not} globally monotone in the continuous variable $H\in(0,\infty)$: it has a local maximum in the open interval $(1,2)$ with $\DS(1)=0.0898<\DS(2)=0.1893$. However, restricted to integer rungs $H\ge 2$, $\DS$ is strictly decreasing: $\DS(2)>\DS(3)>\DS(4)>\cdots$, with $\DS(H)\to 0$ as $H\to\infty$.
\end{proposition}

\begin{proof}[Proof]
Direct computation: $\DS(1)=1-1/\ln 3=0.0898$, $\DS(2)=0.5\cdot(1-1/\ln 5)=0.1893$, $\DS(3)=(1/3)\cdot(1-1/\ln 7)\approx 0.1620$, $\DS(4)=0.25\cdot(1-1/\ln 9)\approx 0.1362$. The continuous derivative $d\DS/dH$ is positive on $(1,H_c)$ and negative on $(H_c,\infty)$ for some $H_c\in(1,2)$; restriction to integer argument yields monotonic decay from $H=2$ onward. The limit $\DS(H)\to 0$ follows from $1/\ln(2H+1)\to 0$ and $1/H\to 0$.\qed
\end{proof}

\begin{proposition}[Relation to the companion ceiling gap $\Dceil$]
The companion paper~\cite{Kato2026c} uses $\Dceil(H):= p-\betap{H}=1-\epsH{H}$. For all $H\in\mathbb{N}_{\ge 1}$,
\[
\Dceil(H)-\DS(H)\;=\;\frac{H-1}{H},
\]
so the two gaps coincide only at $H=1$ and differ by a bounded,
monotonically increasing offset for $H\ge 2$.
\end{proposition}

\begin{proof}
Direct: $\Dceil(H) = 1 - \epsH{H} = 1 - 1/[H\ln(2H+1)]$, and $\DS(H) = (1/H)(1-1/\ln(2H+1)) = 1/H - 1/[H\ln(2H+1)]$. Subtraction gives $\Dceil(H) - \DS(H) = 1 - 1/H = (H-1)/H$. \qed
\end{proof}

\paragraph{General-$H$ bridge formula.}\phantomsection\label{si:P3_P4_general_bridge}
The general bridge theorem quoted in the main text follows immediately from the
identity above.
\begin{theorem}[General $H$ pairwise--Kato bridge]
For every $H\in\mathbb{N}_{\ge 1}$,
\begin{equation}
\betap{H}=2-\DS(H)-\frac{H-1}{H}.
\end{equation}
\end{theorem}
\begin{proof}
By definition, $\Dceil(H)=p-\betap{H}\big|_{p=2}=2-\betap{H}$. Using
$\Dceil(H)-\DS(H)=(H-1)/H$, we obtain
$2-\betap{H}=\DS(H)+(H-1)/H$, hence
$\betap{H}=2-\DS(H)-(H-1)/H$. At $H=1$ the second loss vanishes, recovering
Theorem~\ref{si:H1_vonFoerster_identity}.\qedhere
\end{proof}

\refstepcounter{sisec}
\subsection*{\thesisec. $p$-exclusion theorem for the spatial ladder}\label{si:p_exclusion_theorem}
For cross-paper citation convenience we record the minimal exclusion statement
used by the companion population paper.

\begin{proposition}[$p$-exclusion for finite spatial rungs]
For every finite relay depth $H\in\mathbb{N}_{\ge 1}$,
\begin{equation}
1<\betap{H}<2,
\qquad
0<\betam{H}<1.
\end{equation}
Consequently, the temporal pairwise exponent $p=2$ is not itself a finite Kato
spatial rung and cannot be promoted to a free selector variable inside the
spatial ladder; it remains the unrenormalized pairwise ceiling used only in the
companion temporal framework.
\end{proposition}
\begin{proof}
For $H\ge 1$, the primitive satisfies $0<\epsH{H}=1/[H\ln(2H+1)]<1$ because
$H\ln(2H+1)>1$. Therefore $\betap{H}=1+\epsH{H}\in(1,2)$ and
$\betam{H}=1-\epsH{H}\in(0,1)$. Thus no finite relay rung equals the temporal
ceiling $p=2$, and the latter must be treated as external to the spatial
selector.\qedhere
\end{proof}


% ====================================================================
% V66 additions start here (K9_118 batch, Paper 4 SA reviewer defense).
% Five new SI sections, all inserted before the legacy Graphical Abstract
% block so V65 content downstream is unchanged:
%   §S-MP  Master-parameter ledger  (K9_118.1)
%   §S-SPL Sum-of-power-laws theorem (K9_118.5)
%   §S-CB-diag Finite-window OLS diagnostic (K9_118.3)
%   §si:H_network_uniqueness Network-anchored relay depth (K9_118.4)
%   §S-RA  Reviewer-anticipation audit (K9_118.9 Q6.2)
% ====================================================================

\refstepcounter{sisec}
\subsection*{\thesisec. Master-parameter ledger and reproducibility safeguards}
\label{si:SMP_master_parameter_ledger}

All exponent predictions in this paper are generated from one primitive
relay-depth correction,
\[
\varepsilon(H)=\frac{1}{H\ln(2H+1)},\qquad H\in\mathbb{N}_{\ge1},
\]
with the signed branches
\[
\beta_{\pm}(H)=1\pm\varepsilon(H).
\]
The relay depth \(H\) is a discrete structural rung: it counts irreducible
serial information-transforming relays and is assigned before
\(\beta_{\mathrm{obs}}\) is inspected. Continuous \(H_{\mathrm{eff}}\) is not a
primitive depth. It is a derived coordinate used only after a same-branch or
crossover mixture has been structurally specified. For such a mixture,
\[
\varepsilon(H_{\mathrm{eff}})=\sum_k \pi_k\varepsilon(H_k),
\]
so composition is performed in correction space and only then inverted to
\(H_{\mathrm{eff}}\). If a fractional display coordinate is needed, define
\(b_H:=H_{\mathrm{eff}}-\lfloor H_{\mathrm{eff}}\rfloor\in[0,1)\); this
symbol is distinct from the generalized-exponent baseline
\(b\in\{0,1\}\) in \(\beta=b+s\varepsilon(H)\).

For mixed observables we write
\[
\beta_{\mathrm{eff}}=\sum_k\pi_k\beta_{\sigma_k}(H_k),
\qquad \pi_k\ge0,\qquad \sum_k\pi_k=1,\qquad \sigma_k\in\{+,-\}.
\]
Two kinds of mixture weight must be kept separate. In the canonical
sum-of-subsectors identity, \(\pi_k=Y_k/Y\) is an output share. In the
relay-level cross-branch model, \(\pi_k=w_k/\sum_jw_j\) is a structural
pathway share independent of city size. In the same-depth two-branch
specialization,
\[
\pi_++\pi_-=1,\qquad
\sigma_{gc}:=\pi_+-\pi_-\in[-1,1],
\]
and
\[
\beta_{\mathrm{eff}}
 =\pi_+\beta_+(H)+\pi_-\beta_-(H)
 =1+\sigma_{gc}\varepsilon(H).
\]
Thus \(\sigma_{gc}=0\) is the exact finite-depth linear attractor produced by
balanced generation and coverage pathways. It is not a zero-noise assumption
and should not be confused with residual standard deviations denoted by
\(\sigma\) in statistical simulations.

The reproducibility protocol is: (i) record the independent structural
evidence used to assign \(H_k\), \(\sigma_k\), and \(\pi_k\); (ii) compute
\(\beta_{\mathrm{pred}}\) from the fixed ladder before opening
\(\beta_{\mathrm{obs}}\); (iii) store the processed data, figure-generation
scripts, random seeds, software versions, and output hashes in the public
repository; and (iv) report all sensitivity scans separately from fixed
theoretical constants. The exponent map itself has no fitted coefficient once
\(H\), branch sign, and mixture support are fixed. Statistical nuisance
quantities such as intercepts, residual spreads, bootstrap seeds, and BIC
component counts are therefore documented as reproducibility metadata rather
than as parameters of the Kato ladder.


\refstepcounter{sisec}
\subsection*{\thesisec. Sum-of-Power-Laws Theorem}
\label{si:sum_power_laws_theorem}

This section formalizes the curvature implication of the composition rule
used in the main text. The purpose is to distinguish a true single power law
from an observable that is the aggregate of several relay-depth components.

\begin{theorem}[Sum-of-power-laws]
\label{thm:si_sum_power_laws}
Let
\[
Y(N)=\sum_{k=1}^{K} w_k N^{\beta_k},
\qquad
N>0,\quad w_k>0,\quad \sum_k w_k=1,
\]
with distinct exponents \(\beta_k\). Let \(x=\log N\) and
\(F(x)=\log Y(e^x)\). Define
\[
\pi_k(N)=
\frac{w_k N^{\beta_k}}{\sum_j w_j N^{\beta_j}},
\qquad
\beta_{\mathrm{loc}}(N)=\frac{d\log Y}{d\log N}.
\]
Then
\[
\beta_{\mathrm{loc}}(N)=\sum_k \pi_k(N)\beta_k,
\]
and
\[
\frac{d^2\log Y}{d(\log N)^2}
=
\sum_k \pi_k(N)
\left[\beta_k-\beta_{\mathrm{loc}}(N)\right]^2.
\]
Thus a heterogeneous sum of power laws is not exactly linear on a log--log
plot: it has positive curvature whenever at least two active components have
distinct exponents. Moreover,
\[
\lim_{N\to\infty}\beta_{\mathrm{loc}}(N)=\max_k\beta_k.
\]
\end{theorem}

\begin{proof}
Write
\[
Z(x)=\sum_k w_k e^{\beta_k x},
\qquad F(x)=\log Z(x).
\]
Then
\[
F'(x)=\frac{Z'(x)}{Z(x)}
=
\frac{\sum_k w_k\beta_k e^{\beta_k x}}
     {\sum_j w_j e^{\beta_j x}}
=
\sum_k \pi_k(N)\beta_k.
\]
Also
\[
\pi_k'(x)=\pi_k(x)\left[\beta_k-F'(x)\right].
\]
Therefore
\[
F''(x)
=
\sum_k\beta_k\pi_k'
=
\sum_k\pi_k\beta_k\left[\beta_k-F'(x)\right]
=
\sum_k\pi_k\left[\beta_k-F'(x)\right]^2.
\]
This is a variance and is positive unless all active exponents are equal.
For the large-\(N\) limit, factor out \(N^{\beta_{\max}}\) from numerator and
denominator of \(\pi_k(N)\); the component with \(\beta_{\max}\) receives
probability one as \(N\to\infty\). \qedhere
\end{proof}

\begin{remark}[Finite-window OLS]
The identity above is exact for the local log--log slope. A finite-window
OLS slope is generally an average of the local slope over the observed
population range, not the value at one arbitrary \(N\). For a uniform
continuous design \(x\in[a,b]\),
\[
\beta_{\mathrm{OLS}}[a,b]
=
\int_a^b
\frac{6(x-a)(b-x)}{(b-a)^3}
\beta_{\mathrm{loc}}(e^x)\,dx.
\]
Hence a narrow-window OLS slope converges to the local slope, whereas a
wide window can hide curvature by averaging over changing component shares.
\end{remark}

\begin{corollary}[\(\beta\approx1\) ambiguity]
A near-linear observed exponent does not by itself imply \(H\to\infty\).
It may arise from either a single component with \(\beta_k\approx1\), or
from finite-depth branch cancellation:
\[
\sum_k \pi_k(N)(\beta_k-1)\approx0.
\]
In the same-depth Kato two-branch case,
\[
\pi_+\beta_+(H)+\pi_-\beta_-(H)
=
1+(\pi_+-\pi_-)\epsilon(H),
\qquad
\pi_++\pi_-=1,
\]
so exact balance \(\pi_+=\pi_-=1/2\) gives \(\beta_{\mathrm{eff}}=1\) at
finite \(H\). This is the mechanism used for the near-linear indicators in
Table~\ref{tab:panel_primary_si}.
\end{corollary}

\paragraph{Application to the 22-indicator panel.}
The theorem supports the current interpretation of mixed observables in
Table~\ref{tab:panel_primary_si}. GDP is treated in the revised manuscript
as an \(H=4\)--5 dominated mixed observable, with national
\(\beta_{\mathrm{eff}}=1.113\) and metropolitan mean
\(\beta_{\mathrm{eff}}=1.099\), rather than as a shallow \(H\approx3\)
single-service observable. Patents remain the canonical \(H=2\) innovation
anchor, with aggregate patent panels in 2001 and 2010 remaining consistent
with \(\beta_+(2)\). Road surface is the clean negative-branch example:
\(\beta_{\mathrm{obs}}=0.830\) matches \(\beta_-(3)=0.829\).

\paragraph{Diagnostics for single-term versus multi-term structure.}
For any indicator with city-level data, we recommend three pre-registered
diagnostics.

\begin{enumerate}
\item \emph{Curvature test.} Fit
\[
\log Y_i=\alpha+\beta x_i+\kappa(x_i-\bar x)^2+\varepsilon_i,
\qquad x_i=\log N_i.
\]
A single power law predicts \(\kappa=0\); a heterogeneous output-share sum
of powers predicts \(\kappa>0\).

\item \emph{Population-bin slope test.} Split cities into low-, middle-,
and high-\(N\) bins and estimate \(\beta_{\mathrm{bin}}\). A single term
predicts stable bin slopes; a sum of distinct powers predicts systematic
drift across bins.

\item \emph{Residual-dispersion test.} Compare single-rung, two-rung, and
Kato-constrained mixture models by RMSE, AIC/BIC, and
\[
\mathrm{CV}_{|r|}:=\frac{\operatorname{sd}(|r_i|)}
{\operatorname{mean}(|r_i|)}.
\]
The multi-term model is supported only if it reduces residual dispersion
without introducing unconstrained fitted slopes.
\end{enumerate}

\paragraph{Reviewer-facing interpretation.}
The theorem is not a generic basis-expansion claim. The admissible
\(\beta_k\) values are fixed by the relay-depth ladder
\(\beta_{\pm}(H)=1\pm[H\ln(2H+1)]^{-1}\), and the component weights must be
measured from sector, pathway, CPC, occupation, or input--output structure
before the observed slope is opened. Under that protocol, curvature and
finite-window slope drift are falsifiable consequences of independently
specified network topology, not tautological curve fitting.


\refstepcounter{sisec}
\subsection*{\thesisec. Finite-window OLS diagnostic for literal power-sum mixtures}
\label{si:cross_branch_ols_diagnostic}

The primary cross-branch model in this SI treats \(\pi_k\) as a structural
relay-pathway share. This subsection gives a separate diagnostic for the
case in which an observed aggregate is literally approximated by a finite
sum of power laws.

\begin{proposition}[Finite-window branch-balance diagnostic]
\label{prop:finite_window_branch_balance}
Let
\[
Y(x)=\sum_{k=1}^{K} w_k e^{\beta_k x},\qquad x=\log N,\qquad
w_k>0,\qquad \sum_k w_k=1 .
\]
Fix a reference point \(x_0=\langle x\rangle\), write \(u=x-x_0\), and
define
\[
\alpha_k(x_0)=
\frac{w_k e^{\beta_k x_0}}
{\sum_j w_j e^{\beta_j x_0}},
\qquad
\sum_k\alpha_k(x_0)=1 .
\]
Then the local log-log slope is
\[
\bar\beta(x_0)=\sum_k\alpha_k(x_0)\beta_k .
\]
Moreover, over a centered finite window,
\[
\hat\beta_{\mathrm{OLS}}
=
\bar\beta(x_0)
+
O\!\left(
\operatorname{Var}_{\alpha}(\beta)
\frac{\langle u^3\rangle}{\langle u^2\rangle}
+
\langle |u|^3\rangle
\right).
\]
Consequently, to leading order,
\[
\hat\beta_{\mathrm{OLS}}\approx 1
\quad\Longleftrightarrow\quad
A_+\approx A_-,
\]
where
\[
A_+=\sum_{\beta_k>1}\alpha_k(x_0)(\beta_k-1),\qquad
A_-=\sum_{\beta_k<1}\alpha_k(x_0)(1-\beta_k).
\]
Equivalently,
\[
\pi_+^{\Delta}:=\frac{A_+}{A_++A_-}
\approx
\pi_-^{\Delta}:=\frac{A_-}{A_++A_-}
\approx \frac{1}{2}.
\]
\end{proposition}

\begin{proof}
Expand \(Y\) around \(x_0\):
\[
Y(x_0+u)=
\sum_k w_k e^{\beta_k x_0}e^{\beta_k u}
=
Y(x_0)\sum_k\alpha_k(x_0)e^{\beta_k u}.
\]
Taking the logarithm gives the cumulant-generating function of the exponent
distribution under weights \(\alpha_k(x_0)\):
\[
\log Y(x_0+u)
=
\log Y(x_0)
+
\bar\beta u
+
\frac{1}{2}\operatorname{Var}_{\alpha}(\beta)u^2
+
\frac{1}{6}\kappa_{3,\alpha}(\beta)u^3
+
O(u^4).
\]
The OLS slope of \(\log Y\) on \(x\) is
\[
\hat\beta_{\mathrm{OLS}}
=
\frac{\operatorname{Cov}(u,\log Y(x_0+u))}
{\operatorname{Var}(u)} .
\]
Substituting the Taylor expansion yields the stated approximation. Finally,
\[
\bar\beta-1
=
\sum_k\alpha_k(\beta_k-1)
=
\sum_{\beta_k>1}\alpha_k(\beta_k-1)
-
\sum_{\beta_k<1}\alpha_k(1-\beta_k)
=
A_+-A_-.
\]
Thus the finite-window slope is near one exactly when positive and negative
branch deviations balance to leading order.\qedhere
\end{proof}

\begin{remark}[Why raw \(\beta\)-weighted branch shares are not enough]
The quantities that must balance are the excess deviations
\((\beta_k-1)\) and \((1-\beta_k)\), not the raw sums
\(\sum_{\beta_k\gtrless1} w_k\beta_k\). Raw exponent-weighted branch shares
can be equal even when the net correction away from one is not zero.
For this reason the diagnostic balance variables are
\(\pi_\pm^\Delta\), defined from \(A_\pm\), rather than raw exponent shares.
\end{remark}

\begin{remark}[Diagnostic separation from \(H\to\infty\)]
For a pure \(H\to\infty\) linear limit, all active exponents satisfy
\(\beta_k\approx 1\), so \(\operatorname{Var}_{\alpha}(\beta)\approx0\).
For finite-depth cancellation, \(\bar\beta\approx1\) but
\(\operatorname{Var}_{\alpha}(\beta)>0\). Residual curvature and
population-bin drift therefore distinguish the two mechanisms.
\end{remark}


\refstepcounter{sisec}
\subsection*{\thesisec. Network-anchored relay depth and conditional \(H\) uniqueness}
\label{si:H_network_uniqueness}

\begin{definition}[Relay depth of an indicator]
Let \(Y(N)\) be an urban indicator generated by a typed relay chain
\[
X_0 \xrightarrow{\mathcal A_1} X_1
\xrightarrow{\mathcal A_2}\cdots
\xrightarrow{\mathcal A_H} X_H=Y ,
\]
where \(X_0\) is the primitive agent-level or event-level field and each
\(\mathcal A_h\) is an irreducible aggregation, routing, accounting, or
reporting layer.  The relay depth \(H\) is the number of noncontractible
serial relays in the minimal quotient of this chain. Parallel operations
inside one layer are contracted; serial transformations that change the
information type, institutional frame, or measurement shell are retained.
Under hub scaling, each relay contributes an additive log-slope increment:
\[
\frac{d\log X_h}{d\log N}
=
\frac{d\log X_{h-1}}{d\log N}+\varepsilon_h .
\]
Thus, if \(X_0(N)\asymp N\),
\[
Y(N)\asymp
N^{\,1+\sum_{h=1}^{H}\varepsilon_h}.
\]
In the homogeneous substrate limit this becomes
\(Y(N)\asymp N^{1+H\varepsilon}\).
The Kato closure fixes the total correction as
\[
\sum_{h=1}^{H}\varepsilon_h
=
\epsilon(H)=\frac{1}{H\ln(2H+1)},
\]
yielding the two admissible pure-rung branches
\[
\beta_{+}(H)=1+\epsilon(H),\qquad
\beta_{-}(H)=1-\epsilon(H).
\]
Therefore the exponent map has no free slope parameter once the relay
topology has fixed \(H\).
\end{definition}

\begin{conjecture}[Conditional topology-anchored uniqueness]
For a fixed typed indicator-generation network \(\mathcal R_Y\), let
\(Q(\mathcal R_Y)\) be the minimal serial quotient obtained by contracting
within-layer parallel operations while preserving every irreducible
type-changing relay. Then the primitive relay depth
\[
H(Y)=\#\{\text{noncontractible serial relays in }Q(\mathcal R_Y)\}
\]
is unique up to relabeling of isomorphic layers. Non-integer
\(H_{\mathrm{eff}}\) is not a primitive topology. It is a derived crossover
coordinate:
\[
\epsilon(H_{\mathrm{eff}})
=
(1-b)\epsilon(H)+b\,\epsilon(H+1),
\qquad 0\le b<1 .
\]
The limit \(b=0\) recovers an exact integer rung, while \(0<b<1\) represents
partial occupation of the adjacent network regime.
\end{conjecture}

The RG interpretation is that integer \(H\) values are fixed points of a
coarse-graining flow on the indicator-generation network. If
\(G_{n+1}=\mathcal T(G_n)\) contracts each bottleneck community while
preserving net flow, then for an RG-admissible network the stopping time
\[
H(G)=\min\{n:G_n=G^\ast\}
\]
is integer by construction. A stable integer regime corresponds to one
fixed point of this flow. A fractional \(H_{\mathrm{eff}}\) corresponds
instead to a crossover occupation measure, for example
\((1-b)\delta_H+b\delta_{H+1}\), not to a freely fitted half-layer.
This is the sense in which the assignment is network-anchored: \(H\) is
read from the indicator's generation topology before comparing with
\(\beta_{\mathrm{obs}}\).

Applied to the present panel, aggregate patents have an \(H=2\) innovation
core: inventor/problem-solving interaction \(\rightarrow\) patent document
and claim classification \(\rightarrow\) metropolitan patent count. GDP has
a coarse \(H=3\) accounting lower bound---production \(\rightarrow\) firm
or transaction accounting \(\rightarrow\) value-added aggregation
\(\rightarrow\) statistical account---but the present manuscript refines
GDP into a deep \(H=3,4,5\) composition and therefore places it in the
\(H=4\)--5 corridor. A primitive road-length count would be \(H=1\) if it
were only a direct length measurement; the Bettencourt road-surface row is
instead a serviced infrastructure-network observable and is treated as a
negative-branch \(H\approx3\) case. This distinction is falsifiable: if
future harmonized panels show a stable \(H_{\mathrm{eff}}=2.5\) across
city definitions, countries, estimators, and decades, the proper conclusion
is not to tune \(H\) continuously but to decide among (i) an adjacent-rung
crossover, (ii) a measurement artifact, or (iii) a genuinely intermediate
network regime requiring refinement of the relay-depth theory.


\refstepcounter{sisec}
\subsection*{\thesisec. Reviewer-anticipation audit: assignment, flexibility, and scope}
\label{si:reviewer_anticipation}

\paragraph{Assignment versus exponent prediction.}
The principal circularity risk is the assignment of \(H\), not the exponent
map itself. This section therefore records the admissible workflow:
(i) specify the observable and measurement frame; (ii) assign integer
\(H\) or structural mixture weights from non-slope evidence where
available; (iii) compute \(\beta_{\mathrm{pred}}\); and only then
(iv) compare with \(\beta_{\mathrm{obs}}\). We distinguish structural
evidence (CPC composition, input--output layering, enterprise-size and
occupational structure, graph topology) from slope-audit evidence (ML,
BIC, and nearest-rung diagnostics). The former supports anti-post-hoc
assignment; the latter checks internal coherence.

\paragraph{Integer, composition, and continuous \(H_{\mathrm{eff}}\) score ledger.}
We report three score levels rather than a single headline number. Level~1
is the integer test, \(17/22\). Level~2 applies predeclared composition
corrections to crossover rows and yields \(21/22\). Level~3 reports
continuous \(H_{\mathrm{eff}}\) only as a derived coordinate for
structurally identified crossover cases, yielding \(22/22\). A
McNemar-style paired accounting for Level~1 to Level~2 has four
No\(\to\)Yes changes and zero Yes\(\to\)No changes; exact two-sided
\(p=0.125\), so it is used as a transparent paired-score description rather
than as a primary significance test.

\paragraph{Small-sample and boundary stress tests.}
Because the canonical panel has only 22 rows, this section collects the
dependence-robust checks in one place: leave-one-out stability, repeated
5-fold CV, domain-block bootstrap, ICC-adjusted effective sample size,
hierarchical random-intercept residuals, and random-permutation tests. It
also defines the outlier-gate rule for CPC Section~D, Japan GDP, and other
single-row stress cases. Boundary sensitivity is handled by conditional
invariance: FUA, LUZ, MSA, national, and metropolitan frames are not
pooled unless the measurement frame and mixture weights are explicitly
harmonized.


% ====================================================================
% V61 (2026-04-15): Graphical Abstract Design, civilization-scale instances
% (K9_27.0), multi-observable rank status.
% V62 (2026-04-15): added K9_27.1 (5DAO H=5 stability), K9_27.2 (hexagonal
% H=6 management shell), K9_28.0 (throughput ν identification, conditional
% rank-4 closure).
% ====================================================================

\section{Graphical Abstract Design}
\label{si:graphical_abstract_design}

The submission-ready graphical abstract uses a three-panel left-to-right
narrative: \emph{Input network topology $G$} $\rightarrow$ \emph{Kato primitive
and relay-depth selector} $\rightarrow$ \emph{zero-parameter prediction}. The
visual hierarchy follows the manuscript's main claims: the primitive
correction is $\epsH{H}=1/[H\ln(2H+1)]$, the branch structure is
$\betap{H},\betam{H}=1\pm\epsH{H}$, mixed observables use
$\beff=\sum_k \pi_k \beta(H_k)$, GDP lies in an $H=4$--$5$ corridor, and the
22-indicator confrontation improves from $17/22$ to $21/22$ and then $22/22$
after composition and continuous $H_{\mathrm{eff}}$ refinement. The selector is
the free-energy minimizer $H^{\ast}=\arg\min_H F_G(H;T)$ with
$F_G(H;T)=\Lambda_G/H+\Gamma_G H-T\ln(2H+1)$. The palette is journal-safe and
color-blind-aware: blue for input topology, navy/gold for theory and
selection, and muted red/green for predictions and fit status. The full TikZ
source is archived separately as \texttt{graphical\_abstract\_V3.tex}.

\section{Civilization-Scale Instances of the Kato Ladder}
\label{si:civilization_scale_instances}

This section records three civilization-scale instances of the relay-depth
ladder that appear in a companion non-technical monograph. They are provided
not as additional empirical evidence for Paper~4 but as worked examples of
the universal lens implied by the framework.

\subsection{Six-layer coin as $H=6$ economic implementation}

\begin{definition}[Six-layer coin structure]
The six-layer coin structure (CPT-T, Pow-T, Fiat-T, Gov-T, Orbit-T, Galaxy-T)
of the companion monograph is an $H=6$ implementation of the Kato ladder in
which the local state count is $z(6)=2\cdot 6+1=13$: six usage-right tokens,
six voting-right tokens, and one neutral fiat center. Under this
identification, $\epsH{6}=1/[6\ln 13]\approx 0.06498$ and
$S(6)=6\ln 13$.
\end{definition}

\begin{proposition}[Decimal rigidity theorem]
The $10{:}1$ denomination compression rule of the monograph is serial
additive in the denomination space and therefore introduces no new exponent
parameter. It is equivalent, at the ladder level, to rigid coupling of the
six usage-right layers.
\end{proposition}

\begin{proposition}[Free-energy partition of the six-layer coin]
Let $H_c:= \sqrt{\Lambda_G/\Gamma_G}$. The CPT-T and Pow-T layers contribute to the $\Lambda_G/H$ term (effective-resistance cost of usage rights), while Gov-T, Orbit-T, and Galaxy-T contribute to the $\Gamma_G H$ term (active-edge maintenance of voting rights). Fiat-T implements entropy injection at the $T\ln(2H+1)$ level.
\end{proposition}

\subsection{Five-tier DAO as $H=5$ governance truncation with stability
theorem}
\label{si:K9_27_1_5DAO_H5_stability}

The five-tier DAO pyramid (Christian three-tier $\to$ small-state rung $\to$
galactic upper tier) of the monograph is the $H=5$ governance truncation of
the six-layer coin. We record the stability conditions and the
``four-versus-six'' theorem that accompany this identification.

\begin{definition}[Governance stability]
A governance architecture at depth $H\in\mathbb{N}$ is \emph{stable} if all
three hold: (i) \emph{selector stability}---$H$ is the unique discrete
minimizer of $F(H)=\Lambda/H+\Gamma H-T\ln(2H+1)$; (ii) \emph{functional
completeness}---the realization map $M_H:\mathbb{R}^H\to\mathbb{R}^5$ has full
rank $5$; (iii) \emph{balance compatibility}---there exists a legitimacy
tolerance $\eta\in(0,2\epsH{5})$ with $|\beff-1|\le\eta$.
\end{definition}

\begin{theorem}[Five-tier stability]
Under conditions (A)--(D) of the companion monograph (Kato primitive,
five-coordinate completeness, finite-depth justification, cognitive ceiling
$H\le 5$), and for selector parameters in the local window
\[
\frac{\Lambda}{30}+T\ln\!\frac{13}{11}\;<\;\Gamma\;<\;\frac{\Lambda}{20}+T\ln\!\frac{11}{9},
\]
the value $H=5$ is the unique stable relay depth. In particular,
$F(5)-F(4)<0$ (four-tier is under-resolved, $\operatorname{rank} M_4<5$) and
$F(6)-F(5)>0$ (six-tier is over-resolved beyond human cognitive capacity).
\end{theorem}

\begin{corollary}[Four-under, six-over]
Four-tier architectures fail functional completeness; six-tier architectures
exceed human cognitive capacity. The five-tier DAO sits at the unique
interior minimum of $F$ on the human-admissible range $\{1,\dots,5\}$.
\end{corollary}

\begin{proposition}[Veto as cross-branch balance surrogate]
Each tier veto corresponds to an instance of the $\epsH{H}$ primitive at its
rung. The institutional $33.4\%$ veto threshold is a discretization of the
finite-depth cross-branch balance $\pi_+=\pi_-=1/2$ inside
$\beff=1+(\pi_+-\pi_-)\epsH{H}$; it is the smallest legally enforceable
blocking fraction that prevents $\pi_-\to 0$, not an alternative bookkeeping
of $50\%$.
\end{proposition}

\subsection{Hexagonal management shell at fixed $H=6$}
\label{si:K9_27_2_6type_hexagonal}

The six-type entrepreneur taxonomy of the monograph (three binary axes
$\times$ three-level discretization) populates the six vertices of a
hexagonal shell at fixed $H=6$.

\begin{proposition}[Hexagonal mixture]
Let a managerial style be represented by three normalized coordinates
$\mathbf{x}=(s,d,c)\in[0,1]^3$ with $s$ service orientation, $d$ top-down
orientation, and $c$ capital propensity. A six-state discretization of this
cube projects onto the six vertices of a hexagon that lives entirely on the
$H=6$ level set of the Kato ladder. Individual-firm longitudinal
trajectories across vertices are captured by time-varying mixture weights
$\pi_k(t)$ along this level set.
\end{proposition}

\begin{proposition}[Management scaling law]
For a fixed-shell hexagonal mixture at $H=6$,
$\betap{6}\approx 1.065$ and $\betam{6}\approx 0.935$, with
$\epsH{6}\approx 0.065$. Large mature firms converging to the $H=6$ shell
are predicted to exhibit $\beta_{\mathrm{mgmt}}\approx 1.065$; small
growing firms spread along an $H=4$--$5$ corridor
($\beta_{\mathrm{mgmt}}\in[1.083,1.114]$) before depth-up convergence.
\end{proposition}

\begin{remark}[Book-facing objects as one Kato architecture]
The six-layer coin shell ($H=6$), five-tier DAO shell ($H=5$), hexagonal
executive taxonomy ($C_6$ at fixed $H=6$), and $33.4\%$ veto are four
distinct projections of the same underlying Kato architecture.
\end{remark}

These identifications are structural and do not introduce empirical claims
beyond those already stated in the 22-indicator confrontation; they are
included to make explicit the cross-domain range of the relay-depth ladder.

\section{Multi-Observable Rank and Identification Status}
\label{si:multi_obs_rank_status}
\label{si:scale_gauge_rank_bound}

The selector-only primitive triple $(\Lambda_G/T,\;\Gamma_G/T,\;c)$ attached
to a single GDP-indicator produces a Jacobian of rank at most two by
scale-gauge invariance (Section~\ref{si:scale_gauge_rank_bound}). Two
complementary resolutions are available:

\begin{enumerate}
\item \emph{Landauer normalization.} Imposing
$T_{\mathrm{info}}=k_B T_{\mathrm{phys}}$ at the thermodynamic level unlocks
rank-$3$ closure by anchoring the scale gauge to a universal physical
constant (Section~\ref{si:Landauer_absolute_scale}).
\item \emph{Multi-observable route.} If three graph-independent observables
$(G_1,G_2,G_3)$ (e.g., BEA input-output, CPC product classification,
infrastructure network) have linearly independent
$(\Kf(G_i),\rhoE(G_i))$ pairs, their joint selector system has
rank $3$ without external physical anchoring. The formal independence
conditions and a BEA-CPC-Infrastructure numerical check are deferred to a
forthcoming note (K9\_21.3 retry, in preparation).
\end{enumerate}

Either route closes the identification problem at rank $3$; the Landauer
route is the one invoked in the body of this SI.

\section{Throughput Identification: a Single Clock Observable Closes
bits$/$s and Watt}
\label{si:throughput_nu_identification}
\label{si:Landauer_absolute_scale}

This section resolves the missing identification step left open by the
Landauer-normalized absolute-scale closure. The present source stack already
fixes the thermodynamic scale,
\[
T_{\mathrm{info}}=k_B T_{\mathrm{phys}},\qquad
T_{\mathrm{BEA}}\approx 4.14\times 10^{-21}\,\mathrm{J},\qquad
H^{\ast}_{\mathrm{BEA}}\approx 4.18,
\]
and the serial path entropy $S(H)=H\ln(2H+1)$. Therefore the only missing
ingredient for throughput is not another energy parameter but one clock
observable.

\subsection{Candidate clocks and relay-cascade time}

Among five plausible candidates---(A) annual update, (B) revision cadence,
(C) transaction cadence, (D) completed cascade rate $\nu_{\mathrm{cas}}$,
(E) primitive stage clock $\nu_0$---only (D) and (E) respect the
serial-relay physics of the present source stack; (A) and (B) are
administrative reporting cadences, and (C) is inconsistent with the 71-sector
BEA graph resolution.

\begin{proposition}[Relay-cascade time]
Under the serial-relay architecture with stage equipartition at depth $H$,
the completed-cascade time is additive in stage count,
\begin{equation}
\tau_{\mathrm{cas}}(H)=H\,\tau_0,\qquad
\nu_{\mathrm{cas}}=1/\tau_{\mathrm{cas}}=\nu_0/H.
\label{eq:tau_cascade_linear}
\end{equation}
Diffusive ($\tau=H^2\tau_0$) and entropy-disguised
($\tau=H\ln(2H+1)\tau_0$) alternatives are both excluded by the
architecture.
\end{proposition}

\subsection{Stage versus cascade bookkeeping equivalence}

Let $b_{\mathrm{stage}}(H):=\log_2(2H+1)$,
$q_{\mathrm{stage}}(H):=T_{\mathrm{info}}\ln(2H+1)$,
$b_{\mathrm{cas}}(H):=H\log_2(2H+1)$,
$q_{\mathrm{cas}}(H):=H\,T_{\mathrm{info}}\ln(2H+1)$.

\begin{proposition}[Stage/cascade equivalence]
Under Eq.~\eqref{eq:tau_cascade_linear}, the two bookkeeping conventions
coincide:
\begin{equation}
\dot I=\nu_0\,b_{\mathrm{stage}}(H)=\nu_{\mathrm{cas}}\,b_{\mathrm{cas}}(H),\qquad
W=\nu_0\,q_{\mathrm{stage}}(H)=\nu_{\mathrm{cas}}\,q_{\mathrm{cas}}(H).
\label{eq:stage_cascade_equivalence}
\end{equation}
This resolves the apparent ambiguity between the $\log_2(2H+1)$ form (one
stage) and $H\log_2(2H+1)$ form (one completed cascade).
\end{proposition}

\subsection{BEA 71-sector reference numbers}

At $H^{\ast}_{\mathrm{BEA}}\approx 4.18$,
$b_{\mathrm{stage}}^{\mathrm{BEA}}\approx 3.23$ bits,
$q_{\mathrm{stage}}^{\mathrm{BEA}}\approx 9.26\times 10^{-21}$ J;
$b_{\mathrm{cas}}^{\mathrm{BEA}}\approx 13.49$ bits,
$q_{\mathrm{cas}}^{\mathrm{BEA}}\approx 3.87\times 10^{-20}$ J. Per-sector
averages over the 71-sector network are
$\bar b_{\mathrm{cas}}^{(1\,\mathrm{sec})}\approx 0.190$ bits/update and
$\bar q_{\mathrm{cas}}^{(1\,\mathrm{sec})}\approx 5.45\times 10^{-22}$
J/update.

\subsection{Macro-energy sanity check: not a surrogate}

Even at $\nu=10\,\mathrm{Hz}$ the 71-sector network dissipates only
$W_{\mathrm{BEA}}\approx 3.87\times 10^{-19}\,\mathrm{W}$; reaching
$1\,\mathrm{W}$ would require $\nu_{\mathrm{cas}}\approx 2.58\times
10^{19}\,\mathrm{Hz}$, and matching the $2.85\times 10^{12}\,\mathrm{W}$
world-electricity benchmark would require $\nu_{\mathrm{cas}}\approx
7.36\times 10^{31}\,\mathrm{Hz}$. Hence the Landauer-normalized BEA quantity
is an \emph{informational lower bound} for the abstract 71-sector relay
update, not a surrogate for gross macro energy consumption; no dollar/bit or
dollar/joule conversion is supplied at the present source-stack level.

\subsection{Rank-4 closure status}

\begin{proposition}[Conditional rank-$4$ closure]
The present archive achieves rank-$3$ closure for the selector triple
$(\Lambda_G,\Gamma_G,T_{\mathrm{info}})$. Adding one independent clock
observable---either the primitive stage clock $\nu_0$ or the completed
cascade rate $\nu_{\mathrm{cas}}=\nu_0/H$---immediately closes the system at
rank $4$. The manuscript therefore claims rank-$3$ thermodynamic closure and
conditional rank-$4$ closure once a single clock observable is supplied.
\end{proposition}

The most defensible throughput identification is to treat the unresolved
fourth observable as the primitive stage clock $\nu_0$ and derive the BEA
completed-cascade rate by $\nu_{\mathrm{cas}}=\nu_0/H^{\ast}_{\mathrm{BEA}}$.
Under that choice the present source stack fully determines the
bits-per-update and joules-per-update factors, and all remaining ambiguity
collapses to one measurable clock.

% ====================================================================
% End of V62 additions.
% ====================================================================

% ====================================================================
% V65 additions: Dimensionality conjecture, 4D/5D design, Corollary E.
% ====================================================================

\subsection{Dimensionality Conjecture for Kato Regime Separability}\label{si:dim_conjecture}

\paragraph{Statement (exploratory).}
For $k$ distinct integer-$H$ Kato regimes coexisting in a population of
city-level observables, let $\boldsymbol{\beta}_i=(\beta_{i,1},\ldots,
\beta_{i,d})^{\top}\in\mathbb{R}^d$ denote the theoretical slope vector
of regime $i$ across $d$ scaling indicators, with
$\beta_{i,j}=1\pm\epsilon(H_i)$ and $\epsilon(H)=1/[H\ln(2H+1)]$.
If (a) the slope separation $\Delta_{\beta}=\min_{i\neq j}
\|\boldsymbol{\beta}_i-\boldsymbol{\beta}_j\|>0$, (b) the $d\times k$
matrix $[\boldsymbol{\beta}_1\,\cdots\,\boldsymbol{\beta}_k]$ has full
affine rank, and (c) per-regime residual covariance is non-degenerate,
then $d\geq k$ is sufficient for asymptotic identifiability of regime
membership via integer-$H$-locked Gaussian-mixture clustering, following
the classical identifiability results of
Teicher~\cite{teicher1963identifiability} and
Yakowitz--Spragins~\cite{yakowitz1968identifiability}.

\paragraph{Finite-sample bound.}
Separation in a finite sample of size $n$ holds with probability
$1-o(1)$ once $n\geq N_{0}(\Delta_{\beta},\boldsymbol{\Sigma})=
\mathcal{O}(k^{2}\log k\,/\,\Delta_{\beta}^{2})$. The empirical
2D~$\to$~3D progression observed here (population--patents $k=2$ by BIC
at $n=386$; population--patents--GDP $k=3$ by BIC at $n=378$) is
consistent with this scaling.

\paragraph{Limitations.}
Counter-examples exclude (i) degenerate covariance, (ii) $\sigma=0$
transient branches, (iii) non-Kato outliers, (iv) insufficient sample
size, and (v) non-integer $H$. The conjecture is stated conditionally.

\paragraph{Relation to Bettencourt's single-dimensional analysis.}
The 1D pooled baseline of~\cite{Bettencourt2007} achieves CI consistency
on 17/22 indicators; the 3D cluster-lock procedure in this work raises
this to approximately 22/22 by isolating the dominant integer-$H$
component within each indicator, consistent with the
identifiability-under-full-rank hypothesis above.

\subsection{Proposed 4D/5D Experimental Design}\label{si:dim_4D5D_design}

\paragraph{Motivation.}
The 3D clustering result with $k=3$ suggests at least three integer-$H$
regimes are jointly active in the US-MSA panel. A direct test of the
dimensionality conjecture is to increase the feature dimension to $d=4$
and $d=5$ with curated indicator panels chosen so that the theoretical
$d\times k$ slope matrix attains full affine rank.

\paragraph{Panel construction.}
Proposed $d=4$ panels: $\{$population, patents, GDP, energy$\}$;
$\{$population, patents, GDP, $\mathrm{CO}_2\}$;
$\{$population, patents, GDP, wages$\}$. Proposed $d=5$ panels add
$\mathrm{CO}_2$ or transportation throughput. Each indicator must have
an independently estimable $H$ under the Kato rubric; overlap between
indicators sharing the same $H$ reduces effective rank and should be
avoided.

\paragraph{Falsification protocol.}
Under the conjecture, BIC-preferred $k$ should grow with $d$ up to the
true number of coexisting integer-$H$ regimes and then saturate.
Deviation from saturation---either premature saturation (insufficient
effective rank) or unbounded growth (non-Kato contamination)---
falsifies the conjecture for that panel.

\subsection{Effective rank bound by observable civilization stage}\label{si:dim_corollary_civ}

\paragraph{Corollary E (effective resolvable rank).}
Let $H_{\mathrm{civ}}$ be the integer-$H$ stage characterizing the
civilization from which the panel is drawn, and let
$\{w_{H}\}_{H\in\mathbb{Z}_{\geq 1}}$ denote the population-weighted
regime mixture. For any weight threshold $w_{\mathrm{thr}}>0$ below
which a regime contributes negligibly to the observed signal, the
effective resolvable rank satisfies
\begin{equation}
k_{\mathrm{eff}} \;\approx\;
\bigl|\{H \leq H_{\mathrm{civ}} : w_{H}>w_{\mathrm{thr}}\}\bigr|
\;\leq\; H_{\mathrm{civ}}.
\end{equation}

\paragraph{Practical consequence.}
For a contemporary ($H_{\mathrm{civ}}=3$) panel, regimes with $H=1$
contribute below detection threshold after patent-era mixing; the
empirically resolvable rank collapses to the top two active regimes.
Hence clustering at dimension $d=3$ can recover the dominant $k=2$
structure even when a latent $H=1$ branch exists, and a $d=4$ panel
need not recover $k=4$ unless a currently dormant regime becomes
weight-dominant at the sampled cities.

\paragraph{Caveats.}
The bound is heuristic: $w_{\mathrm{thr}}$ depends on sample size and
noise, and $H_{\mathrm{civ}}$ is itself an empirical assignment. The
corollary is therefore a practical upper bound on clustering rank
rather than a hard identifiability constraint.

% ====================================================================
% End of V65 additions.
% ====================================================================

\end{document}